% Le Panafricanisme et le développement de l'Afrique — Philosophie Tle Tech — Cours signé OnBuch+
% Programme officiel MINESEC. Compiler avec : tectonic N22-panafricanisme-developpement-afrique.tex
\newcommand{\DOCMATIERE}{Philosophie}
\newcommand{\DOCNIVEAU}{Tle Tech}
\newcommand{\DOCMODULE}{Philosophie – classe de Terminale technique}
\newcommand{\DOCLECON}{N22}
\newcommand{\DOCTITRE}{Le Panafricanisme et le développement de l'Afrique}
% OnBuch+ — préambule « cours signés » ENRICHI (compilé avec tectonic / XeLaTeX)
% Système visuel « Pop » d'OnBuch+ : fond crème (jamais blanc), bordures encre
% épaisses, ombres « dures » décalées, titres Archivo Black en capitales.
% Version MATHÉMATIQUES : graphes (pgfplots/tikz), boîtes pédagogiques
% (définition, propriété, méthode, attention, le savais-tu, exercice résolu,
% exercice, TP, bilan), page de garde et signature OnBuch+ ancrée en bas.
%
% Macros attendues dans main.tex AVANT \input{preamble} :
%   \DOCMATIERE  \DOCNIVEAU  \DOCMODULE  \DOCLECON (numéro)  \DOCTITRE
\documentclass[11pt]{article}

\usepackage{fontspec}
\usepackage[french]{babel}
\usepackage[a4paper,margin=2cm,top=3.4cm,bottom=2.8cm,headheight=1.4cm,footskip=1.3cm]{geometry}
\usepackage{amsmath,amssymb,amsfonts}
\usepackage{mathtools} % \xmapsto, \coloneqq... souvent utilisés par les modèles
\usepackage{cancel} % simplifications barrées (bilans de forces)
% \abs{z} (module, valeur absolue) et \norm{u} (norme), utilisés par les modèles
\providecommand{\abs}{}\let\abs\relax\DeclarePairedDelimiter{\abs}{\lvert}{\rvert}
\providecommand{\norm}{}\let\norm\relax\DeclarePairedDelimiter{\norm}{\lVert}{\rVert}
\usepackage[table]{xcolor} % [table] : fournit aussi \rowcolors (alternance de lignes)
\usepackage{pagecolor}
\usepackage{tikz}
\usetikzlibrary{arrows.meta,positioning,calc,shapes.geometric,shapes.misc,shapes.arrows,decorations.pathmorphing,decorations.pathreplacing,decorations.markings,patterns,shadows,mindmap,trees,fit,backgrounds,matrix,intersections,angles,quotes}
\usepackage{pgfplots}
\usepackage[version=4]{mhchem} % équations nucléaires \ce{^{235}_{92}U}
\usepackage[european,siunitx]{circuitikz} % schémas électriques
\pgfplotsset{compat=1.18}
\usepgfplotslibrary{fillbetween,groupplots} % aires entre courbes (calcul intégral)
\usepackage{enumitem}
\usepackage{fancyhdr}
% [most] charge toutes les bibliothèques tcolorbox (listings, minted, raster,
% documentation...) et gonfle inutilement la mémoire TeX sur un document long
% (~300 boîtes) : on ne charge que ce qui est réellement utilisé.
\usepackage[skins,breakable]{tcolorbox}
\usepackage{graphicx}
\usepackage{colortbl}
\usepackage{booktabs}
\usepackage{tabularx}
\usepackage{multirow}
\usepackage{diagbox} % case d'en-tête diagonale des tableaux à double entrée
% Colonnes extensibles centrée / gauche / droite (C, L, R), souvent utilisées par les modèles
\newcolumntype{C}{>{\centering\arraybackslash}X}
\newcolumntype{L}{>{\raggedright\arraybackslash}X}
\newcolumntype{R}{>{\raggedleft\arraybackslash}X}
\usepackage{array}
\usepackage{multicol}
\usepackage{siunitx}
% parse-numbers=false : accepte aussi \SI{2,5 \times 10^{-3}}{...}
\sisetup{locale=FR,output-decimal-marker={,},group-minimum-digits=4,parse-numbers=false}
\DeclareSIUnit\ohm{\ensuremath{\symup{\Omega}}} % U+2126 absent de la police
\DeclareSIUnit\division{div} % réglages d'oscilloscope (ms/div, V/div)
\DeclareSIUnit\tr{tr} % tours (vitesse de rotation, ex. \SI{1500}{\tr\per\minute})
\DeclareSIUnit\molar{M} % concentration molaire (ex. \SI{3}{\molar}, \SI{1}{\micro\molar})
\DeclareSIUnit\an{an} % année (le modèle confond parfois avec \year, un compteur TeX) — ex. \SI{5}{\kilo\meter\per\an}
\DeclareSIUnit\FCFA{FCFA} % franc CFA (ex. \SI{500}{\FCFA\per\kilo\gram})
\DeclareSIUnit\degreeBrix{\textdegree Brix} % teneur en sucre d'un sirop/jus (réfractométrie)
\DeclareSIUnit\month{mois} % le modèle utilise parfois \month, un compteur TeX, comme unité de durée
\DeclareSIUnit\microliter{\micro\liter} % le modèle écrit parfois \microliter en un seul token
\DeclareSIUnit\million{millions} % dénombrement (ex. \SI{15}{\million\per\mL} spermatozoïdes)
\usepackage{qrcode}
% \degree : commande fréquemment utilisée par les modèles pour un angle ou une
% température en dehors de \SI{}{} (non fournie par les paquets chargés).
\providecommand{\degree}{^{\circ}}
% \qed (fin de démonstration), utilisé par les modèles sans amsthm
\providecommand{\qed}{\hfill$\square$}
% \barre : double barre des valeurs interdites dans un tableau de variations (tkz-tab)
\providecommand{\barre}{\ensuremath{\|}}
% environnement proof (amsthm non chargé) utilisé par les modèles
\ifdefined\proof\else\newenvironment{proof}[1][Démonstration]{\par\noindent\textit{#1.}\ }{\qed\par}\fi
\usepackage{lastpage}
\usepackage{hyperref}
\hypersetup{hidelinks}

% ============================================================
% Palette « Pop » OnBuch+ — mêmes hex que la classe P (plus_ui.dart)
% ============================================================
\definecolor{poporange}{HTML}{F59321}
\definecolor{popdark}{HTML}{E07A0C}
\definecolor{popcream}{HTML}{FFF6E8}
\definecolor{popink}{HTML}{1C1714}
\definecolor{poppurple}{HTML}{7A5AE0}
\definecolor{popgreen}{HTML}{2E9E6A}
\definecolor{poppink}{HTML}{E0457B}
\definecolor{popblue}{HTML}{2F7DE1}
\definecolor{popgold}{HTML}{C98A12}
\definecolor{popmuted}{HTML}{8A7F76}
\definecolor{popline}{HTML}{EADFD2}
\definecolor{popgray}{HTML}{8A7F76}
\colorlet{popgrey}{popgray}
% Alias tolérants (le modèle nomme parfois les couleurs différemment)
\colorlet{popwhite}{white}
\colorlet{popblack}{popink}
\colorlet{popred}{poppink}
\colorlet{popyellow}{popgold}
% Teintes claires (fonds de tableaux/encadrés) — le modèle nomme parfois
% les couleurs avec un suffixe L (ex. popblueL) sans les avoir définies.
\colorlet{poporangeL}{poporange!20}
\colorlet{popdarkL}{popdark!20}
\colorlet{poppurpleL}{poppurple!20}
\colorlet{popgreenL}{popgreen!20}
\colorlet{poppinkL}{poppink!20}
\colorlet{popblueL}{popblue!20}
\colorlet{popgoldL}{popgold!20}
\colorlet{popredL}{popred!20}
\colorlet{popgrayL}{popgray!20}
% Teintes claires pour fonds de tableaux / zones
\colorlet{popblueL}{popblue!10!white}
\colorlet{poporangeL}{poporange!14!white}
\colorlet{popgreenL}{popgreen!12!white}
\colorlet{poppurpleL}{poppurple!12!white}
\colorlet{poppinkL}{poppink!12!white}
\colorlet{popgoldL}{popgold!14!white}

\pagecolor{popcream}

% ============================================================
% Typographie
% ============================================================
\defaultfontfeatures{Ligatures=TeX}
\setmainfont{PlusJakartaSans-Medium.ttf}[Path=../../../fonts/,BoldFont=PlusJakartaSans-Bold.ttf]
% Maths en sans-serif (Fira Math) avec chiffres et lettres droites de Plus
% Jakarta, pour que les formules se fondent dans le texte.
\usepackage[math-style=ISO,bold-style=ISO]{unicode-math}
\setmathfont{FiraMath-Regular.otf}[Path=../../../fonts/]
\setmathfont{PlusJakartaSans-Medium.ttf}[Path=../../../fonts/,range={up/num,up/{Latin,latin},\mathup}]
\usepackage{icomma} % virgule décimale sans espace en mode math
% Caractères mathématiques Unicode absents des polices de texte (Plus Jakarta,
% Archivo Black) : rendus par leur équivalent mathématique, en texte comme en
% formule — sinon ils s'affichent en carré vide (ex. « Arithmétique dans ℤ »).
\usepackage{newunicodechar}
\newunicodechar{ℝ}{\ensuremath{\mathbb{R}}}
\newunicodechar{ℤ}{\ensuremath{\mathbb{Z}}}
\newunicodechar{ℂ}{\ensuremath{\mathbb{C}}}
\newunicodechar{ℕ}{\ensuremath{\mathbb{N}}}
\newunicodechar{ℚ}{\ensuremath{\mathbb{Q}}}
\newunicodechar{𝔻}{\ensuremath{\mathbb{D}}}
\newunicodechar{α}{\ensuremath{\alpha}}
\newunicodechar{β}{\ensuremath{\beta}}
\newunicodechar{γ}{\ensuremath{\gamma}}
\newunicodechar{δ}{\ensuremath{\delta}}
\newunicodechar{ε}{\ensuremath{\varepsilon}}
\newunicodechar{θ}{\ensuremath{\theta}}
\newunicodechar{λ}{\ensuremath{\lambda}}
\newunicodechar{μ}{\ensuremath{\mu}}
\newunicodechar{σ}{\ensuremath{\sigma}}
\newunicodechar{τ}{\ensuremath{\tau}}
\newunicodechar{φ}{\ensuremath{\varphi}}
\newunicodechar{ψ}{\ensuremath{\psi}}
\newunicodechar{ω}{\ensuremath{\omega}}
\newunicodechar{Σ}{\ensuremath{\Sigma}}
% majuscules produites par \MakeUppercase dans les titres (α-aminés -> Α)
\newunicodechar{Α}{\ensuremath{\alpha}}
\newunicodechar{Β}{\ensuremath{\beta}}
\newunicodechar{Γ}{\ensuremath{\Gamma}}
\newunicodechar{Δ}{\ensuremath{\Delta}}
\newunicodechar{Ε}{\ensuremath{\varepsilon}}
\newunicodechar{Θ}{\ensuremath{\Theta}}
\newunicodechar{Λ}{\ensuremath{\Lambda}}
\newunicodechar{Μ}{\ensuremath{\mu}}
\newunicodechar{Τ}{\ensuremath{\tau}}
\newunicodechar{Φ}{\ensuremath{\Phi}}
\newunicodechar{Ψ}{\ensuremath{\Psi}}
\newunicodechar{Ω}{\ensuremath{\Omega}}
\newunicodechar{↦}{\ensuremath{\mapsto}}
\newunicodechar{∈}{\ensuremath{\in}}
\newunicodechar{∉}{\ensuremath{\notin}}
\newunicodechar{∧}{\ensuremath{\wedge}}
\newunicodechar{⇔}{\ensuremath{\Leftrightarrow}}
\newunicodechar{⟺}{\ensuremath{\Longleftrightarrow}}
\newunicodechar{⇒}{\ensuremath{\Rightarrow}}
\newunicodechar{⟹}{\ensuremath{\Longrightarrow}}
\newunicodechar{∘}{\ensuremath{\circ}}
\newunicodechar{∪}{\ensuremath{\cup}}
\newunicodechar{∩}{\ensuremath{\cap}}
\newunicodechar{≡}{\ensuremath{\equiv}}
\newunicodechar{⊂}{\ensuremath{\subset}}
\newunicodechar{⊥}{\ensuremath{\perp}}
\newunicodechar{∃}{\ensuremath{\exists}}
\newunicodechar{∀}{\ensuremath{\forall}}
\newunicodechar{∛}{\ensuremath{\sqrt[3]{\,}}}
\newunicodechar{∜}{\ensuremath{\sqrt[4]{\,}}}
\newunicodechar{⃗}{\ensuremath{{}^{\to}}}
\newunicodechar{ⁿ}{\ifmmode^{n}\else\textsuperscript{\ensuremath{n}}\fi}
\newunicodechar{ˣ}{\ifmmode^{x}\else\textsuperscript{\ensuremath{x}}\fi}
\newunicodechar{ᵇ}{\ifmmode^{b}\else\textsuperscript{\ensuremath{b}}\fi}
\newunicodechar{ʳ}{\ifmmode^{r}\else\textsuperscript{\ensuremath{r}}\fi}
\newunicodechar{ᵘ}{\ifmmode^{u}\else\textsuperscript{\ensuremath{u}}\fi}
\newunicodechar{ʸ}{\ifmmode^{y}\else\textsuperscript{\ensuremath{y}}\fi}
\newunicodechar{ᵃ}{\ifmmode^{a}\else\textsuperscript{\ensuremath{a}}\fi}
\newunicodechar{ᵉ}{\ifmmode^{e}\else\textsuperscript{\ensuremath{e}}\fi}
\newunicodechar{ᵏ}{\ifmmode^{k}\else\textsuperscript{\ensuremath{k}}\fi}
\newunicodechar{ᵖ}{\ifmmode^{p}\else\textsuperscript{\ensuremath{p}}\fi}
\newunicodechar{ᴺ}{\ifmmode^{N}\else\textsuperscript{\ensuremath{N}}\fi}
\newunicodechar{ᶜ}{\ifmmode^{c}\else\textsuperscript{\ensuremath{c}}\fi}
\newunicodechar{ʰ}{\ifmmode^{h}\else\textsuperscript{\ensuremath{h}}\fi}
\newunicodechar{ᵠ}{\ifmmode^{q}\else\textsuperscript{\ensuremath{q}}\fi}
\newunicodechar{ₙ}{\ifmmode_{n}\else\textsubscript{\ensuremath{n}}\fi}
\newunicodechar{ₐ}{\ifmmode_{a}\else\textsubscript{\ensuremath{a}}\fi}
\newunicodechar{ᵢ}{\ifmmode_{i}\else\textsubscript{\ensuremath{i}}\fi}
\newunicodechar{ᵣ}{\ifmmode_{r}\else\textsubscript{\ensuremath{r}}\fi}
\newunicodechar{ₖ}{\ifmmode_{k}\else\textsubscript{\ensuremath{k}}\fi}
\newunicodechar{ᵦ}{\ifmmode_{\beta}\else\textsubscript{\ensuremath{\beta}}\fi}
\newunicodechar{ₓ}{\ifmmode_{x}\else\textsubscript{\ensuremath{x}}\fi}
\newfontfamily\popdisplay{ArchivoBlack-Regular.ttf}[Path=../../../fonts/]
\newfontfamily\poplogo{SpaceGrotesk-Bold.ttf}[Path=../../../fonts/]
\newfontfamily\popsemi{PlusJakartaSans-SemiBold.ttf}[Path=../../../fonts/]
\newfontfamily\popxbold{PlusJakartaSans-ExtraBold.ttf}[Path=../../../fonts/]
\newfontfamily\popxboldls{PlusJakartaSans-ExtraBold.ttf}[Path=../../../fonts/,LetterSpace=3.0]

\setlist[enumerate]{leftmargin=1.7em,itemsep=5pt,topsep=4pt}
\setlist[itemize]{leftmargin=1.7em,itemsep=4pt,topsep=4pt,label={\color{poporange}\textbullet}}
\setlength{\parindent}{0pt}
\setlength{\parskip}{5pt}
\renewcommand{\arraystretch}{1.25}

% pgfplots : style Pop par défaut
\pgfplotsset{
  every axis/.append style={axis line style={popink,line width=1pt},tick style={popink},
    grid style={popline,line width=0.5pt},label style={font=\small},tick label style={font=\footnotesize}},
  popcurve/.style={poporange,line width=1.8pt,no markers,smooth},
}
% Même style disponible dans un \draw TikZ simple (hors pgfplots)
\tikzset{popcurve/.style={poporange,line width=1.8pt}}
\tikzset{popgrid/.style={popline,very thin}}
\colorlet{popcurve}{poporange} % le nom du style est parfois utilisé comme couleur

% ============================================================
% Pill / squiggle
% ============================================================
\newcommand{\poppill}[3]{%
  \tikz[baseline=-0.55ex]\node[draw=popink,line width=1.2pt,rounded corners=2.6pt,%
    fill=#2,inner xsep=6.5pt,inner ysep=3pt]{%
    \popxboldls\fontsize{7}{7}\selectfont\color{#3}\MakeUppercase{#1}};%
}
\newcommand{\popsquiggle}[1]{%
  \tikz[baseline=-0.4ex]{
    \draw[#1,line width=2.6pt,line cap=round]
      plot [smooth] coordinates {(0,0.09) (0.34,0.01) (0.68,0.21) (1.02,0.01) (1.36,0.10)};
  }%
}
% Mot-clé surligné (marqueur orange sous le texte)
\newcommand{\cle}[1]{{\popxbold\color{popdark}#1}}

% ============================================================
% En-tête / pied
% ============================================================
\pagestyle{fancy}
\fancyhf{}
\renewcommand{\headrulewidth}{0pt}
\renewcommand{\footrulewidth}{0pt}
\lhead{\raisebox{-0.15ex}{\poplogo\fontsize{13}{13}\selectfont\color{popink}OnBuch\color{poporange}+}}
\rhead{\poppill{\DOCMATIERE\ \textbullet\ \DOCNIVEAU\ \textbullet\ Leçon \DOCLECON}{white}{popink}}
\lfoot{\popxbold\fontsize{7.6}{7.6}\selectfont\color{popmuted}\MakeUppercase{Cours signé OnBuch+}\ \textbullet\ {\popsemi\DOCTITRE}}
\rfoot{\tikz[baseline=-0.6ex]\node[rounded corners=8.5pt,fill=popink,minimum height=17pt,minimum width=17pt,inner xsep=5pt,inner sep=0]{\popxbold\fontsize{8}{8}\selectfont\color{popcream}\thepage\,/\,\pageref*{LastPage}};}

% ============================================================
% Titres
% ============================================================
\newcommand{\chaptitre}[2]{%
  \begin{tcolorbox}[enhanced,colback=poporange,colframe=popink,boxrule=2.6pt,arc=11pt,
      drop shadow=popink,left=15pt,right=15pt,top=12pt,bottom=11pt,before skip=3pt,after skip=18pt]
    \poppill{#2}{popink}{popcream}\par\vspace{9pt}
    {\raggedright\hyphenpenalty=10000\exhyphenpenalty=10000\popdisplay\fontsize{22}{24}\selectfont\color{popcream}\MakeUppercase{#1}\par}
  \end{tcolorbox}
}
% Section numérotée : pastille orange avec le numéro + titre Archivo
\newcounter{popsec}
\newcommand{\coursec}[1]{%
  \stepcounter{popsec}\setcounter{popsub}{0}%
  \par\vspace{16pt}\noindent
  \begin{minipage}{\linewidth}
  \tikz[baseline=-0.7ex]\node[draw=popink,line width=1.6pt,fill=poporange,rounded corners=4pt,minimum size=22pt,inner sep=0]{\popdisplay\fontsize{12}{12}\selectfont\color{popcream}\thepopsec};%
  \hspace{8pt}{\popdisplay\fontsize{15}{17}\selectfont\color{popink}\MakeUppercase{#1}}\par
  \vspace{4pt}\noindent{\color{poporange}\rule{36pt}{3.6pt}}
  \end{minipage}\par\nopagebreak\vspace{8pt}
}
\newcounter{popsub}
\newcommand{\courssub}[1]{%
  \stepcounter{popsub}%
  \par\vspace{9pt}\noindent{\popxbold\fontsize{11.5}{13}\selectfont\color{popink}{\color{poporange}\thepopsec.\thepopsub}\hspace{6pt}#1}\par\nopagebreak\vspace{4pt}
}

% ============================================================
% Boîtes pédagogiques — même recette : fond blanc, bordure épaisse colorée,
% ombre dure encre, badge plein. Titre optionnel [#1].
% ============================================================
\tcbset{popbase/.style={breakable,enhanced,colback=white,boxrule=2.3pt,arc=9pt,
  shadow={3pt}{-3pt}{0pt}{popink},left=14pt,right=14pt,top=11pt,bottom=11pt,before skip=11pt,after skip=15pt,
  fontupper=\popsemi\fontsize{10.6}{14.5}\selectfont\color{popink},
  fontlower=\popsemi\fontsize{10.6}{14.5}\selectfont\color{popink}}}
% Coche dessinée (le glyphe ✓ n'existe pas dans les polices du document)
\newcommand{\popcheck}[1]{\tikz[baseline=-0.5ex]\draw[#1,line width=1.6pt,line cap=round,line join=round](-0.13,0.0)--(-0.04,-0.09)--(0.13,0.1);}
\providecommand{\coche}{\popcheck{popgreen}}
\providecommand{\resultat}[1]{\par\smallskip\noindent\fcolorbox{popgreen}{popgreenL}{\parbox{\dimexpr\linewidth-2\fboxsep-2\fboxrule\relax}{\textbf{Résultat.} #1}}\par\smallskip}
\newcommand{\popboxhead}[3]{\poppill{#1}{#2}{white}\ifx\relax#3\relax\else\hspace{6pt}{\popxbold\fontsize{10.6}{12}\selectfont\color{popink}#3}\fi\par\vspace{7pt}}

\newtcolorbox{aretenir}[1][]{popbase,colframe=popgreen,before upper={\popboxhead{À retenir}{popgreen}{#1}}}
\newtcolorbox{exemplebox}[1][]{popbase,colframe=poporange,before upper={\popboxhead{Exemple}{poporange}{#1}}}
\newtcolorbox{definition}[1][]{popbase,colframe=popblue,before upper={\popboxhead{Définition}{popblue}{#1}}}
\newtcolorbox{propriete}[1][]{popbase,colframe=poppurple,before upper={\popboxhead{Propriété}{poppurple}{#1}}}
\newtcolorbox{methode}[1][]{popbase,colframe=popgold,before upper={\popboxhead{Méthode}{popgold}{#1}}}
\newtcolorbox{attention}[1][]{popbase,colframe=poppink,before upper={\popboxhead{Attention}{poppink}{#1}}}
\newtcolorbox{experience}[1][]{popbase,colframe=popblue,colback=popblueL,before upper={\popboxhead{Expérience}{popblue}{#1}}}
% « Le savais-tu ? » : carte violette pleine (contexte, culture, Cameroun)
\newtcolorbox{savaistu}[1][]{popbase,colback=poppurple,colframe=popink,
  fontupper=\popsemi\fontsize{10.4}{14.2}\selectfont\color{white},
  before upper={\poppill{Le savais-tu\,?}{white}{poppurple}\ifx\relax#1\relax\else\hspace{6pt}{\popxbold\color{white}#1}\fi\par\vspace{7pt}}}
% Exercice résolu : énoncé en haut, \tcblower puis la solution sur fond crème
\newcounter{popexo}\newcounter{popexr}
\newtcolorbox{exoresolu}[1][]{popbase,colframe=poporange,colbacklower=popcream,
  segmentation style={popink,line width=1.2pt,dashed},
  before upper={\stepcounter{popexr}\popboxhead{Exercice résolu \thepopexr}{poporange}{#1}},
  before lower={\poppill{Solution}{popink}{popcream}\par\vspace{6pt}}}
% Exercice (non résolu en place) — la correction est dans la partie Corrigés
\newtcolorbox{exercice}[1][]{popbase,colframe=popink,
  before upper={\stepcounter{popexo}\popboxhead{Exercice \thepopexo}{popink}{#1}}}
% Corrigé
\newtcolorbox{corrige}[1][]{popbase,colframe=popgreen,colback=popgreenL,
  before upper={\popboxhead{Corrigé}{popgreen}{#1}}}
% Objectifs / prérequis / bilan
\newtcolorbox{objectifs}{popbase,colframe=popink,colback=poporangeL,
  before upper={\popboxhead{Objectifs de la leçon}{popink}{}}}
\newtcolorbox{prerequis}{popbase,colframe=popink,colback=popblueL,
  before upper={\popboxhead{Prérequis}{popblue}{}}}
\newtcolorbox{bilan}{popbase,colframe=popink,colback=white,boxrule=2.8pt,
  before upper={\popboxhead{Fiche bilan}{poporange}{L'essentiel à connaître pour l'examen}}}
% Figure encadrée avec légende
\newtcolorbox{popfigure}[1][]{enhanced,breakable=false,colback=white,colframe=popline,boxrule=1.4pt,arc=8pt,
  left=10pt,right=10pt,top=10pt,bottom=8pt,before skip=10pt,after skip=12pt,halign=center,
  lower separated=false,halign lower=center,
  fontlower=\popsemi\fontsize{9}{11.5}\selectfont\color{popmuted},
  #1}
\newcommand{\legende}[1]{\par\vspace{6pt}{\popsemi\fontsize{9}{11.5}\selectfont\color{popmuted}#1}}

% ============================================================
% Signature OnBuch+
% ============================================================
\newcommand{\poplogoinline}[1]{{\poplogo\fontsize{#1}{#1}\selectfont\color{popink}OnBuch\color{poporange}+}}
\newcommand{\signatureonbuch}{%
  \begin{tcolorbox}[enhanced,colback=popink,colframe=popink,boxrule=2.2pt,arc=10pt,
      drop shadow=poporange,left=14pt,right=14pt,top=10pt,bottom=10pt,before skip=0pt,after skip=0pt]
    \tikz[baseline=-0.7ex]\node[circle,fill=popgreen,draw=popcream,line width=1.3pt,minimum size=22pt,inner sep=0]{\popcheck{white}};%
    \hspace{9pt}\begin{minipage}[c]{0.62\linewidth}
      {\popxbold\fontsize{12}{13}\selectfont\color{popcream}Signé\ \textbullet\ {\poplogo OnBuch\color{poporange}+}}\par\vspace{2pt}
      {\raggedright\popsemi\fontsize{8}{10.5}\selectfont\color{popline}\MakeUppercase{Équipe pédagogique OnBuch+}\\\MakeUppercase{Conforme au programme officiel MINESEC}\par}
    \end{minipage}\hfill
    {\poplogo\fontsize{22}{22}\selectfont\color{popcream}OnBuch\color{poporange}+}
  \end{tcolorbox}%
}

% ============================================================
% Page de garde
% #1 titre · #2 matière · #3 niveau · #4 module · #5 n° leçon · #6 durée ·
% #7 description / objectifs (texte ou itemize)
% ============================================================
\newcommand{\pagegarde}[7]{%
  \thispagestyle{empty}
  \newgeometry{a4paper,margin=1.7cm,top=1.6cm,bottom=1.4cm}%
  \begin{tikzpicture}[remember picture,overlay]
    \fill[poporange!18] ([xshift=-3.2cm,yshift=-3.6cm]current page.north east) circle (4.6cm);
    \fill[poppurple!14] ([xshift=2.4cm,yshift=7.6cm]current page.south west) circle (3.8cm);
  \end{tikzpicture}%
  {\poplogo\fontsize{30}{30}\selectfont\color{popink}OnBuch\color{poporange}+}\hfill
  \raisebox{0.6ex}{\poppill{Cours signé \textbullet\ Premium}{popink}{popcream}}\par
  \vspace{1pt}\popsquiggle{poppurple}\par
  \vspace{8pt}
  \begin{tcolorbox}[enhanced,colback=poporange,colframe=popink,boxrule=2.8pt,arc=13pt,
      drop shadow=popink,left=18pt,right=18pt,top=12pt,bottom=13pt,after skip=10pt]
    \poppill{#2 \textbullet\ #3}{popink}{popcream}\hspace{5pt}\poppill{#4}{white}{popink}\par\vspace{8pt}
    {\popdisplay\fontsize{14}{14}\selectfont\color{popink}LEÇON #5}\par\vspace{4pt}
    {\raggedright\hyphenpenalty=10000\exhyphenpenalty=10000\popdisplay\fontsize{25}{27}\selectfont\color{popcream}\MakeUppercase{#1}\par}
    \vspace{8pt}
    {\popxbold\fontsize{9.5}{9.5}\selectfont\color{popink}\MakeUppercase{Durée conseillée : #6}}
  \end{tcolorbox}
  \begin{tcolorbox}[enhanced,colback=white,colframe=popink,boxrule=2.2pt,arc=10pt,
      drop shadow=popink,left=16pt,right=16pt,top=10pt,bottom=10pt,after skip=8pt]
    \poppill{Dans cette leçon}{poppurple}{white}\par\vspace{5pt}
    \popsemi\fontsize{9.3}{12.6}\selectfont\color{popink}#7
  \end{tcolorbox}
  \noindent\begin{minipage}[t]{0.487\linewidth}
    \begin{tcolorbox}[enhanced,colback=popgreen,colframe=popink,boxrule=2.2pt,arc=10pt,
        drop shadow=popink,left=11pt,right=11pt,top=10pt,bottom=10pt,height=3.1cm,valign=center]
      \tcbox[colback=white,colframe=popink,boxrule=1.3pt,arc=5pt,boxsep=0pt,nobeforeafter,
        left=3pt,right=3pt,top=3pt,bottom=3pt,valign=center]{\qrcode[height=1.9cm]{https://play.google.com/store/apps/details?id=com.luvvix.onbuch&hl=fr}}%
      \hspace{8pt}\begin{minipage}[c]{0.5\linewidth}
        {\popdisplay\fontsize{11}{12.5}\selectfont\color{white}\MakeUppercase{Télécharge OnBuch}}\par\vspace{4pt}
        {\raggedright\popsemi\fontsize{8.4}{11}\selectfont\color{white}Gratuit \textbullet\ Google Play\\Tuteur IA, annales, concours, résultats\par}
      \end{minipage}
    \end{tcolorbox}
  \end{minipage}\hfill
  \begin{minipage}[t]{0.487\linewidth}
    \begin{tcolorbox}[enhanced,colback=poppurple,colframe=popink,boxrule=2.2pt,arc=10pt,
        drop shadow=popink,left=11pt,right=11pt,top=10pt,bottom=10pt,height=3.1cm,valign=center]
      \tikz[baseline=-0.7ex]\node[circle,fill=white,draw=popink,line width=1.3pt,minimum size=1.6cm,inner sep=0]{\popdisplay\fontsize{24}{24}\selectfont\color{poppurple}+};%
      \hspace{8pt}\begin{minipage}[c]{0.6\linewidth}
        {\popdisplay\fontsize{11}{12.5}\selectfont\color{white}\MakeUppercase{Passe à OnBuch+}}\par\vspace{4pt}
        {\raggedright\popsemi\fontsize{8.4}{11}\selectfont\color{white}Tous les cours signés, Tuteur IA illimité, fiches PDF — onglet OnBuch+ de l'app\par}
      \end{minipage}
    \end{tcolorbox}
  \end{minipage}
  \vspace{8pt}
  \popcontenu
  \vfill
  \signatureonbuch
  \restoregeometry
  \newpage
}

% Rangée « Ce cours contient » (page de garde)
\newcommand{\poptile}[3]{% #1 couleur #2 gros texte #3 légende
  \begin{tcolorbox}[enhanced,colback=#1,colframe=popink,boxrule=2pt,arc=9pt,shadow={2.5pt}{-2.5pt}{0pt}{popink},
      left=6pt,right=6pt,top=9pt,bottom=9pt,halign=center,valign=center,height=1.75cm,width=\linewidth,nobeforeafter]
    {\popdisplay\fontsize{12.5}{13}\selectfont\color{white}\MakeUppercase{#2}}\par\vspace{3pt}
    {\centering\popsemi\fontsize{7.8}{9.5}\selectfont\color{white}#3\par}
  \end{tcolorbox}}
\newcommand{\popcontenu}{%
  \par\noindent{\popxboldls\fontsize{7.5}{7.5}\selectfont\color{popmuted}CE COURS SIGNÉ CONTIENT}\par\vspace{7pt}
  \noindent\begin{minipage}[t]{0.235\linewidth}\poptile{popblue}{Cours}{complet, illustré, pas à pas}\end{minipage}\hfill
  \begin{minipage}[t]{0.235\linewidth}\poptile{poporange}{Méthodes}{et exercices résolus}\end{minipage}\hfill
  \begin{minipage}[t]{0.235\linewidth}\poptile{poppink}{Exercices}{type examen + corrigés}\end{minipage}\hfill
  \begin{minipage}[t]{0.235\linewidth}\poptile{popgreen}{Fiche bilan}{l'essentiel à retenir}\end{minipage}\par}

% Sommaire
\newtcolorbox{sommaire}{enhanced,colback=white,colframe=popink,boxrule=2.2pt,arc=10pt,
  drop shadow=popink,left=18pt,right=18pt,top=13pt,bottom=13pt,before skip=6pt,after skip=20pt,
  before upper={\poppill{Sommaire}{poppurple}{white}\par\vspace{9pt}\popsemi\fontsize{10.6}{14.5}\selectfont\color{popink}}}

% Signature de fin de cours — poussée tout en bas de la dernière page
\newcommand{\findecours}{%
  \par\vfill\nopagebreak
  \begin{center}\popsquiggle{poporange}\end{center}\vspace{4pt}
  \signatureonbuch
}
\AtBeginDocument{\def\llbracket{\mathopen{[\mkern-3.5mu[}}\def\rrbracket{\mathclose{]\mkern-3.5mu]}}} % intervalles d'entiers (glyphe ⟦ absent de la police)
% Glyphes absents de Fira Math : redéfinis à partir de glyphes disponibles
\AtBeginDocument{%
  \def\setminus{\mathbin{\backslash}}\let\smallsetminus\setminus
  \def\vdots{\mathord{\vbox{\baselineskip3.2pt\lineskiplimit0pt\kern4pt\hbox{.}\hbox{.}\hbox{.}}}}%
  \def\ddots{\mathinner{\mkern1mu\raise6.5pt\hbox{.}\mkern2mu\raise3.6pt\hbox{.}\mkern2mu\raise0.7pt\hbox{.}\mkern1mu}}%
  \def\triangle{\mathord{\scalebox{0.9}{$\symup{\Delta}$}}}%
  \def\nabla{\mathord{\raisebox{\depth}{\scalebox{1}[-1]{$\symup{\Delta}$}}}}%
  \def\checkmark{\popcheck{popdark}}%
  \def\star{\mathbin{\ast}}\def\bigstar{\mathord{\ast}}%
  \def\diamond{\mathbin{\scalebox{0.6}{$\Diamond$}}}%
  \def\bigcup{\mathop{\mathchoice{\raisebox{-0.3em}{\scalebox{1.7}{$\cup$}}}{\scalebox{1.25}{$\cup$}}{\cup}{\cup}}\displaylimits}%
  \def\bigcap{\mathop{\mathchoice{\raisebox{-0.3em}{\scalebox{1.7}{$\cap$}}}{\scalebox{1.25}{$\cap$}}{\cap}{\cap}}\displaylimits}%
  \def\lesseqqgtr{\mathrel{\lessgtr}}\def\bigoplus{\mathop{\oplus}\displaylimits}%
}
\providecommand{\ssi}{si et seulement si} % abréviation fréquente
\AtBeginDocument{\providecommand{\wideparen}{\overparen}} % arc de cercle

\begin{document}

\pagegarde{Le Panafricanisme et le développement de l'Afrique}{Philosophie}{Tle Tech}{Philosophie – classe de Terminale technique}{N22}{2 séances (fiche de progression harmonisée)}{Cette leçon interroge les notions de développement et de sous-développement appliquées à l'Afrique, puis présente le panafricanisme comme réponse philosophique et politique au retard africain. Elle en retrace l'historique, en dégage les objectifs et en analyse lucidement les difficultés, à partir de situations concrètes camerounaises et africaines.
\vspace{4pt}
\begin{multicols}{2}\small
\begin{itemize}
\item À la fin de cette leçon, je sais définir les concepts de développement et de sous-développement
\item À la fin de cette leçon, je sais distinguer croissance économique et développement humain à partir d'indicateurs (PIB, IDH)
\item À la fin de cette leçon, je sais identifier les manifestations du sous-développement dans la vie quotidienne camerounaise
\item À la fin de cette leçon, je sais retracer les grandes étapes historiques du panafricanisme
\item À la fin de cette leçon, je sais citer les principales figures du panafricanisme et leurs idées
\item À la fin de cette leçon, je sais dégager les objectifs politiques, économiques et culturels du panafricanisme
\end{itemize}\end{multicols}}

\begin{sommaire}
\begin{multicols}{2}
\begin{enumerate}
\item Le développement : définition et indicateurs
\item Le sous-développement de l'Afrique : manifestations et causes
\item Historique du panafricanisme
\item Objectifs du panafricanisme comme philosophie du développement
\item Difficultés et limites du panafricanisme
\item Synthèse : le panafricanisme, utopie ou voie d'avenir ?
\item Activité d'intégration
\item Méthodes et exercices résolus
\item Exercices
\item Corrigés des exercices
\item Fiche bilan
\end{enumerate}
\end{multicols}
\end{sommaire}

\begin{objectifs}
\begin{itemize}
\item À la fin de cette leçon, je sais définir les concepts de développement et de sous-développement
\item À la fin de cette leçon, je sais distinguer croissance économique et développement humain à partir d'indicateurs (PIB, IDH)
\item À la fin de cette leçon, je sais identifier les manifestations du sous-développement dans la vie quotidienne camerounaise
\item À la fin de cette leçon, je sais retracer les grandes étapes historiques du panafricanisme
\item À la fin de cette leçon, je sais citer les principales figures du panafricanisme et leurs idées
\item À la fin de cette leçon, je sais dégager les objectifs politiques, économiques et culturels du panafricanisme
\item À la fin de cette leçon, je sais analyser les difficultés qui freinent l'unité africaine
\item À la fin de cette leçon, je sais appliquer la notion d'unité à des situations concrètes de la vie courante (CEMAC, Union africaine, coopératives locales)
\item À la fin de cette leçon, je sais rédiger et justifier une conclusion argumentée sur l'avenir du panafricanisme
\end{itemize}
\end{objectifs}

\begin{prerequis}
\begin{itemize}
\item Notion de colonisation et de décolonisation de l'Afrique (histoire de 1ère)
\item Notions de liberté, d'État et de justice (leçons de philosophie de Terminale)
\item Notion de culture et d'identité (leçon sur la culture)
\item Repères géographiques de base sur l'Afrique et le Cameroun
\end{itemize}
\end{prerequis}

\newpage

\coursec{Le développement : définition et indicateurs}

En 2023, le Cameroun importe encore plus de 80 \% de ses produits manufacturés, alors que le continent africain regorge de matières premières : pétrole, cacao, cobalt, bois. Un élève de Yaoundé se demande : pourquoi l'Afrique, si riche, reste-t-elle présentée comme « sous-développée » ? Et si la solution passait par l'unité des Africains, comme le rêvaient Nkrumah, Senghor ou Nyerere ? Le panafricanisme est-il une utopie dépassée ou une philosophie du développement encore d'actualité ?

Avant de répondre à ces grandes questions, il faut d'abord s'entendre sur les mots. Que signifie exactement « développement » ? Comment le mesure-t-on ? Un pays qui vend beaucoup de pétrole est-il pour autant développé ? Telle est la problématique de cette première section : \emph{qu'est-ce que le développement et comment peut-on évaluer le niveau de développement d'un pays comme le Cameroun ?}

\courssub{Définir le développement}

\textbf{L'intuition d'abord.} Imagine deux villages de la région de l'Est. Dans le premier, une compagnie étrangère exploite le bois : les camions passent, l'argent circule entre quelques mains, mais l'école reste sans bancs, le centre de santé sans médicaments, et les jeunes partent à Yaoundé. Dans le second, une coopérative locale transforme le cacao sur place : les revenus financent un collège, un forage d'eau potable et une maternité ; les jeunes restent et créent des ateliers. Lequel des deux villages est « développé » ? Intuitivement, tu réponds : le second. Pourtant, le premier rapporte peut-être plus d'argent au niveau national. Cette intuition contient déjà toute la leçon : le développement n'est pas seulement une question d'argent, mais une transformation globale de la vie d'une population.

\begin{definition}[Le développement]
Le \cle{développement} est l'ensemble des transformations quantitatives et qualitatives — économiques, sociales, culturelles et politiques — qui améliorent durablement les conditions de vie matérielles et spirituelles d'une société. C'est un \emph{processus global} : il touche à la fois la production des richesses, leur répartition, la santé, l'éducation, les libertés et la culture d'un peuple.
\end{definition}

Analysons les termes de cette définition. « Transformations \textbf{quantitatives} » : il faut produire davantage de biens et de services (routes, usines, récoltes). « Transformations \textbf{qualitatives} » : il faut aussi changer les structures — moderniser l'agriculture, scolariser les filles, rendre la justice accessible, transformer sur place les matières premières au lieu de les exporter brutes. « \textbf{Durablement} » : un développement qui épuise les forêts et les sols, ou qui dépend entièrement de l'aide extérieure, n'en est pas un. Enfin, « conditions de vie matérielles \emph{et} spirituelles » : l'être humain ne vit pas seulement de pain ; il a besoin de dignité, de culture, de participation à la vie de la cité. Comme le soulignait le prix Nobel d'économie Amartya Sen, développer, c'est élargir les \emph{libertés réelles} des personnes : pouvoir se nourrir, se soigner, apprendre, choisir sa vie.

\courssub{Croissance économique et développement : une distinction capitale}

La confusion la plus répandue consiste à identifier le développement à la \cle{croissance économique}. Or ce sont deux notions distinctes.

\begin{definition}[Croissance économique]
La \cle{croissance économique} est l'augmentation quantitative et soutenue de la production de biens et de services d'un pays, mesurée par l'évolution du \cle{Produit Intérieur Brut} (PIB), c'est-à-dire la valeur totale des richesses produites sur le territoire en une année.
\end{definition}

La croissance est donc une affaire de \emph{quantité} : le PIB augmente de 3 \%, de 5 \%... Le développement est une affaire de \emph{qualité} : à qui profite cette richesse ? Sert-elle à construire des hôpitaux, des écoles, des industries ? Transforme-t-elle les mentalités et les institutions ?

\begin{attention}[Erreur fréquente : confondre croissance et développement]
Un pays peut connaître une forte croissance \emph{sans} se développer. Exemple classique : un pays pétrolier voit son PIB bondir grâce à l'exportation du brut, mais cette richesse reste concentrée entre quelques mains, les inégalités explosent, la pauvreté rurale persiste et l'économie reste dépendante d'un seul produit. À l'inverse, la croissance est généralement une \emph{condition} du développement (difficile de financer écoles et hôpitaux sans richesses), mais elle n'en est pas une condition \emph{suffisante}. Dans une dissertation, ne jamais écrire « le pays se développe » quand on veut dire « le PIB augmente ».
\end{attention}

\begin{exemplebox}[Le « paradoxe pétrolier » en Afrique]
Plusieurs pays africains exportateurs de pétrole ont affiché des taux de croissance élevés (parfois supérieurs à 5 \%) tout en conservant une majorité de la population vivant avec moins de deux dollars par jour. Le Cameroun lui-même exporte du pétrole, du cacao et du bois bruts, mais importe l'essentiel de ses produits manufacturés : la richesse extraite du sol ne se transforme pas automatiquement en bien-être pour tous. C'est exactement ce que dénonçait Kwame Nkrumah dans \emph{Le néo-colonialisme, dernier stade de l'impérialisme} (1965) : une indépendance politique sans maîtrise économique reste une indépendance inachevée.
\end{exemplebox}

\begin{popfigure}
\begin{center}
\begin{tikzpicture}
\draw[fill=popgoldL,draw=popgold,thick] (0,0) circle (3.1);
\draw[fill=popblueL,draw=popblue,thick] (0,0) circle (2.1);
\draw[fill=poporangeL,draw=poporange,thick] (0,0) circle (1.1);
\node[font=\small\bfseries, align=center] at (0,0){\shortstack{Croissance\\économique\\(PIB)}};
\node[font=\small\bfseries,text=popdark, align=center] at (0,1.55){\shortstack{Développement humain\\(santé, éducation, revenu)}};
\node[font=\small\bfseries,text=popdark, align=center] at (0,2.62){\shortstack{Développement durable\\(environnement, équité, avenir)}};
\end{tikzpicture}
\end{center}
\legende{Schéma en cercles concentriques : la croissance économique (au centre) est le noyau nécessaire mais insuffisant ; le développement humain l'élargit à la santé, à l'éducation et au revenu ; le développement durable l'inscrit dans le respect de l'environnement et des générations futures.}
\end{popfigure}

\courssub{Les indicateurs du développement}

Comment comparer objectivement le niveau de développement du Cameroun, du Nigeria ou de la France ? Les économistes ont construit des \cle{indicateurs}, c'est-à-dire des instruments de mesure.

\textbf{1. Le PIB par habitant.} On divise le PIB par le nombre d'habitants. Cela donne une idée du revenu moyen. Limite majeure : c'est une \emph{moyenne} qui masque les inégalités. Si dix personnes gagnent 1 million et quatre-vingt-dix gagnent 10 000 FCFA, la moyenne sera trompeuse.

\textbf{2. L'Indice de Développement Humain (IDH).} Créé pour dépasser la seule logique du revenu, l'\cle{IDH} combine trois dimensions : la \emph{santé} (mesurée par l'espérance de vie à la naissance), l'\emph{éducation} (durée moyenne et durée attendue de scolarisation) et le \emph{niveau de vie} (revenu national brut par habitant). Il varie entre 0 et 1 : plus il est proche de 1, plus le développement humain est élevé.

\begin{savaistu}[D'où vient l'IDH ?]
Le concept d'IDH a été créé en 1990, dans le cadre du Programme des Nations Unies pour le Développement (PNUD), par l'économiste pakistanais Mahbub ul Haq, avec la contribution décisive d'Amartya Sen, prix Nobel d'économie en 1998. Leur idée révolutionnaire : « les gens sont la vraie richesse d'une nation ». Le développement ne se mesure pas d'abord en dollars, mais en capacités humaines — vivre longtemps, être instruit, disposer d'un revenu décent.
\end{savaistu}

\textbf{3. Les indicateurs complémentaires.} L'\cle{espérance de vie} à la naissance traduit l'état sanitaire d'un pays (environ 60 ans au Cameroun, plus de 82 ans en France). Le \cle{taux d'alphabétisation} mesure la part des adultes sachant lire et écrire. On utilise aussi le taux de scolarisation, l'accès à l'eau potable ou à l'électricité.

\begin{exemplebox}[Le Cameroun en chiffres]
L'IDH du Cameroun est d'environ 0,58, ce qui place le pays autour du 150\textsuperscript{e} rang mondial, dans la catégorie des pays à « développement humain moyen ». La France affiche un IDH voisin de 0,90. Le Cameroun est classé parmi les pays à \emph{revenu intermédiaire}, mais cette moyenne cache de fortes inégalités internes : les villes de Douala et Yaoundé concentrent les services, les industries et les emplois, tandis que de nombreuses zones rurales manquent encore d'électricité, de routes praticables et de centres de santé. Un même indicateur national peut donc masquer des réalités très contrastées.
\end{exemplebox}

\begin{popfigure}
\begin{center}
\begin{tikzpicture}
\begin{axis}[
  ybar, bar width=0.35cm,
  width=0.95\linewidth, height=6.2cm,
  ymin=0,
  symbolic x coords={Cameroun,Nigeria,France,Corée du Sud},
  xtick=data,
  x tick label style={font=\small},
  ylabel={Valeur},
  legend style={at={(0.5,-0.18)},anchor=north,legend columns=2,font=\small},
  nodes near coords, every node near coord/.append style={font=\scriptsize},
  axis lines=left, enlarge x limits=0.18
]
\addplot+[ybar,fill=popblue] coordinates {(Cameroun,1.6) (Nigeria,2.2) (France,43) (Corée du Sud,33)};
\addplot+[ybar,fill=poporange] coordinates {(Cameroun,0.58) (Nigeria,0.55) (France,0.90) (Corée du Sud,0.93)};
\legend{PIB/habitant (milliers USD), IDH (échelle 0--1)}
\end{axis}
\end{tikzpicture}
\end{center}
\legende{Comparaison indicative du PIB par habitant (en milliers de dollars, ordres de grandeur) et de l'IDH pour quatre pays. On remarque que le classement par revenu et le classement par IDH se recoupent largement, mais que l'IDH donne une image plus fidèle du bien-être réel (santé et éducation incluses). Sources : ordres de grandeur PNUD/Banque mondiale, années récentes.}
\end{popfigure}

\begin{methode}[Analyser le niveau de développement d'un pays]
Pour éviter les conclusions hâtives, procède toujours en trois étapes :
\begin{enumerate}
\item \textbf{Croiser au moins trois indicateurs} : le revenu (PIB par habitant), la santé (espérance de vie) et l'éducation (taux d'alphabétisation ou de scolarisation). L'IDH les synthétise déjà.
\item \textbf{Interroger les moyennes} : demande-toi qui profite de la richesse. Cherche les inégalités (ville/campagne, hommes/femmes, riches/pauvres).
\item \textbf{Situer dans le temps} : un indicateur isolé ne dit rien ; c'est son évolution sur dix ou vingt ans qui révèle une dynamique de développement ou de stagnation.
\end{enumerate}
\end{methode}

\courssub{Les dimensions du développement}

Le développement est un phénomène \emph{multidimensionnel}. On en distingue classiquement cinq dimensions solidaires :

\begin{itemize}
\item \textbf{Dimension économique} : modernisation de la production, industrialisation, transformation locale des matières premières, infrastructures (routes, ports, énergie).
\item \textbf{Dimension sociale} : réduction de la pauvreté et des inégalités, accès de tous à la santé, à l'eau potable, au logement décent.
\item \textbf{Dimension humaine} : éducation, formation professionnelle, épanouissement des capacités de chacun — c'est le cœur de la pensée d'Amartya Sen.
\item \textbf{Dimension culturelle} : un développement authentique ne copie pas servilement les modèles étrangers ; il s'enracine dans les valeurs, les langues et les créativités propres du peuple. Senghor insistait sur ce point : l'Afrique doit se développer « sans se renier ».
\item \textbf{Dimension environnementale} : c'est l'exigence du \cle{développement durable}, défini par le rapport Brundtland (1987) comme « un développement qui répond aux besoins du présent sans compromettre la capacité des générations futures à répondre aux leurs ». Exploiter la forêt du bassin du Congo sans la reboiser, c'est hypothéquer l'avenir.
\end{itemize}

\begin{aretenir}[L'essentiel]
Le développement est un processus global (économique, social, humain, culturel, environnemental) d'amélioration durable des conditions de vie. Il ne se réduit pas à la croissance du PIB : il se mesure par des indicateurs croisés (PIB/habitant, IDH, espérance de vie, alphabétisation) et se juge à la répartition des richesses autant qu'à leur volume.
\end{aretenir}

\courssub{Le développement endogène : compter d'abord sur soi-même}

Face aux modèles de développement importés — plans conçus à Washington ou à Paris, dépendance à l'aide extérieure —, de nombreux penseurs africains défendent l'idée d'un \cle{développement endogène}.

\begin{definition}[Développement endogène]
Le \cle{développement endogène} est une stratégie qui consiste à se développer en s'appuyant d'abord sur ses propres ressources (matières premières, savoir-faire, marché intérieur) et sur ses propres valeurs culturelles, au lieu d'attendre le salut de l'aide extérieure ou d'imiter des modèles étrangers. « Endogène » signifie littéralement « qui vient de l'intérieur ».
\end{definition}

Cette idée est au cœur de la pensée de Julius Nyerere (Tanzanie), qui prônait l'\emph{Ujamaa} ou socialisme communautaire fondé sur l'entraide villageoise traditionnelle, et de Kwame Nkrumah, pour qui l'unité africaine était la condition d'une industrialisation autonome. L'exemple camerounais de la transformation locale du cacao — fabriquer le chocolat à Douala plutôt qu'exporter les fèves brutes — illustre concrètement ce que serait un développement endogène : la richesse créée reste sur place, forme des techniciens, fait vivre des familles.

\begin{exoresolu}[Croissance ou développement ?]
\textbf{Énoncé.} Un pays A voit son PIB augmenter de 7 \% par an grâce à l'exportation de pétrole brut, mais 40 \% de sa population vit sous le seuil de pauvreté, l'espérance de vie stagne à 55 ans et le taux d'alphabétisation reste à 60 \%. Un pays B affiche une croissance de 4 \%, mais son IDH passe de 0,55 à 0,70 en quinze ans grâce à des investissements massifs dans l'éducation et la santé. Lequel des deux pays connaît un véritable développement ? Justifie ta réponse en t'appuyant sur les notions du cours.
\tcblower
\textbf{Solution détaillée.} Le pays A connaît une forte \emph{croissance économique} (augmentation quantitative du PIB), mais pas un développement : la richesse produite ne se transforme ni en bien-être social (pauvreté massive), ni en progrès humain (espérance de vie et alphabétisation stagnantes). C'est le cas typique du « paradoxe pétrolier ». Le pays B, malgré une croissance plus modeste, connaît un véritable \emph{développement} : la progression de l'IDH de 0,55 à 0,70 atteste une amélioration durable de la santé, de l'éducation et du niveau de vie de la population. Conclusion : conformément à la distinction établie en cours, la croissance est une condition du développement, mais seule la transformation qualitative des structures et du bien-être permet de parler de développement. C'est donc le pays B qui est en voie de développement.
\end{exoresolu}

\begin{exercice}[À toi de jouer]
1. Définis le développement et explique en quoi il diffère de la simple croissance économique.\\
2. Pourquoi le PIB par habitant est-il un indicateur insuffisant ? Donne un exemple concret.\\
3. Le Cameroun a un IDH d'environ 0,58. Quelles sont les trois dimensions prises en compte par cet indice ?\\
4. Rédige un paragraphe argumenté (10 à 15 lignes) sur le sujet : « Le développement endogène est-il une voie réaliste pour l'Afrique ? »
\end{exercice}

\begin{corrige}[Exercice]
1. Le développement est l'ensemble des transformations quantitatives et qualitatives (économiques, sociales, culturelles, politiques) qui améliorent durablement les conditions de vie matérielles et spirituelles d'une société. La croissance n'est que l'augmentation quantitative du PIB ; le développement suppose en plus une transformation qualitative des structures et une amélioration du bien-être de tous.\\
2. Le PIB par habitant est une moyenne qui masque les inégalités de répartition : une richesse concentrée entre quelques mains peut produire une moyenne élevée alors que la majorité vit dans la pauvreté (cas des pays pétroliers). Il ignore aussi la santé, l'éducation et l'environnement.\\
3. L'IDH combine : la santé (espérance de vie à la naissance), l'éducation (durée moyenne et attendue de scolarisation) et le niveau de vie (revenu national brut par habitant).\\
4. Le paragraphe doit définir le développement endogène (s'appuyer d'abord sur ses propres ressources et valeurs), présenter un argument en sa faveur (transformation locale des matières premières, exemple du cacao camerounais ; référence possible à Nkrumah ou Nyerere), une difficulté (manque de capitaux, de technologies, poids des importations) et une conclusion nuancée : voie réaliste à condition d'investir dans l'éducation, l'industrie locale et l'unité africaine — ce qui annonce la suite de la leçon sur le panafricanisme.
\end{corrige}

Ainsi, mesurer le développement exige de dépasser les chiffres bruts pour interroger la vie concrète des populations. C'est précisément parce que l'Afrique, riche de ses ressources, reste marquée par le sous-développement que des penseurs comme Nkrumah, Senghor et Nyerere ont fait du panafricanisme une philosophie du développement : c'est ce que nous étudierons dans les sections suivantes.

\coursec{Le sous-développement de l'Afrique : manifestations et causes}

Dans la section précédente, nous avons défini le \cle{développement} comme un processus global de transformation économique, sociale et culturelle, mesurable par des indicateurs comme le PIB par habitant ou l'IDH. Mais cette définition appelle immédiatement son envers : si certains pays se développent, d'autres semblent enfermés dans le retard. C'est le cas de la plupart des pays africains. Faut-il y voir une simple étape, un « retard à rattraper », ou bien un état produit par des mécanismes précis ? C'est toute la question du \cle{sous-développement}, que nous abordons maintenant en trois temps : sa définition, ses manifestations visibles au quotidien, puis ses causes, externes et internes.

\courssub{Définition du sous-développement : un état, mais surtout un processus}

Commençons par une intuition simple. Deux enfants naissent le même jour, l'un à Yaoundé, l'autre à Oslo. Le second a statistiquement plus de chances d'être soigné, scolarisé longtemps, et de vivre plus de 80 ans ; le premier risque davantage la malnutrition, le chômage ou l'émigration forcée. Cette inégalité de départ n'a rien de « naturel » : elle traduit un état de fait que les économistes appellent le sous-développement.

\begin{definition}[Le sous-développement]
Le \cle{sous-développement} désigne un état de retard économique et social (faiblesse des revenus, de l'industrialisation, des services de santé et d'éducation) résultant de causes à la fois \textbf{internes} et \textbf{externes}, et non d'une quelconque infériorité naturelle des peuples concernés.
\end{definition}

Mais la définition la plus profonde est celle de l'économiste \textbf{André Gunder Frank}. Dans sa théorie dite de la \emph{dépendance}, Frank soutient que le sous-développement n'est pas un « retard » que les pays pauvres rattraperaient en copiant les pays riches : c'est un \textbf{processus historiquement produit}. Autrement dit, le sous-développement de la « périphérie » (l'Afrique, l'Amérique latine) et le développement du « centre » (l'Europe, l'Amérique du Nord) sont les deux faces d'un même mouvement : l'insertion inégale des continents dans l'économie mondiale. Le sous-développement n'est donc pas l'absence de développement ; il est le \emph{produit} d'un certain développement mondial.

\begin{aretenir}[L'idée-force de Frank]
On ne naît pas sous-développé, on le devient. Le sous-développement est le résultat de rapports historiques de domination (traite, colonisation, échanges inégaux), auxquels s'ajoutent des facteurs internes. Il est donc, en principe, \textbf{réversible}.
\end{aretenir}

\courssub{Les manifestations du sous-développement au quotidien : l'exemple camerounais}

Le sous-développement n'est pas une abstraction : il se voit et se vit. Menons l'enquête comme le ferait un sociologue, en observant la réalité camerounaise.

\begin{itemize}
\item \textbf{La pauvreté.} Une part importante de la population vit avec moins de deux dollars par jour. Dans les quartiers populaires de Douala ou les villages de l'Extrême-Nord, l'accès à l'eau potable reste un combat quotidien.
\item \textbf{Le chômage des jeunes.} Des milliers de diplômés — techniciens, ingénieurs, licenciés — ne trouvent pas d'emploi correspondant à leur formation. Beaucoup survivent grâce au secteur informel (moto-taxis, petit commerce), faute d'industries pour les absorber.
\item \textbf{L'exode rural.} Les jeunes quittent les villages pour les villes, puis parfois pour l'étranger, vidant les campagnes de leurs forces vives et grossissant des bidonvilles surpeuplés.
\item \textbf{L'insuffisance des infrastructures.} Routes dégradées qui isolent les zones de production agricole, hôpitaux sous-équipés, coupures d'électricité qui freinent les entreprises : le manque d'infrastructures bloque toute dynamique économique.
\item \textbf{La dépendance alimentaire.} Paradoxalement, un pays au sol fertile comme le Cameroun importe du riz, du poisson congelé ou de la farine de blé, alors qu'il pourrait nourrir sa population et exporter.
\end{itemize}

\begin{exemplebox}[Observer pour comprendre]
Un jeune planteur de cacao de la région du Centre récolte 500 kg de fèves. Il les vend brutes à un acheteur qui les exporte vers l'Europe. Là-bas, ces fèves deviennent du chocolat vendu dix à vingt fois plus cher — et une partie revient au Cameroun sous forme de tablettes importées. Le planteur reste pauvre non parce qu'il travaille mal, mais parce qu'il est enfermé dans le maillon le moins rentable de la chaîne : la production brute.
\end{exemplebox}

\courssub{Les causes externes du sous-développement}

Première piste de notre investigation : l'histoire des relations entre l'Afrique et le reste du monde.

\textbf{1. La traite négrière.} Pendant près de quatre siècles, des millions d'Africains en âge de travailler ont été déportés. Le continent a perdu une part considérable de sa force de travail, ses structures sociales ont été désorganisées, tandis que cette main-d'œuvre enrichissait les plantations d'Amérique et les ports d'Europe.

\textbf{2. La colonisation.} De la fin du XIX\textsuperscript{e} siècle aux indépendances, les puissances européennes ont organisé l'économie africaine pour leurs propres besoins : cultures d'exportation (cacao, café, coton) au détriment des cultures vivrières, routes construites pour évacuer les matières premières vers les ports et non pour relier les régions entre elles, main-d'œuvre exploitée à bas prix. L'Afrique indépendante a hérité d'économies tournées vers l'extérieur et déformées.

\textbf{3. Le partage inégal du commerce mondial.} Les pays africains vendent des matières premières dont les prix, fixés sur les marchés de Londres ou de New York, chutent facilement, et achètent des produits manufacturés dont les prix montent : c'est la \emph{détérioration des termes de l'échange}.

\textbf{4. L'endettement.} Les pays africains, endettés auprès des institutions financières internationales, consacrent une part importante de leurs budgets au remboursement de la dette, au détriment des écoles et des hôpitaux. Les plans d'ajustement structurel des années 1980-1990 ont souvent aggravé la situation sociale.

\textbf{5. L'exportation de matières premières brutes.} Sans transformation locale, la plus grande partie de la valeur ajoutée — et des emplois — échappe au continent.

\begin{savaistu}[Walter Rodney : l'Europe s'est développée contre l'Afrique]
Dans son ouvrage \emph{Comment l'Europe a sous-développé l'Afrique} (1972), l'historien guyanais \textbf{Walter Rodney} soutient une thèse radicale : le développement de l'Europe et le sous-développement de l'Afrique sont liés comme cause et effet. La traite et la colonisation ont transféré les richesses et les forces vives de l'Afrique vers l'Europe. Pour Rodney, comprendre le sous-développement exige donc de penser l'histoire des rapports de domination, et non d'accuser les Africains eux-mêmes.
\end{savaistu}

\courssub{Les causes internes du sous-développement}

Une analyse honnête ne peut pas tout mettre sur le compte de l'extérieur. L'investigation doit aussi se retourner vers l'intérieur.

\begin{itemize}
\item \textbf{La mauvaise gouvernance.} Détournements de fonds publics, administration inefficace, absence de planification à long terme : l'argent qui devrait financer routes et écoles s'évapore parfois.
\item \textbf{La corruption.} Elle augmente le coût de tout (marchés publics gonflés, pots-de-vin) et décourage les investisseurs sérieux.
\item \textbf{Les conflits.} Guerres civiles et instabilité politique détruisent les infrastructures, font fuir les capitaux et interrompent la production.
\item \textbf{La faible industrialisation.} Faute d'usines de transformation, les pays africains restent des fournisseurs de matières brutes et des importateurs de produits finis.
\item \textbf{La fuite des cerveaux.} Médecins, ingénieurs et chercheurs formés en Afrique — souvent aux frais de l'État — partent exercer en Europe ou en Amérique du Nord, privant le continent de ses compétences les plus précieuses.
\end{itemize}

\begin{attention}[Erreur fréquente à l'examen]
Attribuer le sous-développement à une \textbf{seule} cause — la colonisation seule, ou la gouvernance seule — est une faute d'analyse. Le correcteur attend une démarche \textbf{équilibrée} : montrer les causes externes (héritage historique, échanges inégaux) \emph{et} les causes internes (gouvernance, corruption), puis discuter leur articulation. Un devoir manichéen est un devoir faible.
\end{attention}

\courssub{Le paradoxe africain : un continent riche mais des populations pauvres}

Voici le fait le plus déroutant de notre enquête : l'Afrique est immensément riche. Le cobalt de la République démocratique du Congo est indispensable aux batteries du monde entier ; le pétrole du Nigeria et de l'Angola alimente l'économie mondiale ; le cacao de la Côte d'Ivoire et du Cameroun fait le chocolat de la planète ; l'or du Ghana et du Mali brille partout. Et pourtant, les revenus des populations restent parmi les plus faibles du monde.

\begin{exemplebox}[L'exemple chiffré à retenir]
L'Afrique détient environ \textbf{30 \%} des réserves mondiales de minerais, mais ne capte qu'environ \textbf{3 \%} du commerce mondial. Cet écart spectaculaire entre la richesse du sol et la pauvreté des revenus est le cœur du paradoxe africain : la richesse existe, mais sa valeur est captée ailleurs, dans les usines et les bourses des pays industrialisés.
\end{exemplebox}

\begin{popfigure}
\begin{center}
\begin{tikzpicture}[scale=0.9]
% Contour très simplifié de l'Afrique
\draw[thick, popblue, fill=popcream]
  (0.2,4.6) -- (2.2,5.2) -- (4.2,4.8) -- (5.6,4.2) -- (6.2,3.2)
  -- (5.6,2.2) -- (5.2,1.0) -- (4.4,-0.4) -- (3.6,-1.2) -- (3.0,-1.6)
  -- (2.4,-1.0) -- (1.8,0.2) -- (1.2,1.4) -- (0.6,2.4) -- (-0.2,3.4)
  -- cycle;
% Ressources (points légendés)
\fill[popdark] (1.0,3.6) circle (0.09); % pétrole Nigeria/Golfe de Guinée
\fill[popdark] (2.6,4.4) circle (0.09); % pétrole Sahara
\fill[poppurple] (3.4,1.6) circle (0.09); % cobalt RDC
\fill[popgold] (1.6,2.6) circle (0.09); % or Ghana/Mali
\fill[popgold] (3.4,-0.6) circle (0.09); % or Afrique du Sud
\fill[popgreen] (1.4,2.2) circle (0.09); % cacao Côte d'Ivoire/Cameroun
\fill[popgreen] (2.0,2.0) circle (0.09);
% Nord
\draw[->, popmuted] (6.6,4.6) -- (6.6,5.4) node[above]{\scriptsize N};
\end{tikzpicture}
\end{center}
\legende{Carte très simplifiée de l'Afrique localisant quelques ressources majeures : \textcolor{popdark}{$\bullet$} pétrole (Nigeria, golfe de Guinée, Sahara) ; \textcolor{poppurple}{$\bullet$} cobalt (RDC) ; \textcolor{popgold}{$\bullet$} or (Ghana, Mali, Afrique du Sud) ; \textcolor{popgreen}{$\bullet$} cacao (Côte d'Ivoire, Cameroun). Croquis indicatif, sans prétention à la précision cartographique.}
\end{popfigure}

\begin{popfigure}
\begin{center}
\begin{tikzpicture}
\begin{axis}[
  ybar, bar width=0.55cm,
  width=0.85\linewidth, height=5.5cm,
  ymin=0, ymax=45,
  ylabel={Part du commerce mondial (en \%)},
  symbolic x coords={Afrique, Asie, Europe, Am\'erique},
  xtick=data,
  nodes near coords, nodes near coords align={vertical},
  every node near coord/.append style={font=\scriptsize},
  ymajorgrids=true, grid style={popline!40},
]
\addplot[fill=popblue, draw=popblue!70!black] coordinates {(Afrique,3) (Asie,32) (Europe,38) (Am\'erique,20)};
\end{axis}
\end{tikzpicture}
\end{center}
\legende{Part approximative des continents dans le commerce mondial des marchandises (ordres de grandeur, en \%). L'Afrique, malgré ses ressources, ne pèse qu'environ 3 \% : elle vend brut et achète transformé.}
\end{popfigure}

\courssub{Débat philosophique : fatalité ou domination réversible ?}

Nous pouvons maintenant formuler le problème philosophique de cette section : \emph{le sous-développement de l'Afrique est-il une fatalité, ou le produit de rapports de domination que l'on peut renverser ?}

\textbf{Première thèse : la fatalité.} Certains discours, anciens ou actuels, présentent le sous-développement comme une donnée quasi naturelle : climat difficile, « mentalités » rétrogrades, démographie galopante. Cette thèse est philosophiquement suspecte : elle naturalise un état historique, elle déresponsabilise les acteurs (puissances étrangères comme élites locales) et elle décourage toute action. Or aucune loi de la nature n'assigne un continent à la pauvreté : des pays comme la Corée du Sud, très pauvre en 1960, ont montré que des trajectoires peuvent changer.

\textbf{Deuxième thèse : la domination réversible.} Pour Frank et Rodney, le sous-développement est un \emph{rapport}, non une essence : il résulte de mécanismes historiques (traite, colonisation, échanges inégaux) et de choix politiques (gouvernance, politiques économiques). Ce qui a été produit par des hommes peut être transformé par des hommes : transformation locale des matières premières, intégration régionale, bonne gouvernance, investissement dans l'éducation. C'est précisément le pari du \cle{panafricanisme}, que nous étudierons dans la suite de la leçon.

\begin{methode}[Analyser un sujet sur le sous-développement]
Face à un sujet de dissertation portant sur le sous-développement de l'Afrique, structurez toujours votre analyse en trois temps :
\begin{enumerate}
\item \textbf{Les manifestations} : décrire concrètement le phénomène (pauvreté, chômage, infrastructures, dépendance alimentaire) avec des exemples camerounais ;
\item \textbf{Les causes externes} : traite, colonisation, échanges inégaux, endettement (pensez à Frank et Rodney) ;
\item \textbf{Les causes internes} : gouvernance, corruption, conflits, fuite des cerveaux.
\end{enumerate}
Terminez par le débat (fatalité ou réversibilité) pour ouvrir vers une position argumentée.
\end{methode}

\begin{exoresolu}[Exercice d'application]
\textbf{Énoncé.} Un élève déclare : « L'Afrique est pauvre parce que les Africains ne travaillent pas assez. » En vous appuyant sur la thèse d'André Gunder Frank et sur l'exemple du cacao camerounais, montrez en quelques lignes pourquoi cette affirmation est philosophiquement insuffisante.
\tcblower
\textbf{Solution détaillée.} L'affirmation repose sur une erreur : elle explique un état historique par une prétendue nature (la « paresse »), ce qui est un préjugé et non une analyse. Or le planteur de cacao camerounais travaille dur : il récolte ses 500 kg de fèves et les vend. Sa pauvreté ne vient pas de son manque de travail, mais de sa \emph{position} dans l'économie mondiale : il vend brut, sans transformation, à un prix qu'il ne fixe pas, tandis que la valeur ajoutée du chocolat est captée en Europe. C'est exactement la thèse de Frank : le sous-développement n'est pas un retard naturel mais un \emph{processus historiquement produit} par l'insertion inégale de l'Afrique dans le commerce mondial — auquel s'ajoutent des causes internes (absence d'usines de transformation, gouvernance). Conclusion : le sous-développement s'explique par des rapports économiques et historiques, réversibles, et non par une infériorité des peuples.
\end{exoresolu}

\begin{exercice}[Pour s'entraîner]
Sujet de réflexion : « Le sous-développement de l'Afrique est-il une fatalité ? » Construisez le plan détaillé d'une dissertation (introduction problématisée, deux ou trois parties, conclusion) en mobilisant les notions de cette section et les auteurs cités (Frank, Rodney).
\end{exercice}

\begin{corrige}[Exercice]
\textbf{Introduction.} Amorce : le contraste entre la richesse des ressources africaines (30 \% des minerais mondiaux) et la pauvreté des revenus (3 \% du commerce mondial). Problématique : le sous-développement est-il un destin inscrit dans la nature du continent, ou le résultat de rapports que l'on peut transformer ? Annonce du plan.

\textbf{I. Le sous-développement semble une fatalité.} Manifestations massives et persistantes (pauvreté, chômage, infrastructures) ; discours qui naturalisent la pauvreté ; poids de l'histoire (traite, colonisation).

\textbf{II. Mais il est un processus historiquement produit.} Thèse de Frank (le sous-développement comme produit de l'échange inégal) ; thèse de Rodney (le développement de l'Europe contre l'Afrique) ; rôle des causes internes (gouvernance, corruption, fuite des cerveaux) : ce qui est produit par des choix peut être changé par d'autres choix.

\textbf{III. Vers un dépassement : la voie panafricaine.} Transformation locale, unité africaine, bonne gouvernance : le sous-développement est réversible, ce qui ouvre sur la suite de la leçon.

\textbf{Conclusion.} Le sous-développement n'est pas une fatalité mais un rapport de domination réversible ; dire qu'il est réversible n'est pas dire qu'il est facile à vaincre, mais que l'action collective a un sens.
\end{corrige}

Ainsi, le sous-développement de l'Afrique apparaît comme un état de retard économique et social, mais surtout comme le produit d'une histoire de domination combinée à des défaillances internes. Cette double causalité — et l'idée que ce qui a été fait peut être défait — est le socle sur lequel s'est construite la grande réponse politique et philosophique du continent : le panafricanisme, dont nous étudierons l'historique dans la prochaine section.

\coursec{Historique du panafricanisme}

Dans la section précédente, nous avons montr\'e que le sous-d\'eveloppement de l'Afrique n'est pas une fatalit\'e : il r\'esulte de causes historiques (traite n\'egri\`ere, colonisation) et structurelles (d\'ependance \'economique, faiblesse de l'industrie, morcellement des march\'es). Une question s'impose alors : comment l'Afrique peut-elle s'en sortir ? La r\'eponse que des g\'en\'erations de penseurs et d'hommes d'\'Etat africains ont donn\'ee tient en un mot : \emph{l'unit\'e}. C'est l'id\'ee-force du \cle{panafricanisme}, dont nous allons retracer l'histoire. Comprendre cet historique, c'est comprendre comment une id\'ee n\'ee dans la souffrance de la diaspora est devenue un projet politique continental, de Paris 1919 \`a la ZLECAf de 2021.

\begin{definition}[Le panafricanisme]
Le \cle{panafricanisme} est un mouvement politique, culturel et philosophique visant l'\cle{unit\'e} et la \cle{solidarit\'e} de tous les Africains, ceux du continent comme ceux de la \cle{diaspora} (descendants des Africains d\'eport\'es en Am\'erique, aux Cara\"ibes et en Europe). Le pr\'efixe grec \emph{pan} signifie « tout » : le panafricanisme affirme que tous les Noirs forment une m\^eme communaut\'e de destin et que l'\'emancipation et le d\'eveloppement de l'Afrique passent par leur union.
\end{definition}

\courssub{Les origines dans la diaspora : de la r\'evolte des esclaves au « Back to Africa »}

\textbf{L'intuition de d\'epart.} Le panafricanisme ne na\^it pas en Afrique, mais dans les Am\'eriques, chez les descendants d'esclaves. Pourquoi ? Parce que la d\'eportation a m\'elang\'e des millions d'Africains de toutes ethnies dans une m\^eme condition : r\'eduis en esclavage, ils ont d\'ecouvert qu'aux yeux du monde ils \'etaient d'abord des « Noirs », au-del\`a de leurs appartenances particuli\`eres. La conscience d'une identit\'e commune est n\'ee de l'oppression commune.

\textbf{Les r\'evoltes d'esclaves.} D\`es le XVII\textsuperscript{e} si\`ecle, les r\'evoltes d'esclaves (marrons de la Jama\"ique, r\'evoltes au Br\'esil) expriment le refus de la d\'eshumanisation. L'\'ev\'enement fondateur est la r\'evolution de Saint-Domingue (Ha\"iti, 1791-1804) : men\'ee par Toussaint Louverture, elle aboutit \`a la premi\`ere r\'epublique noire ind\'ependante du monde. Elle prouve que les Africains d\'eport\'es peuvent se lib\'erer eux-m\^emes, et elle devient un symbole pour toute la diaspora.

\textbf{Le mouvement « Back to Africa » de Marcus Garvey.} Au d\'ebut du XX\textsuperscript{e} si\`ecle, le Jama\"icain \cle{Marcus Garvey} (1887-1940) fonde l'UNIA (\emph{Universal Negro Improvement Association}) et lance le mouvement « \cle{Back to Africa} » (« Retour en Afrique »). Son id\'ee : les Noirs ne seront jamais respect\'es en Am\'erique ; ils doivent retourner en Afrique et y b\^atir une nation forte. Garvey cr\'ee une compagnie maritime, la \emph{Black Star Line}, pour concr\'etiser ce retour. Son mouvement rassemble des millions d'adh\'erents et diffuse la fiert\'e noire (« \emph{Black is beautiful} » avant la lettre). M\^eme si le projet de retour massif \'echoue, Garvey a pos\'e le principe : la lib\'eration des Noirs passe par la lib\'eration de l'Afrique.

\textbf{Les p\`eres fondateurs.} Deux intellectuels structurent le mouvement :
\begin{itemize}
\item \cle{Henry Sylvester Williams}, avocat de Trinidad, organise en 1900 \`a Londres la premi\`ere \emph{Conf\'erence panafricaine} : il invente le mot m\^eme de « panafricanisme » et d\'enonce les violences coloniales.
\item \cle{W.E.B. Du Bois} (1868-1963), sociologue am\'ericain, premier Afro-Am\'ericain docteur de Harvard, devient le v\'ritable organisateur des congr\`es panafricains. Sa formule c\'el\`ebre : « le probl\`eme du XX\textsuperscript{e} si\`ecle est le probl\`eme de la ligne de couleur ». Il finira sa vie au Ghana, pays de Nkrumah, symbole du lien entre diaspora et continent.
\end{itemize}

\begin{attention}[Pi\`ege fr\'equent]
Ne pas confondre \cle{panafricanisme} et \cle{n\'egritude}. Le panafricanisme est un mouvement d'\emph{id\'ees politiques} visant l'unit\'e et la solidarit\'e de tous les Africains ; la n\'egritude est un mouvement \emph{culturel et litt\'eraire} de revalorisation de la culture noire. La n\'egritude est une \emph{composante} du panafricanisme, non son \'equivalent. Dire « le panafricanisme est un mouvement litt\'eraire » est une erreur de copie.
\end{attention}

\courssub{Les congr\`es panafricains : de Paris 1919 \`a Manchester 1945}

Entre 1919 et 1945, cinq congr\`es panafricains se succ\`edent. Ils marquent le passage d'une simple solidarit\'e raciale \`a une revendication politique claire : l'ind\'ependance.

\textbf{Le congr\`es de Paris (1919).} Organis\'e par Du Bois en marge de la conf\'erence de paix qui suit la Premi\`ere Guerre mondiale, il r\'unit 57 d\'el\'egu\'es. Ses revendications restent mod\'er\'ees : il demande que les colonies allemandes ne soient pas partag\'ees entre les vainqueurs et que les Africains soient mieux trait\'es. On ne parle pas encore d'ind\'ependance, mais d'une administration plus juste. Les congr\`es suivants (Londres 1921, Londres-Lisbonne 1923, New York 1927) restent dans cette ligne r\'eformiste.

\textbf{Le congr\`es de Manchester (1945) : le tournant.} Le 5\textsuperscript{e} congr\`es, tenu \`a Manchester en octobre 1945, change tout. Le contexte : la Seconde Guerre mondiale vient de s'achever ; des soldats africains ont combattu pour la libert\'e de la France et de l'Angleterre ; la charte de l'ONU proclame le droit des peuples \`a disposer d'eux-m\^emes. Les jeunes leaders africains pr\'esents — le Ghan\'een \cle{Kwame Nkrumah}, le K\'enyan \cle{Jomo Kenyatta} — imposent un ton nouveau : fini les demandes polies, place \`a l'exigence.

\begin{exemplebox}[Manchester 1945 : l'exigence d'ind\'ependance]
Le 5\textsuperscript{e} congr\`es panafricain de Manchester (octobre 1945) adopte une r\'esolution exigeant l'\cle{ind\'ependance imm\'ediate} des colonies africaines. Les d\'el\'egu\'es d\'eclarent : « Nous sommes d\'etermin\'es \`a \^etre libres » et appellent, si n\'ecessaire, \`a la d\'esob\'eissance civile et \`a la gr\`eve. Quinze ans plus tard, Nkrumah devient pr\'esident du Ghana ind\'ependant (1957) et Kenyatta premier pr\'esident du Kenya (1963) : les revendications de Manchester sont devenues r\'ealit\'e.
\end{exemplebox}

\courssub{La n\'egritude : la dimension culturelle du combat}

Parall\`elemment aux congr\`es, un mouvement litt\'eraire donne au panafricanisme sa dimension culturelle : la \cle{n\'egritude}. Le mot est forg\'e par le Martiniquais \cle{Aim\'e C\'esaire} dans son \emph{Cahier d'un retour au pays natal} (1939). Autour de lui, le S\'en\'egalais \cle{L\'eopold S\'edar Senghor} et le Guyanais \cle{L\'eon-Gontran Damas} fondent \`a Paris, dans les ann\'ees 1930, une po\'esie de la fiert\'e noire.

\textbf{L'id\'ee centrale.} La colonisation a enseign\'e aux Noirs \`a m\'epriser leur propre culture, pr\'esent\'ee comme « primitive ». La n\'egritude renverse la perspective : elle affirme que la culture noire — ses valeurs de communaut\'e, de rythme, de relation vivante \`a la nature — a sa dignit\'e propre et peut apporter sa contribution \`a la civilisation universelle. Senghor le r\'esume : « L'\'emotion est n\`egre, comme la raison est hell\`ene » — formule discutable, mais qui exprime la volont\'e de r\'ehabiliter ce que le colonisateur d\'evalorisait.

\textbf{Le lien avec le d\'eveloppement.} La n\'egritude n'est pas un simple folklore : elle restaure la \emph{confiance en soi}, condition psychologique de tout d\'eveloppement. Un peuple convaincu de son inf\'eriorit\'e ne peut pas se d\'evelopper ; un peuple fier de son identit\'e peut entreprendre. C'est en ce sens que la n\'egritude est une composante du panafricanisme comme philosophie du d\'eveloppement.

\courssub{Les ind\'ependances et le d\'ebat Casablanca--Monrovia}

\textbf{L'ann\'ee 1960.} \`A partir de 1957 (Ghana), les colonies africaines acc\`edent massivement \`a l'ind\'ependance ; l'ann\'ee 1960, dite « \emph{ann\'ee de l'Afrique} », voit na\^itre dix-sept \'Etats ind\'ependants, dont le Cameroun (1\textsuperscript{er} janvier 1960). Mais l'ind\'ependance politique pose aussit\^ot une question : que faire de cette libert\'e retrouv\'ee ? Deux strat\'egies s'affrontent.

\begin{center}
\begin{tabularx}{\linewidth}{|l|X|X|}
\hline
\rowcolor{popblueL} & \textbf{Groupe de Casablanca} & \textbf{Groupe de Monrovia} \\
\hline
Pays principaux & Ghana, Guin\'ee, Mali, Maroc, \'Egypte, Alg\'erie & Nigeria, Lib\'eria, \'Ethiopie, la plupart des \'Etats francophones \\
\hline
Leader embl\'ematique & Kwame Nkrumah & Chefs d'\'Etat mod\'er\'es (Ha\"il\'e S\'elass\'e, Tubman) \\
\hline
Doctrine & \cle{Unit\'e imm\'ediate} : cr\'eer tout de suite les « \'Etats-Unis d'Afrique », avec gouvernement, arm\'ee et monnaie communes & \cle{Coop\'eration progressive} : respecter la souverainet\'e de chaque \'Etat, coop\'erer \'etape par \'etape sans fusion politique \\
\hline
Argument cl\'e & S\'epar\'es, les \'Etats africains restent faibles et d\'ependants ; seule l'unit\'e rend l'ind\'ependance r\'eelle & Chaque \'Etat vient de conqu\'erir sa souverainet\'e ; il est irr\'ealiste de la dissoudre aussit\^ot \\
\hline
\end{tabularx}
\end{center}

Ce d\'ebat est philosophiquement important : il oppose une conception \emph{f\'ed\'eraliste} de l'unit\'e (Nkrumah) \`a une conception \emph{inter\'etatique} et prudente. L'histoire donnera raison, provisoirement, aux partisans de Monrovia — mais les \'etapes r\'ecentes (ZLECAf) montrent que l'intuition de Nkrumah n'a jamais disparu.

\courssub{De l'OUA \`a l'Union africaine}

\textbf{La cr\'eation de l'OUA (1963).} Le 25 mai 1963, \`a \cle{Addis-Abeba} (\'Ethiopie), 32 chefs d'\'Etat africains signent la charte de l'\cle{Organisation de l'unit\'e africaine} (OUA). C'est un compromis : l'unit\'e est affirm\'ee, mais la charte consacre la souverainet\'e des \'Etats et le principe de \emph{non-ing\'erence} dans les affaires int\'erieures. Ses objectifs : achever la d\'ecolonisation (lutte contre l'apartheid en Afrique du Sud), promouvoir la coop\'eration, d\'efendre l'int\'egrit\'e territoriale des \'Etats. Le 25 mai devient la « Journ\'ee de l'Afrique ».

\textbf{Bilan et transformation.} L'OUA r\'eussit la d\'ecolonisation totale du continent (apartheid aboli en 1994), mais elle \'echoue sur le plan \'economique : faute de pouvoir contraignant, elle ne peut ni emp\^echer les conflits ni imposer l'int\'egration. D'o\`u sa transformation : le 9 juillet 2002, \`a Durban (Afrique du Sud), l'OUA devient l'\cle{Union africaine} (UA), sur le mod\`ele de l'Union europ\'eenne. L'UA dispose d'organes plus int\'egr\'es : Commission, Parlement panafricain, Cour africaine de justice, Conseil de paix et de s\'ecurit\'e. Elle passe du principe de non-ing\'erence \`a celui de « non-indiff\'erence » face aux crises.

\begin{savaistu}[Le r\^eve audacieux de Nkrumah]
D\`es 1963, lors du sommet fondateur de l'OUA, Kwame Nkrumah proposait un v\'ritable \emph{gouvernement continental} : une arm\'ee africaine commune, une monnaie unique africaine, un march\'e unifi\'e. Ses pairs jug\`erent le projet pr\'ematur\'e. Soixante ans plus tard, la ZLECAf et les discussions sur une monnaie panafricaine montrent que cette vision, longtemps trait\'ee d'utopie, inspire toujours l'int\'egration africaine.
\end{savaistu}

\courssub{Les \'etapes r\'ecentes : NEPAD et ZLECAf}

\textbf{Le NEPAD (2001).} Le \cle{Nouveau Partenariat pour le d\'eveloppement de l'Afrique} (NEPAD), adopt\'e en 2001 et int\'egr\'e \`a l'UA, marque un changement de m\'ethode : les Africains prennent eux-m\^emes en charge leur d\'eveloppement (infrastructures, agriculture, \'education, bonne gouvernance) en \'echange d'un partenariat renouvel\'e avec les pays riches. Le NEPAD introduit m\^eme un m\'ecanisme d'\'evaluation par les pairs : les \'Etats africains s'\'evaluent mutuellement en mati\`ere de gouvernance.

\textbf{La ZLECAf (2019-2021).} La \cle{Zone de libre-\'echange continentale africaine} (ZLECAf), dont l'accord est sign\'e \`a Kigali en 2018, entre en vigueur en 2019 et devient \emph{op\'erationnelle le 1\textsuperscript{er} janvier 2021}. Objectif : cr\'eer un march\'e unique de plus de 1,3 milliard de consommateurs en supprimant progressivent les droits de douane entre pays africains. Enjeu d\'ecisif : aujourd'hui, les \'echanges \emph{entre} pays africains ne repr\'esentent qu'environ 15 \% du commerce total du continent (contre pr\`es de 60 \% en Europe) ; chacun commercerait davantage avec l'ext\'erieur qu'avec ses voisins. La ZLECAf veut inverser cette logique de d\'ependance h\'erit\'ee de la colonisation. C'est la traduction \'economique concr\`ete du vieux r\^eve panafricain.

\begin{popfigure}
\begin{center}
\begin{tikzpicture}[scale=1]
\draw[very thick,popdark,->] (0,0) -- (13.2,0);
\foreach \x/\annee/\evt/\col in {
  0.6/1919/{Congr\`es de Paris\\ (Du Bois)}/popblue,
  3.1/1945/{Congr\`es de Manchester\\ (Nkrumah, Kenyatta)}/poporange,
  5.6/1960/{Ann\'ee de l'Afrique\\ (17 ind\'ependances)}/popgreen,
  8.1/1963/{Cr\'eation de l'OUA\\ (Addis-Abeba)}/poppurple,
  10.4/2002/{OUA $\rightarrow$\\ Union africaine}/poppink,
  12.5/2021/{ZLECAf\\ op\'erationnelle}/popgold}
{
  \fill[\col] (\x,0) circle (0.09);
  \draw[\col,thick] (\x,0) -- (\x,0.55);
  \node[above,align=center,font=\small,text=\col] at (\x,0.6) {\textbf{\annee}\\ \evt};
}
\node[below,font=\small\itshape,text=popmuted] at (6.6,-0.35) {Du congr\`es des id\'ees au march\'e commun : un si\`ecle de panafricanisme};
\end{tikzpicture}
\end{center}
\legende{Frise chronologique du panafricanisme : 1919 (Paris) $\rightarrow$ 1945 (Manchester) $\rightarrow$ 1960 (ind\'ependances) $\rightarrow$ 1963 (OUA) $\rightarrow$ 2002 (Union africaine) $\rightarrow$ 2021 (ZLECAf).}
\end{popfigure}

\begin{popfigure}
\begin{center}
\begin{tikzpicture}[scale=0.62]
% Contour tr\`es simplifi\'e de l'Afrique
\draw[very thick,popdark,fill=popcream] (2.2,7.6) -- (4.6,8.6) -- (7.2,8.2) -- (9.4,7.0) -- (10.6,5.6) -- (10.2,3.6) -- (9.0,1.6) -- (7.4,-0.6) -- (6.4,-1.8) -- (5.4,-0.6) -- (4.6,1.2) -- (3.4,2.2) -- (2.0,2.8) -- (0.8,3.4) -- (0.2,4.6) -- (1.0,6.0) -- cycle;
% Blocs r\'egionaux
\fill[popblue!45] (1.0,5.2) -- (3.6,5.4) -- (3.8,3.4) -- (1.4,3.2) -- (0.4,4.2) -- cycle;
\fill[poporange!50] (3.8,4.6) -- (6.2,4.6) -- (6.2,2.4) -- (4.2,2.2) -- (3.8,3.4) -- cycle;
\fill[popgreen!45] (4.6,2.0) -- (7.6,2.0) -- (7.2,-0.4) -- (6.3,-1.5) -- (5.5,-0.5) -- (4.7,1.1) -- cycle;
\node[font=\small\bfseries,text=popblue] at (2.0,4.4) {CEDEAO};
\node[font=\small\bfseries,text=poporange!80!black] at (5.0,3.5) {CEMAC};
\node[font=\small\bfseries,text=popgreen!60!black] at (6.1,0.8) {SADC};
% Villes rep\`eres
\fill[popdark] (4.9,3.0) circle (0.12); \node[right,font=\scriptsize] at (5.05,3.0) {Yaound\'e};
\fill[popdark] (9.6,4.4) circle (0.12); \node[right,font=\scriptsize] at (9.75,4.4) {Addis-Abeba};
\fill[popdark] (6.3,-1.3) circle (0.12); \node[right,font=\scriptsize] at (6.45,-1.3) {Le Cap};
% Nord
\draw[->,thick,popdark] (11.6,6.6) -- (11.6,7.8); \node[above,font=\small\bfseries] at (11.6,7.8) {N};
\end{tikzpicture}
\end{center}
\legende{Carte sch\'ematique de l'Afrique et de ses principaux blocs r\'egionaux : CEDEAO (Afrique de l'Ouest), CEMAC (Afrique centrale, dont le Cameroun), SADC (Afrique australe). Ces communaut\'es \'economiques r\'egionales sont les « briques » de l'int\'egration continentale voulue par l'Union africaine (si\`ege : Addis-Abeba). Croquis indicatif, sans pr\'ecision cartographique.}
\end{popfigure}

\begin{methode}[M\'emoriser l'historique : 4 dates-cl\'es]
Pour retenir l'essentiel de l'historique du panafricanisme, fixez quatre dates rep\`eres, puis rattachez les autres \'ev\'enements autour :
\begin{enumerate}
\item \textbf{1945} : congr\`es de Manchester — le panafricanisme exige l'ind\'ependance (Nkrumah, Kenyatta).
\item \textbf{1963} : cr\'eation de l'OUA \`a Addis-Abeba — l'unit\'e devient une organisation (compromis Casablanca/Monrovia).
\item \textbf{2002} : transformation de l'OUA en Union africaine — vers une int\'egration plus pouss\'ee.
\item \textbf{2021} : entr\'ee en fonctionnement de la ZLECAf — l'unit\'e devient march\'e commun.
\end{enumerate}
Astuce : 1945 = les id\'ees ; 1963 = l'organisation ; 2002 = l'union ; 2021 = l'\'economie.
\end{methode}

\begin{exoresolu}[Analyser un document historique]
\textbf{\'Enonc\'e.} « Nous disons aux peuples colonis\'es d'Afrique : nous sommes d\'etermin\'es \`a \^etre libres. Nous exigeons l'ind\'ependance maintenant. » (R\'esolution du 5\textsuperscript{e} congr\`es panafricain, Manchester, 1945.) Question : en quoi ce texte marque-t-il un tournant dans l'histoire du panafricanisme ?
\tcblower
\textbf{Solution d\'etaill\'ee.} Trois \'etapes :
\begin{enumerate}
\item \emph{Situer le document} : il s'agit de la r\'esolution finale du congr\`es de Manchester (octobre 1945), organis\'e par Du Bois avec la participation de jeunes leaders africains comme Nkrumah et Kenyatta, au lendemain de la Seconde Guerre mondiale.
\item \emph{D\'egager la nouveaut\'e} : contrairement au congr\`es de Paris (1919), qui demandait seulement une administration coloniale plus juste, Manchester \emph{exige} l'ind\'ependance imm\'ediate. Le vocabulaire change : « nous exigeons », « maintenant », « d\'etermin\'es ». Le panafricanisme passe du r\'eformisme \`a la revendication politique radicale, et le centre de gravit\'e du mouvement se d\'eplace de la diaspora vers le continent.
\item \emph{\'Evaluer la port\'ee} : ce tournant annonce les ind\'ependances des ann\'ees 1960 et la cr\'eation de l'OUA en 1963. Le panafricanisme cesse d'\^etre un courant d'id\'ees pour devenir un programme d'action, puis une organisation continentale.
\end{enumerate}
\emph{Conclusion} : Manchester 1945 est le moment o\`u le panafricanisme passe de la solidarit\'e culturelle au combat pour la souverainet\'e politique.
\end{exoresolu}

\begin{aretenir}[L'essentiel de la section]
\begin{itemize}
\item Le panafricanisme na\^it dans la \cle{diaspora} (r\'evoltes d'esclaves, Ha\"iti, Garvey et le « Back to Africa », Sylvester Williams, Du Bois).
\item Les congr\`es panafricains structurent le mouvement ; \cle{Manchester 1945} exige l'ind\'ependance imm\'ediate.
\item La \cle{n\'egritude} (Senghor, C\'esaire, Damas) en est la composante culturelle : fiert\'e et revalorisation de la culture noire.
\item Apr\`es les ind\'ependances (1960), le d\'ebat \cle{Casablanca} (unit\'e imm\'ediate, Nkrumah) contre \cle{Monrovia} (coop\'eration progressive) aboutit au compromis de l'\cle{OUA} (1963, Addis-Abeba), devenue \cle{Union africaine} en 2002.
\item Les \'etapes r\'ecentes — \cle{NEPAD} (2001) et \cle{ZLECAf} (op\'erationnelle en 2021) — traduisent le panafricanisme en strat\'egie \'economique de d\'eveloppement.
\end{itemize}
\end{aretenir}

Cet historique montre que le panafricanisme n'est pas une id\'ee abstraite : c'est un mouvement qui a obtenu des r\'esultats concrets, de la d\'ecolonisation au march\'e commun. Reste \`a examiner, dans la section suivante, quels sont pr\'ecis\'ement ses objectifs comme philosophie du d\'eveloppement de l'Afrique.

\coursec{Objectifs du panafricanisme comme philosophie du développement}

Nous avons vu, dans la section précédente, que le panafricanisme est né d'un long cheminement historique : des congrès de la diaspora (Sylvester-Williams, Du Bois) aux indépendances, de la création de l'OUA en 1963 à sa transformation en Union africaine en 2002. Mais une question demeure : \emph{que veut} exactement le panafricanisme ? Un mouvement ne se comprend pleinement que par ses finalités. Or le panafricanisme n'est pas seulement un courant politique : il se présente comme une véritable \cle{philosophie du développement}, c'est-à-dire comme une réponse réfléchie à la question : comment l'Afrique peut-elle sortir du sous-développement ? C'est cette question que nous allons maintenant examiner, en distinguant quatre grands objectifs : politique, économique, culturel et social.

\begin{definition}[Philosophie du développement]
On appelle \cle{philosophie du développement} toute conception réfléchie qui définit ce qu'est le développement authentique d'un peuple et indique les moyens d'y parvenir. Le panafricanisme est une philosophie du développement en ce sens qu'il affirme que le développement de l'Afrique passe par l'\cle{unité}, l'\cle{endogénéité} (le fait de puiser en soi-même ses forces et ses modèles) et la \cle{solidarité}, plutôt que par la dépendance extérieure (aide, tutelle, modèles importés).
\end{definition}

\begin{popfigure}
\begin{center}
\begin{tikzpicture}[scale=1]
\node[circle, draw=popink, fill=popblue!25, thick, minimum size=2.9cm, align=center, font=\small\bfseries] (centre) at (0,0) {Panafricanisme\\ (l'unité comme\\ condition du\\ développement)};
\node[rectangle, draw=popblue, fill=popblueL, rounded corners, thick, align=center, font=\small, minimum width=3.4cm] (pol) at (-4.6,2.4) {\textbf{Objectif politique}\\ unité, souveraineté,\\ une seule voix};
\node[rectangle, draw=poporange, fill=poporangeL, rounded corners, thick, align=center, font=\small, minimum width=3.4cm] (eco) at (4.6,2.4) {\textbf{Objectif économique}\\ intégration, industrie,\\ monnaie commune};
\node[rectangle, draw=popgreen, fill=popgreenL, rounded corners, thick, align=center, font=\small, minimum width=3.4cm] (cul) at (-4.6,-2.4) {\textbf{Objectif culturel}\\ dignité, identité,\\ décolonisation des\\ mentalités};
\node[rectangle, draw=poppurple, fill=poppurpleL, rounded corners, thick, align=center, font=\small, minimum width=3.4cm] (soc) at (4.6,-2.4) {\textbf{Objectif social}\\ éducation, santé,\\ épanouissement humain};
\draw[-{Stealth[length=3mm]}, very thick, popblue] (centre) -- (pol);
\draw[-{Stealth[length=3mm]}, very thick, poporange] (centre) -- (eco);
\draw[-{Stealth[length=3mm]}, very thick, popgreen] (centre) -- (cul);
\draw[-{Stealth[length=3mm]}, very thick, poppurple] (centre) -- (soc);
\end{tikzpicture}
\end{center}
\legende{Schéma en étoile des quatre objectifs du panafricanisme : au centre, le principe unificateur (l'unité comme condition du développement) ; sur les quatre branches, les objectifs politique, économique, culturel et social.}
\end{popfigure}

\courssub{Objectif politique : l'unité et la souveraineté de l'Afrique}

Commençons par une intuition simple. Imaginez un élève seul qui réclame une amélioration des conditions de vie au lycée : sa voix porte peu. Mais si les cinquante élèves de la classe parlent d'une seule voix par l'intermédiaire de leur délégué, on les écoute. C'est exactement le raisonnement politique du panafricanisme, transposé à l'échelle d'un continent.

L'observation historique est frappante : divisée en plus de cinquante États, dont beaucoup comptent moins d'habitants qu'une grande ville d'Asie ou d'Europe, l'Afrique pèse peu dans les négociations internationales (commerce, climat, institutions financières). Chaque État pris isolément est faible face aux grandes puissances et aux firmes multinationales. Le premier objectif du panafricanisme est donc politique : réaliser l'\cle{unité} de l'Afrique afin de garantir sa \cle{souveraineté} réelle.

Il faut distinguer ici indépendance et souveraineté. Les États africains ont obtenu l'\emph{indépendance juridique} dans les années 1960 : ils ont un drapeau, un hymne, un siège à l'ONU. Mais la \emph{souveraineté effective} — la capacité de décider librement de sa politique économique, monétaire et diplomatique — reste limitée par la dépendance financière et par la faiblesse de chaque État isolé. Kwame Nkrumah, premier président du Ghana et théoricien majeur du panafricanisme, résumait cette exigence par une formule restée célèbre : \emph{l'Afrique doit s'unir ou périr}. Pour lui, l'unité politique (il rêvait d'États-Unis d'Afrique) est la condition sans laquelle les indépendances demeurent formelles.

Concrètement, cet objectif se traduit par :
\begin{itemize}
\item la création d'organisations continentales : l'OUA (1963), puis l'\cle{Union africaine} (2002), dotée d'une Commission, d'un Parlement panafricain et d'un Conseil de paix et de sécurité ;
\item la volonté de \emph{parler d'une seule voix} sur la scène mondiale : positions communes dans les négociations commerciales, revendication d'un siège permanent pour l'Afrique au Conseil de sécurité de l'ONU ;
\item la gestion africaine des conflits africains, à travers les forces d'interposition de l'Union africaine, plutôt que l'attente d'interventions extérieures.
\end{itemize}

\courssub{Objectif économique : l'intégration et la transformation locale}

L'intuition économique est tout aussi accessible. Un cultivateur de cacao de la région du Centre qui vend seul sa récolte à un acheteur étranger n'a aucun pouvoir de négociation : il subit le prix qu'on lui impose. S'il se regroupe avec d'autres producteurs en coopérative, il peut stocker, négocier, voire transformer lui-même son cacao en poudre ou en beurre de cacao, qui se vendent beaucoup plus cher. L'union crée la force économique.

Le deuxième objectif du panafricanisme est donc économique. Il repose sur un constat : l'Afrique possède d'immenses richesses (pétrole, minerais, produits agricoles), mais elle les exporte \emph{bruts}, à bas prix, et réimporte les produits manufacturés à prix élevé. Ce mécanisme, hérité de la colonisation, entretient le sous-développement. Le panafricanisme économique vise à rompre ce cercle par quatre moyens :
\begin{enumerate}
\item l'\cle{intégration des marchés} : abaisser les barrières douanières entre pays africains pour que les Africains commercent d'abord entre eux (le commerce intra-africain ne représente qu'environ 15 \% du commerce total du continent, contre plus de 60 \% en Europe) ;
\item l'\cle{industrialisation} : construire des usines africaines plutôt que d'importer les produits finis ;
\item la \cle{transformation locale des matières premières} : transformer le cacao en chocolat, le coton en tissu, le pétrole en produits raffinés, afin que la valeur ajoutée profite aux Africains ;
\item la perspective d'une \cle{monnaie commune}, qui affranchirait le continent de la dépendance monétaire (comme celle du franc CFA, arrimé à l'euro et garanti par le Trésor français).
\end{enumerate}

\begin{exemplebox}[La ZLECAf : le plus grand marché commun du monde]
Entrée en vigueur en 2021, la \cle{Zone de libre-échange continentale africaine} (ZLECAf) réunit 54 pays et environ 1,4 milliard de personnes : c'est la plus grande zone de libre-échange du monde en nombre d'États participants. Son PIB combiné est estimé à environ 3 400 milliards de dollars. Son ambition : supprimer progressivement les droits de douane sur 90 \% des marchandises échangées entre pays africains, stimuler l'industrie continentale et créer des millions d'emplois.
\end{exemplebox}

\begin{popfigure}
\begin{center}
\begin{tikzpicture}
\begin{axis}[
ybar, bar width=1.1cm,
width=0.85\linewidth, height=6.2cm,
ymin=0, ymax=30000,
ylabel={PIB combiné (milliards de dollars)},
symbolic x coords={ZLECAf, UE, US},
xtick=data,
xticklabels={ZLECAf (54 pays), Union européenne, États-Unis},
x tick label style={font=\small, align=center},
nodes near coords,
every node near coord/.append style={font=\small},
ymajorgrids=true, grid style={popline!40},
]
\addplot[fill=poporange, draw=popdark] coordinates {(ZLECAf,3400) (UE,17000) (US,27000)};
\end{axis}
\end{tikzpicture}
\end{center}
\legende{Potentiel économique de la ZLECAf : un marché de 1,4 milliard de consommateurs et un PIB combiné d'environ 3 400 milliards de dollars (ordres de grandeur, années 2020). Le graphique montre que l'Afrique unie devient un acteur économique significatif, alors que chacun de ses États isolé pèse très peu.}
\end{popfigure}

\courssub{Objectif culturel : la restauration de la dignité et de l'identité africaines}

Le développement ne se réduit pas à l'argent. Un peuple qui méprise sa propre culture, ses langues et son histoire ne peut se développer de manière authentique : il imite au lieu de créer. C'est pourquoi le panafricanisme assigne un troisième objectif, culturel : restaurer la \cle{dignité} et l'\cle{identité} africaines, que la colonisation et l'esclavage ont profondément blessées.

La colonisation n'a pas seulement exploité les richesses : elle a imposé l'idée que les cultures africaines étaient inférieures, que le progrès consistait à ressembler à l'Occident. Cette intériorisation du mépris de soi est ce que les penseurs africains appellent l'\emph{aliénation culturelle}. D'où l'exigence d'une \cle{décolonisation des mentalités}. L'écrivain kényan \cle{Ngũgĩ wa Thiong'o} en donne l'exemple le plus frappant : ayant d'abord écrit en anglais, il décida d'écrire ses œuvres dans sa langue maternelle, le gikuyu, affirmant dans \emph{Décoloniser l'esprit} (1986) que la langue est le porteur de la culture et qu'un peuple qui abandonne ses langues abandonne une part de son âme. Dans la même lignée, la négritude de Senghor et de Césaire avait déjà affirmé la valeur de la civilisation négro-africaine, et Cheikh Anta Diop avait travaillé à restituer à l'Afrique son ancienneté historique (l'Égypte antique comme civilisation négro-africaine).

Concrètement, cet objectif culturel suppose : la valorisation des langues africaines dans l'enseignement, la réécriture de l'histoire de l'Afrique par les Africains eux-mêmes, la promotion des arts, des musiques et des savoirs traditionnels, et la confiance en soi collective.

\begin{savaistu}[La tontine, un panafricanisme de tous les jours]
Au Cameroun, la \cle{tontine} — association d'épargne et de crédit rotatif où chaque membre cotise et reçoit à tour de rôle la somme réunie — est une forme traditionnelle de solidarité économique. Sans banque, sans contrat écrit, elle repose sur la confiance et l'entraide. Elle illustre localement l'esprit panafricain : unir ses forces pour réaliser ensemble ce qu'aucun ne pourrait accomplir seul. Les économistes parlent aujourd'hui de « microfinance » pour désigner ce que les villages africains pratiquent depuis des siècles.
\end{savaistu}

\courssub{Objectif social : le développement humain et l'épanouissement de la personne}

Enfin, à quoi serviraient l'unité politique et la richesse économique si elles ne se traduisaient pas dans la vie concrète des hommes ? Le quatrième objectif du panafricanisme est social : le \cle{développement humain}, c'est-à-dire l'amélioration effective des conditions de vie de chaque Africain.

Cet objectif comprend :
\begin{itemize}
\item l'\cle{éducation} : scolarisation de masse, formation technique et professionnelle, universités africaines capables de former les cadres dont le continent a besoin, au lieu de subir la fuite des cerveaux ;
\item la \cle{santé} : hôpitaux accessibles, lutte contre les grandes endémies (paludisme, VIH/sida), production locale de médicaments — la crise de la Covid-19 a montré le danger de dépendre entièrement des vaccins importés ;
\item l'\cle{épanouissement de la personne africaine} : liberté, dignité, participation citoyenne, égalité hommes-femmes. Le développement véritable n'est pas seulement la croissance du PIB, mais l'élargissement des capacités de chacun à mener une vie qu'il a des raisons de valoriser.
\end{itemize}

On voit ainsi que les quatre objectifs forment un tout cohérent : l'unité politique donne la force, l'intégration économique donne les moyens, la restauration culturelle donne la confiance, et le développement social donne le but final — l'épanouissement de la personne humaine, fin dernière de tout développement authentique.

\courssub{Le panafricanisme comme philosophie : l'unité, condition du développement}

Réunissons maintenant ces objectifs en une thèse philosophique unique. Le panafricanisme affirme que l'\emph{unité n'est pas un luxe mais une condition} : divisée, l'Afrique reste faible, dépendante, méprisée ; unie, elle devient capable de se développer par elle-même. C'est le sens de la formule de Nkrumah : \emph{l'Afrique doit s'unir ou périr}. Le raisonnement peut se formuler ainsi :
\begin{enumerate}
\item le développement exige la souveraineté réelle et la force économique ;
\item or aucun État africain isolé ne possède cette force ;
\item donc le développement de l'Afrique passe par l'unité des Africains.
\end{enumerate}

Cette philosophie repose sur deux principes : l'\cle{endogénéité} (le développement doit partir des réalités, des besoins et des ressources propres de l'Afrique, non de modèles importés) et la \cle{solidarité} (les Africains partagent un destin commun : la faiblesse de l'un affaiblit tous). Le panafricanisme prolonge ainsi, à l'échelle du continent, les valeurs communautaires traditionnelles africaines de l'entraide — ce que la philosophie sud-africaine de l'\emph{Ubuntu} exprime par la formule : « je suis parce que nous sommes ».

\begin{attention}[L'unité n'est pas l'uniformité]
Piège fréquent dans les copies : présenter l'unité africaine comme la suppression des différences. C'est une erreur. Le panafricanisme ne vise pas à effacer la diversité des langues, des cultures et des traditions africaines, qui est une richesse. Il vise une \emph{unité dans la diversité} : des peuples différents qui coopèrent autour d'intérêts et d'un destin communs, comme des instruments différents jouent ensemble dans un même orchestre. Ne confondez jamais unité (coopération volontaire) et uniformité (identité imposée).
\end{attention}

\courssub{Application à la vie courante : l'esprit d'unité autour de nous}

La notion d'unité panafricaine n'est pas une abstraction réservée aux sommets des chefs d'État : elle s'incarne dans des pratiques concrètes que l'élève camerounais connaît bien.

\begin{methode}[Appliquer la notion d'unité à une situation concrète]
Pour analyser une situation de la vie courante à la lumière du panafricanisme, procédez en trois étapes :
\begin{enumerate}
\item identifiez les \textbf{acteurs séparés} : qui agit isolément et avec quelles limites ?
\item dégagez ce que \textbf{l'union leur permettrait} : quelle force, quel gain, quelle autonomie nouvelle ?
\item repérez les \textbf{obstacles à lever} : méfiance, égoïsme, frontières, rivalités, intérêts particuliers.
\end{enumerate}
\end{methode}

\begin{exoresolu}[Les tontines et les marchés transfrontaliers : un panafricanisme vécu]
\textbf{Énoncé.} À l'aide de la méthode en trois étapes, montrez que la tontine villageoise et le marché transfrontalier entre le Cameroun et le Nigeria illustrent concrètement les objectifs du panafricanisme.
\tcblower
\textbf{Solution détaillée.}

\emph{1. Les acteurs séparés.} Dans le village, chaque habitant dispose d'une épargne trop faible pour financer seul un projet (acheter une moto, payer une dot, lancer un commerce). De même, les commerçants camerounais et nigérians des zones frontalières (Kumba, Mamfé, Ikom côté nigérian) vendent chacun sur de petits marchés locaux, limités et pauvres.

\emph{2. Ce que l'union permet.} Grâce à la tontine, les cotisations réunies offrent à chaque membre, à tour de rôle, un capital important : l'union des petites épargnes crée la capacité d'investir — c'est l'objectif économique panafricain (mise en commun des forces) à l'échelle du village. Grâce aux échanges transfrontaliers, le commerçant camerounais écoule son maïs, son huile de palme ou son bétail sur le vaste marché nigérian, et reçoit en retour des produits manufacturés : c'est l'intégration des marchés prônée par la ZLECAf, déjà pratiquée spontanément par les populations.

\emph{3. Les obstacles à lever.} Pour la tontine : la méfiance et le risque de défaillance d'un membre, que la confiance et les règles du groupe permettent de surmonter. Pour les échanges transfrontaliers : les barrières douanières, les tracasseries aux frontières, la méfiance entre États — précisément les obstacles que la CEMAC (Communauté économique et monétaire de l'Afrique centrale) et la ZLECAf cherchent à supprimer.

\emph{Conclusion.} Ces pratiques montrent que le panafricanisme n'est pas d'abord une idéologie venue d'en haut : il prolonge et généralise, à l'échelle du continent, un esprit d'unité et de solidarité que les Africains vivent déjà au quotidien. L'unité apparaît bien comme la condition du développement, du village jusqu'au continent.
\end{exoresolu}

\begin{aretenir}[L'essentiel de la section]
\begin{itemize}
\item Le panafricanisme est une \cle{philosophie du développement} : il affirme que l'Afrique se développera par l'unité, l'endogénéité et la solidarité, non par la dépendance extérieure.
\item Quatre objectifs : \textbf{politique} (unité, souveraineté, une seule voix), \textbf{économique} (intégration des marchés, industrialisation, transformation locale, monnaie commune), \textbf{culturel} (dignité, identité, décolonisation des mentalités — Ngũgĩ wa Thiong'o), \textbf{social} (éducation, santé, épanouissement humain).
\item Formule-clé de Nkrumah : \emph{l'Afrique doit s'unir ou périr}.
\item La ZLECAf (54 pays, 1,4 milliard de personnes) incarne aujourd'hui l'objectif économique.
\item L'unité recherchée respecte la diversité : unité n'est pas uniformité.
\item Tontines, coopératives, CEMAC et marchés transfrontaliers sont des formes concrètes de l'esprit panafricain.
\end{itemize}
\end{aretenir}

\begin{exercice}[Pour s'entraîner]
\begin{enumerate}
\item Définissez la « philosophie du développement » et montrez en quoi le panafricanisme en est une.
\item Expliquez la formule de Nkrumah : « l'Afrique doit s'unir ou périr ». Quels objectifs du panafricanisme justifient cette exigence ?
\item Sujet de dissertation : « Le développement de l'Afrique passe-t-il nécessairement par l'unité des Africains ? » Vous construirez un plan en deux ou trois parties en mobilisant les quatre objectifs étudiés.
\end{enumerate}
\end{exercice}

\begin{corrige}[Exercice n° 2]
La formule de Nkrumah signifie que l'unité est pour l'Afrique une question de survie, non un simple idéal. Elle se justifie par les quatre objectifs : \emph{politiquement}, des États isolés ne peuvent défendre leur souveraineté face aux grandes puissances ; \emph{économiquement}, des marchés trop étroits ne permettent ni industrialisation ni transformation locale des matières premières, ce qui perpétue la dépendance ; \emph{culturellement}, seule une Afrique unie et confiante peut restaurer sa dignité bafouée par la colonisation ; \emph{socialement}, la mise en commun des ressources est nécessaire pour financer l'éducation et la santé de masse. « Périr » désigne donc le maintien dans le sous-développement et la dépendance ; « s'unir » désigne la condition du développement authentique.
\end{corrige}

\coursec{Difficultés et limites du panafricanisme}

Dans la section précédente, nous avons montré que le panafricanisme se donne des objectifs ambitieux : l'unité politique du continent, l'indépendance économique réelle, la revalorisation des cultures africaines et la solidarité de tous les peuples d'ascendance africaine. Mais entre le projet et sa réalisation s'interpose l'épreuve du réel. Plus d'un demi-siècle après les indépendances et la création de l'Organisation de l'unité africaine (OUA) en 1963, l'unité politique de l'Afrique reste un horizon plutôt qu'un fait accompli. Il faut donc examiner lucidement les obstacles qui freinent le panafricanisme, puis dresser un bilan nuancé : ni triomphalisme naïf, ni défaite définitive. C'est précisément l'attitude dialectique qu'exige la dissertation philosophique.

\courssub{Les héritages coloniaux : des difficultés structurelles}

\begin{definition}[Héritage colonial]
On appelle \cle{héritage colonial} l'ensemble des structures (frontières, langues, institutions, circuits économiques) léguées par la colonisation et qui continuent, après les indépendances, d'organiser la vie des États africains. Ces héritages constituent des difficultés \emph{structurelles} : elles ne dépendent pas directement de la volonté des dirigeants actuels.
\end{definition}

L'intuition est simple : on ne bâtit pas une maison commune sur des fondations posées par d'autres pour d'autres fins. L'Afrique unie rêvée par Nkrumah ou Nyerere devait s'édifier sur un continent découpé, administré et exploité par l'Europe. Trois héritages pèsent particulièrement lourd.

\textbf{1. Les frontières artificielles de Berlin.} La conférence de Berlin (novembre 1884 -- février 1885) a partagé l'Afrique entre les puissances européennes selon des tracés géométriques (lignes droites, méridiens, parallèles) ignorant les réalités humaines : un même peuple se retrouva séparé entre deux ou trois colonies, tandis que des groupes historiquement rivaux furent enfermés dans le même État. À l'indépendance, l'OUA a consacré le principe de l'\emph{intangibilité des frontières} (résolution du Caire, 1964) pour éviter des guerres sans fin : les États africains héritent donc de frontières qu'ils n'ont ni choisies ni le droit de remettre en cause.

\begin{savaistu}[Berlin 1884-1885 : l'Afrique partagée sans les Africains]
La conférence de Berlin, organisée par le chancelier allemand Bismarck, a réuni quatorze puissances européennes (plus les États-Unis en observateurs) pour régler le partage de l'Afrique. \cle{Aucun Africain} n'était présent ni même consulté. Les tracés décidés à la règle sur des cartes imprécises divisent encore aujourd'hui des peuples entiers : les Ewondo-Bulu entre le Cameroun, le Gabon et la Guinée équatoriale, les peuples du bassin du Logone entre le Cameroun et le Tchad, les Maasaï entre le Kenya et la Tanzanie. Beaucoup de conflits et de tensions identitaires actuels plongent leurs racines dans ce partage.
\end{savaistu}

\textbf{2. Les langues officielles imposées.} Le français, l'anglais, le portugais ou l'espagnol sont restés les langues de l'administration, du droit et de l'école. Conséquence : un Camerounais francophone et un Nigérian anglophone, voisins de frontière, communiquent moins facilement entre eux que chacun avec Paris ou Londres. Le continent se trouve fragmenté en blocs linguistiques (francophonie, Commonwealth, lusophonie) qui reproduisent les clivages coloniaux et compliquent toute institution commune.

\textbf{3. Les liens économiques verticaux.} La colonisation a organisé les économies africaines comme des appendices des métropoles : chemins de fer allant de la mine vers le port (jamais d'un pays africain vers son voisin), cultures d'exportation (cacao, café, coton) destinées à l'Europe, monnaies adossées aux monnaies coloniales. Après l'indépendance, ces circuits \emph{verticaux} (colonie-métropole) ont survécu au détriment des échanges \emph{horizontaux} (Afrique-Afrique). Douala communique économiquement plus avec Marseille qu'avec Lagos, pourtant toute proche.

\begin{popfigure}
\begin{center}
\begin{tikzpicture}[scale=0.9]
% Contour très simplifié de l'Afrique
\draw[thick, popdark, fill=popcream] (0,4.5) -- (1.2,5.6) -- (2.6,6.1) -- (4.2,5.9) -- (5.4,6.4) -- (6.6,6.0) -- (7.4,5.2) -- (7.8,4.2) -- (8.4,3.2) -- (8.0,2.2) -- (7.2,1.4) -- (6.4,0.4) -- (5.4,-0.6) -- (4.6,-1.2) -- (4.0,-0.6) -- (3.4,0.4) -- (2.6,1.2) -- (1.8,1.8) -- (1.0,2.2) -- (0.2,2.8) -- (-0.4,3.6) -- cycle;
% Corne de l'Afrique
\draw[thick, popdark, fill=popcream] (7.8,4.2) -- (9.4,4.8) -- (10.0,4.2) -- (8.6,3.4) -- (8.4,3.2);
% Frontières artificielles (lignes droites géométriques)
\draw[popmuted, dashed] (2.0,6.05) -- (2.2,1.6);
\draw[popmuted, dashed] (4.0,5.95) -- (3.9,0.2);
\draw[popmuted, dashed] (5.8,6.15) -- (5.6,0.0);
\draw[popmuted, dashed] (0.6,4.4) -- (8.0,4.3);
\draw[popmuted, dashed] (1.2,3.2) -- (8.3,3.0);
\draw[popmuted, dashed] (7.0,5.5) -- (7.3,1.5);
% Zones de conflit (hachurées)
\fill[poppink, opacity=0.55] (1.6,3.4) circle (0.55); % Sahel ouest
\fill[poppink, opacity=0.55] (5.2,3.6) circle (0.6); % bassin lac Tchad
\fill[poppink, opacity=0.55] (6.2,2.6) circle (0.5); % RDC est
\fill[poppink, opacity=0.55] (8.6,4.1) circle (0.45); % Corne
% Points villes
\fill[popblue] (1.9,2.6) circle (0.07) node[below left, font=\scriptsize, popdark]{Douala};
\fill[popblue] (5.9,5.6) circle (0.07) node[above, font=\scriptsize, popdark]{Le Caire};
\fill[popblue] (4.7,-0.5) circle (0.07) node[below, font=\scriptsize, popdark]{Le Cap};
\fill[popblue] (0.6,4.0) circle (0.07) node[left, font=\scriptsize, popdark]{Dakar};
% Nord
\draw[-{Stealth}, thick, popdark] (9.6,5.6) -- (9.6,6.4) node[above, font=\scriptsize]{N};
\end{tikzpicture}
\end{center}
\legende{Carte schématique de l'Afrique. En pointillés : frontières héritées de la colonisation, souvent tracées à la règle sans tenir compte des peuples. En rose : zones marquées par des conflits ou l'insécurité (Sahel, bassin du lac Tchad avec Boko Haram, est de la RDC, Corne de l'Afrique). Croquis indicatif, sans précision cartographique.}
\end{popfigure}

\courssub{Les égoïsmes nationaux : la souveraineté contre l'unité}

\begin{definition}[Égoïsme national]
On appelle \cle{égoïsme national} l'attitude d'un État qui fait passer ses intérêts propres (souveraineté, prestige, avantages économiques immédiats) avant l'intérêt commun de l'ensemble continental. En philosophie politique, c'est l'application à la scène internationale du problème classique de l'individu et de la cité : chacun veut jouir de sa liberté sans accepter la loi commune.
\end{definition}

Dès 1963, le débat entre le « groupe de Casablanca » (Nkrumah, partisan d'un gouvernement continental immédiat) et le « groupe de Monrovia » (partisans d'une simple coopération entre États souverains) s'est soldé par la victoire des seconds : l'OUA fut fondée sur le principe de \emph{non-ingérence} et de respect de la souveraineté de chaque État. Autrement dit, l'organisation de l'unité africaine a été construite sur la négation de l'unité politique réelle : aucun organe supranational ne peut imposer une décision à un État membre.

À cela s'ajoutent les \cle{rivalités entre leaders} : compétition pour le leadership continental (entre l'Afrique du Sud, le Nigeria, l'Égypte ou l'Éthiopie), susceptibilités personnelles, crainte des chefs d'État de perdre leur pouvoir au profit d'une autorité commune. Les nationalismes, légitimes comme fierté identitaire, deviennent un obstacle lorsqu'ils se ferment : xénophobie contre les migrants africains, fermetures de frontières, expulsions de ressortissants de pays voisins. L'unité africaine se heurte ainsi au paradoxe suivant : les États nés de la lutte anticoloniale ont fait de la souveraineté leur bien le plus précieux, et c'est précisément cette souveraineté qu'il faudrait partiellement abandonner pour réaliser l'unité.

\courssub{Les faiblesses économiques : un continent qui ne commerçait pas avec lui-même}

L'observation statistique est ici décisive : la part des échanges \emph{intra-africains} dans le commerce total de l'Afrique n'est que d'environ \num{15}~\% , contre environ \num{60}~\% en Europe et \num{55}~\% en Asie. Autrement dit, les pays africains vendent et achètent essentiellement hors du continent : matières premières vers l'Europe, la Chine ou l'Amérique ; produits manufacturés importés de ces mêmes régions.

\begin{exemplebox}[L'exemple chiffré des échanges intra-africains]
Sur 100 francs de commerce réalisé par les pays africains, seulement environ 15 francs correspondent à des échanges \emph{entre pays africains}. En Europe, ce chiffre atteint environ 60 francs. Conséquence concrète : il est souvent moins coûteux d'importer du riz d'Asie au Cameroun que du riz du Nigeria voisin, faute de routes, de voies ferrées et d'accords commerciaux efficaces. Cette faiblesse prive le continent de la force de négociation qu'un grand marché unifié lui donnerait face aux puissances extérieures.
\end{exemplebox}

\begin{popfigure}
\begin{center}
\begin{tikzpicture}
\begin{axis}[
    ybar, bar width=1.1cm,
    width=0.85\linewidth, height=6.5cm,
    ymin=0, ymax=70,
    ylabel={Part des échanges intra-régionaux (en \%)},
    symbolic x coords={Afrique, Asie, Europe},
    xtick=data,
    nodes near coords, nodes near coords align={vertical},
    every node near coord/.append style={font=\small\bfseries},
    ymajorgrids=true, grid style={popline, dashed},
    axis lines=left,
    enlarge x limits=0.3
]
\addplot[fill=poporange, draw=popdark] coordinates {(Afrique,15) (Asie,55) (Europe,60)};
\end{axis}
\end{tikzpicture}
\end{center}
\legende{Part approximative des échanges commerciaux réalisés \emph{à l'intérieur} de chaque région du monde (ordres de grandeur, années 2010-2020). L'Afrique (\(\approx 15\,\%\)) commerce trois à quatre fois moins avec elle-même que l'Europe (\(\approx 60\,\%\)) ou l'Asie (\(\approx 55\,\%\)).}
\end{popfigure}

Les causes de cette faiblesse sont connues : \cle{infrastructures de liaison insuffisantes} (routes et chemins de fer conçus pour évacuer les matières premières vers les ports, non pour relier les pays entre eux), similarité des productions (les économies africaines exportent toutes des matières premières et sont donc \emph{concurrentes} plutôt que complémentaires), droits de douane et lourdeurs administratives aux frontières. La Zone de libre-échange continentale africaine (ZLECAf), entrée en vigueur en 2021, vise précisément à corriger cette faiblesse : nous y reviendrons dans le bilan.

\courssub{Les dépendances extérieures : l'unité empêchée de l'extérieur}

\begin{definition}[Dépendance]
La \cle{dépendance} désigne une situation dans laquelle un État, bien que politiquement indépendant, voit ses décisions essentielles (économiques, monétaires, militaires) déterminées ou fortement contraintes par des puissances ou des institutions extérieures. Les théoriciens de la dépendance (A. G. Frank, S. Amin) parlent de \emph{néocolonialisme} : la domination change de forme, elle ne disparaît pas.
\end{definition}

Quatre mécanismes entretiennent cette dépendance et divisent le continent :

\begin{enumerate}
\item \textbf{L'aide publique au développement.} Présentée comme une générosité, elle crée souvent une relation de tutelle : le donateur oriente les politiques du pays aidé, et l'aide peut être suspendue comme sanction politique.
\item \textbf{La dette extérieure.} Le service de la dette (remboursements et intérêts) absorbe une part considérable des budgets : des milliards qui pourraient financer écoles, hôpitaux et routes partent chaque année vers les créanciers étrangers. Les plans d'ajustement structurel imposés par le FMI et la Banque mondiale dans les années 1980-1990 ont imposé privatisations et réductions des dépenses sociales, affaiblissant les États.
\item \textbf{Le franc CFA.} Utilisé par quatorze pays d'Afrique centrale et de l'Ouest, dont le Cameroun, cette monnaie héritée de la colonisation reste arrimée à l'euro et garantie par le Trésor français. Ses partisans soulignent la stabilité monétaire ; ses détracteurs (comme l'économiste sénégalais Ndongo Samba Sylla) y voient une servitude monétaire qui prive les États d'un instrument essentiel de souveraineté et divise le continent en zones monétaires hétérogènes.
\item \textbf{Les ingérences étrangères.} Bases militaires, soutien à des régimes dociles ou à des rebellions, accaparement des terres et des matières premières stratégiques (pétrole, cobalt, uranium, coltan) : les puissances extérieures entretiennent parfois délibérément la division africaine, car une Afrique unie négocierait ses ressources avec un rapport de force inversé.
\end{enumerate}

\courssub{Conflits et instabilité : quand la survie prime sur le développement}

Un pays en guerre ne construit pas : il détruit. Les conflits détournent vers la défense et la sécurité des ressources qui devraient aller à l'éducation, à la santé et aux infrastructures. Ils détruisent aussi le capital humain (morts, déplacés, réfugiés) et découragent tout investissement.

L'Afrique a connu depuis les indépendances de nombreuses guerres civiles (Nigeria-Biafra, Rwanda, RDC, Soudan), des coups d'État récurrents et, plus récemment, le terrorisme. Le Cameroun en offre un exemple direct : depuis 2014, la secte \cle{Boko Haram} mène des attaques dans l'Extrême-Nord (massacres, enlèvements, attentats-suicides), obligeant l'État à mobiliser d'importantes forces armées et provoquant le déplacement de centaines de milliers de personnes. À l'ouest, la crise anglophone illustre comment l'héritage colonial (la double tutelle française et britannique) peut se transformer en conflit interne durable. Chaque franc dépensé pour la guerre est un franc soustrait au développement ; chaque État déstabilisé est un maillon faible de l'édifice panafricain.

\begin{methode}[Distinguer difficulté structurelle et difficulté conjoncturelle]
Pour analyser rigoureusement une difficulté du panafricanisme dans une dissertation, procédez en trois étapes :
\begin{enumerate}
\item \textbf{Identifier} la difficulté et la formuler clairement (exemple : faiblesse des échanges intra-africains).
\item \textbf{Qualifier} sa nature : est-elle \emph{structurelle} (héritée de l'histoire, inscrite dans les structures : frontières de Berlin, langues, circuits économiques coloniaux) ou \emph{conjoncturelle} (liée à des choix politiques actuels, donc réversible : rivalités de leaders, corruption, mauvaise gouvernance) ?
\item \textbf{Conclure} sur la marge d'action : une difficulté structurelle exige un travail de longue durée (réformes monétaires, infrastructures continentales) ; une difficulté conjoncturelle relève de la volonté politique immédiate. Cette distinction montre que les obstacles ne sont pas une fatalité.
\end{enumerate}
\end{methode}

\courssub{Bilan nuancé : ni échec total, ni succès accompli}

\begin{attention}[Erreur fréquente : conclure que le panafricanisme a « échoué »]
Le piège classique de la copie consiste à accumuler les difficultés précédentes et à conclure : « le panafricanisme est un échec ». C'est une erreur de méthode et de pensée. Erreur de méthode, car le devoir de philosophie exige un \emph{bilan dialectique} : peser les limites ET les avancées, puis trancher par un jugement argumenté. Erreur de pensée, car les faits contredisent le verdict d'échec : l'Union africaine existe, la ZLECAf fonctionne depuis 2021, les organisations régionales (CEMAC, CEDEAO, SADC) intègrent progressivement leurs économies, et l'idéal panafricain continue d'inspirer les jeunesses du continent. Une difficulté n'est pas une impossibilité.
\end{attention}

Dressons donc le bilan. Du côté des \cle{limites} : l'unité politique immédiate rêvée par Nkrumah (États-Unis d'Afrique, armée continentale, gouvernement commun) n'a pas eu lieu ; l'UA, comme l'OUA avant elle, reste une organisation interétatique sans pouvoir de contrainte ; les dépendances extérieures et les conflits persistent.

Du côté des \cle{avancées réelles} :
\begin{itemize}
\item la transformation de l'OUA en \textbf{Union africaine} (2002), dotée d'un Parlement panafricain, d'une Cour africaine des droits de l'homme et d'un Conseil de paix et de sécurité, et qui a abandonné le principe absolu de non-ingérence en cas de crimes graves ;
\item la \textbf{ZLECAf} (Zone de libre-échange continentale africaine), opérationnelle depuis le 1\textsuperscript{er} janvier 2021, qui constitue le plus grand marché commun du monde par le nombre de pays et vise à porter les échanges intra-africains bien au-delà des \num{15}~\% actuels ;
\item les \textbf{organisations régionales} (CEMAC et CEEAC en Afrique centrale, CEDEAO à l'ouest, SADC au sud) qui réalisent concrètement, pas à pas, l'intégration : monnaies communes, libre circulation partielle, forces d'interposition ;
\item la victoire historique déjà remportée : c'est la solidarité panafricaine qui a permis d'achever la décolonisation politique du continent et d'isoler l'apartheid sud-africain jusqu'à sa chute en 1994.
\end{itemize}

\begin{exoresolu}[Sujet de dissertation : « Le panafricanisme a-t-il échoué ? »]
Énoncé. Un élève doit traiter le sujet suivant : « Le panafricanisme a-t-il échoué ? » Construisez le plan dialectique de la réponse.
\tcblower
\textbf{Introduction.} Accroche : plus de soixante ans après l'OUA, l'Afrique reste fragmentée et dépendante. Définition du panafricanisme (mouvement visant l'unité et la libération des peuples africains). Problématique : l'écart entre les objectifs et les résultats signifie-t-il l'échec du projet, ou seulement sa lenteur ?
\par\textbf{Thèse (I).} On peut soutenir l'échec : frontières coloniales intangibles, égoïsmes nationaux et rivalités de leaders, échanges intra-africains de seulement \(\approx 15\,\%\), dépendances (dette, franc CFA, ingérences), conflits (Boko Haram, guerres civiles). L'unité politique n'a pas eu lieu.
\par\textbf{Antithèse (II).} Mais l'échec n'est pas total : l'UA dispose d'institutions communes, la ZLECAf crée un marché continental, les organisations régionales avancent, et le panafricanisme a déjà gagné sa première bataille (décolonisation, fin de l'apartheid). Beaucoup de difficultés sont structurelles et demandent du temps, non l'abandon du projet.
\par\textbf{Synthèse (III).} Le panafricanisme n'a pas échoué : il est \emph{inachevé}. Comme tout grand projet politique (l'unité européenne a mis un demi-siècle), il progresse par étapes. Son échec ne serait réel que si les Africains y renonçaient ; or la ZLECAf prouve le contraire. Le panafricanisme est moins un bilan qu'une tâche.
\end{exoresolu}

\begin{aretenir}[L'essentiel de la section]
\begin{itemize}
\item Les difficultés du panafricanisme sont d'abord \textbf{structurelles} : frontières artificielles de Berlin (1885), langues coloniales, circuits économiques verticaux hérités de la colonisation.
\item Elles sont aussi \textbf{politiques} : égoïsmes nationaux, attachement à la souveraineté, rivalités entre leaders.
\item Elles sont \textbf{économiques} : échanges intra-africains d'environ \num{15}~\% seulement, infrastructures de liaison insuffisantes.
\item Elles sont \textbf{extérieures} : aide, dette, franc CFA, ingérences des puissances.
\item Les \textbf{conflits} (guerres civiles, coups d'État, Boko Haram au Cameroun) détournent les ressources du développement.
\item Le bilan est \textbf{dialectique} : échec de l'unité politique immédiate, mais avancées réelles (UA, ZLECAf, intégrations régionales). Ne jamais conclure à l'échec pur et simple.
\end{itemize}
\end{aretenir}

\begin{exercice}[Pour s'entraîner]
1. Définissez les termes : héritage colonial, égoïsme national, dépendance.\par
2. Pour chacune des difficultés suivantes, dites si elle est structurelle ou conjoncturelle, en justifiant : (a) les frontières de Berlin ; (b) la rivalité entre deux chefs d'État ; (c) le franc CFA ; (d) un coup d'État.\par
3. Expliquez en dix lignes pourquoi la faiblesse des échanges intra-africains (\(\approx 15\,\%\)) constitue un obstacle au développement du continent, et ce que la ZLECAf peut y changer.
\end{exercice}

\begin{corrige}[Exercice]
\textbf{1.} Héritage colonial : ensemble des structures (frontières, langues, circuits économiques) léguées par la colonisation et qui persistent après l'indépendance. Égoïsme national : attitude d'un État qui privilégie ses intérêts propres au détriment de l'intérêt commun continental. Dépendance : situation où un État formellement souverain voit ses décisions essentielles contraintes par des puissances extérieures.
\par\textbf{2.} (a) Structurelle : héritée de la colonisation, figée par le principe d'intangibilité. (b) Conjoncturelle : elle relève des personnes au pouvoir et peut disparaître avec un changement de volonté politique. (c) Structurelle : institution monétaire héritée, dont la réforme exige une négociation de long terme (même si la décision d'en sortir serait un acte politique). (d) Conjoncturelle : événement politique contingent, même s'il s'enracine parfois dans des fragilités structurelles.
\par\textbf{3.} Commercer peu entre soi prive les pays africains de débouchés proches, maintient la dépendance envers les marchés extérieurs, empêche l'industrialisation par le marché régional et affaiblit le rapport de force dans les négociations mondiales. La ZLECAf, en supprimant progressivement les droits de douane entre pays africains, crée un marché de plus d'un milliard de consommateurs : elle peut stimuler la production locale, attirer les investissements et faire passer la part des échanges intra-africains bien au-delà de \num{15}~\%, à condition d'être accompagnée d'infrastructures de liaison.
\end{corrige}

\coursec{Synthèse : le panafricanisme, utopie ou voie d'avenir ?}

Après avoir examiné les difficultés et les limites du panafricanisme — divisions linguistiques et politiques, rivalités entre États, dépendances économiques persistantes, faiblesse des institutions continentales — une question s'impose avec force : faut-il en conclure que le projet d'unité africaine est une \cle{utopie}, c'est-à-dire un rêve irréalisable ? Ou bien le panafricanisme demeure-t-il, malgré ses échecs, la \cle{voie d'avenir} la plus cohérente pour le développement du continent ? Cette dernière section propose de confronter les deux thèses, puis de les dépasser dans une position dialectique qui prépare directement à la conclusion de la dissertation.

\courssub{Rappel des notions clés : développement, sous-développement et unité}

\begin{definition}[Développement]
Le \cle{développement} est le processus global de transformation économique, sociale, politique et culturelle d'une société, qui améliore durablement les conditions de vie de l'ensemble de sa population. Il ne se réduit pas à la croissance économique (augmentation du PIB) : il implique l'accès à l'éducation, à la santé, à la justice et à la participation politique. Le \cle{sous-développement} désigne la situation d'une société qui n'a pas atteint ces seuils, non par fatalité naturelle, mais le plus souvent à la suite de causes historiques (colonisation, traite) et structurelles (dépendance, mauvaise gouvernance, éclatement des marchés).
\end{definition}

Le lien avec l'unité africaine est direct : l'Afrique compte 55 États, dont beaucoup sont de petits marchés intérieurs, dépendants des exportations de matières premières brutes. Le sous-développement du continent n'est pas seulement un problème national : il est aussi un problème de \textbf{fragmentation}. Un État isolé de quelques millions d'habitants n'a ni la puissance industrielle, ni le poids diplomatique, ni les capacités de négociation commerciale d'un ensemble unifié de plus d'un milliard d'habitants. C'est précisément l'intuition fondatrice du panafricanisme : \emph{l'unité est la condition du développement}. Mais cette intuition est-elle réaliste ? Deux réponses s'affrontent.

\courssub{Première thèse : le panafricanisme est une utopie}

L'\cle{utopie} (du grec \emph{ou-topos}, « nulle part ») désigne un projet idéal qui ne correspond à aucun lieu ni à aucune réalité possible. Appliquée au panafricanisme, cette thèse s'appuie sur des constats solides.

\begin{itemize}
\item \textbf{Des divisions profondes et anciennes.} Le continent est traversé par des fractures linguistiques (anglophones, francophones, lusophones, arabophones), religieuses, ethniques et politiques. Les frontières héritées de la colonisation ont découpé des peuples et enfermé des groupes rivaux dans les mêmes États. Les conflits (Nigéria-Biafra, Soudan, RDC) montrent que l'unité nationale elle-même est déjà difficile : à plus forte raison l'unité continentale.
\item \textbf{Des échecs répétés de l'unité politique.} Le rêve de Nkrumah d'un « gouvernement unioniste africain » a été rejeté dès 1963 au profit du compromis minimal de l'OUA. Les unions régionales ont souvent échoué ou stagné. L'Union africaine, malgré ses ambitions, reste dépendante de financements extérieurs et peine à faire appliquer ses décisions.
\item \textbf{Des dépendances persistantes.} Soixante ans après les indépendances, l'Afrique exporte toujours des matières premières brutes et importe des produits manufacturés ; le commerce intra-africain reste faible (environ 15 à 18 \% des échanges du continent, contre près de 60 \% en Europe). Les élites africaines elles-mêmes sont souvent plus tournées vers Paris, Londres, Washington ou Pékin que vers leurs voisins.
\end{itemize}

Pour les partisans de cette thèse, vouloir l'unité politique de 55 États souverains, jaloux de leurs prérogatives, relève du mythe mobilisateur, non de la politique réaliste.

\courssub{Deuxième thèse : le panafricanisme est une nécessité}

Face à ce scepticisme, la thèse inverse soutient que l'unité n'est pas un luxe idéaliste mais une \textbf{condition de survie} dans le monde contemporain.

\begin{itemize}
\item \textbf{Le poids démographique.} L'Afrique comptera environ 2,5 milliards d'habitants en 2050, soit un quart de l'humanité, avec la population la plus jeune du monde. Dispersée en 55 marchés, cette force est neutralisée ; unie, elle devient un acteur majeur de l'histoire mondiale.
\item \textbf{Le poids économique.} La ZLECAf (Zone de libre-échange continentale africaine, entrée en vigueur en 2019, début des échanges en 2021) crée un marché unique de plus de 1,3 milliard de consommateurs. Aucun État africain seul ne peut industrialiser son économie face à la Chine, aux États-Unis ou à l'Union européenne ; ensemble, les États africains peuvent négocier, transformer localement leurs ressources (cacao, cobalt, pétrole) et capter la valeur ajoutée.
\item \textbf{Le poids diplomatique.} Dans les institutions internationales (ONU, OMC, COP sur le climat), un continent qui parle d'une seule voix obtient ce qu'aucun État isolé ne peut obtenir : siège permanent au Conseil de sécurité, justice climatique, conditions commerciales équitables.
\end{itemize}

Pour les nécessitaristes, l'alternative est brutale : \emph{s'unir ou rester dominé}. Les blocs régionaux (Union européenne, ASEAN, Mercosur) montrent que l'intégration est la loi du monde contemporain ; l'Afrique ne peut pas faire exception.

\begin{popfigure}
\begin{center}
\begin{tabularx}{\linewidth}{|X|X|}
\hline
\rowcolor{popblueL} \textbf{Le panafricanisme, une utopie ?} & \textbf{Le panafricanisme, une nécessité ?} \\
\hline
Divisions linguistiques, ethniques et politiques profondes entre les 55 États. & Démographie : 2,5 milliards d'habitants en 2050, force considérable si elle est unie. \\
\hline
Échecs répétés de l'unité politique (projet unioniste de Nkrumah rejeté en 1963, OUA minimaliste). & Économie : ZLECAf = marché unique de 1,3 milliard de consommateurs ; transformation locale des ressources. \\
\hline
Dépendances persistantes : exportations brutes, faible commerce intra-africain (15-18 \%), institutions financées de l'extérieur. & Diplomatie : une seule voix africaine pèse à l'ONU, à l'OMC et dans les négociations climatiques. \\
\hline
\multicolumn{2}{|c|}{\cellcolor{poporangeL}\textbf{Synthèse :} un idéal régulateur — horizon qui oriente des avancées concrètes et progressives (UA, ZLECAf, intégrations régionales).} \\
\hline
\end{tabularx}
\end{center}
\legende{Confrontation des deux thèses : l'utopie (à gauche) et la nécessité (à droite), dépassées dans une synthèse dialectique.}
\end{popfigure}

\courssub{Position dialectique : le panafricanisme comme idéal régulateur}

La démarche dialectique consiste à ne pas choisir brutalement entre la thèse et l'antithèse, mais à montrer que chacune saisit une part de la vérité, puis à les dépasser dans une \cle{synthèse}.

\begin{definition}[Idéal régulateur]
Un \cle{idéal régulateur} est un but qui n'est peut-être jamais pleinement atteint, mais qui \textbf{oriente et organise l'action} : il sert de boussole, non de destination immédiate. On juge alors un projet non à sa réalisation totale, mais au mouvement réel qu'il engendre.
\end{definition}

Appliquée au panafricanisme, cette catégorie résout l'opposition. Oui, l'unité politique totale du continent est une utopie si on l'imagine comme un État unifié réalisé demain : la thèse sceptique a raison sur ce point. Mais non, le panafricanisme n'est pas pour autant stérile : il fonctionne comme un \textbf{horizon} qui a déjà produit des avancées concrètes — la décolonisation coordonnée, la condamnation unanime de l'apartheid, la création de l'Union africaine, les communautés économiques régionales (CEDEAO, CEEAC), la ZLECAf. La thèse de la nécessité a raison sur ce point. Le panafricanisme est donc une \emph{utopie féconde} : il ne décrit pas un lieu réel, mais il indique une direction, et chaque étape vers lui est un progrès mesurable du développement.

\begin{popfigure}
\begin{center}
\begin{tikzpicture}[node distance=1.1cm and 1.6cm,
  box/.style={draw=popdark, rounded corners=4pt, fill=popblueL, text width=4.6cm, align=center, minimum height=1.6cm, font=\small},
  box2/.style={draw=popdark, rounded corners=4pt, fill=poporangeL, text width=4.6cm, align=center, minimum height=1.6cm, font=\small},
  box3/.style={draw=popdark, rounded corners=4pt, fill=popgreenL, text width=6.2cm, align=center, minimum height=1.8cm, font=\small},
  arr/.style={-{Stealth[length=3mm]}, thick, popdark}]
\node[box, align=center] (these){\textbf{Thèse}\\ Le panafricanisme est une utopie : divisions, échecs, dépendances.};
\node[box2, right=of these, align=center] (antithese){\textbf{Antithèse}\\ Le panafricanisme est une nécessité : poids démographique, économique, diplomatique.};
\node[box3, below=1.4cm of $(these)!0.5!(antithese)$, align=center] (synthese){\textbf{Synthèse}\\ Un idéal régulateur : horizon qui guide des avancées concrètes et progressives (UA, ZLECAf).};
\draw[arr] (these.south) -- (synthese.north west);
\draw[arr] (antithese.south) -- (synthese.north east);
\draw[arr, dashed, popmuted] (these) -- (antithese) node[midway, above, font=\footnotesize\itshape, popmuted]{confrontation};
\end{tikzpicture}
\end{center}
\legende{La démarche dialectique : la confrontation de la thèse et de l'antithèse produit une synthèse qui dépasse l'opposition.}
\end{popfigure}

\courssub{Le rôle de la jeunesse : les nouveaux vecteurs de l'unité}

La synthèse dialectique n'est pas une résignation : elle assigne une tâche, et cette tâche revient d'abord à la jeunesse, majoritaire sur le continent. Trois vecteurs concrets renouvellent le panafricanisme aujourd'hui.

\begin{itemize}
\item \textbf{L'entrepreneuriat.} Les jeunes entrepreneurs africains créent des entreprises qui ignorent les frontières : fintech (paiement mobile panafricain), agritech, start-ups numériques de Lagos à Nairobi, de Dakar à Douala. Chaque entreprise qui vend dans dix pays africains réalise pratiquement l'intégration que les sommets politiques peinent à décréter.
\item \textbf{Le panafricanisme numérique.} Internet et les réseaux sociaux créent un espace public continental : les jeunes de Yaoundé, Abidjan ou Johannesburg débattent des mêmes sujets, partagent les mêmes causes (justice climatique, gouvernance), contournant les vieilles divisions linguistiques par la traduction et la création de contenus communs.
\item \textbf{La culture.} La musique (afrobeat, makossa, coupé-décalé, amapiano), le cinéma (Nollywood, festivals panafricains comme le FESPACO), la littérature et le sport (la Coupe d'Afrique des nations) tissent quotidiennement une identité africaine vécue et joyeuse, plus efficace que bien des traités : on ne s'unit pas d'abord par des institutions, mais par des émotions et des références partagées.
\end{itemize}

\begin{savaistu}[Complément utile au Bac : Kant et les utopies rationnelles]
Dans l'\emph{Idée d'une histoire universelle d'un point de vue cosmopolitique} (1784), le philosophe Emmanuel Kant défend l'idée que l'histoire humaine progresse vers des buts rationnels — paix perpétuelle, société civile universelle — qui ne seront peut-être jamais pleinement réalisés, mais qui doivent \textbf{guider} l'action politique comme des « idées régulatrices ». Citer Kant permet de donner une assise philosophique solide à la synthèse : le panafricanisme, comme la paix perpétuelle kantienne, est une utopie rationnelle qui oriente le progrès réel de l'humanité.
\end{savaistu}

\courssub{Ouverture : de l'unité africaine à la solidarité mondiale}

Si l'unité est la condition du développement de l'Afrique, la même logique vaut à l'échelle de la planète : les grands problèmes contemporains — réchauffement climatique, pandémies, inégalités, migrations — ne connaissent pas de frontières et ne peuvent être résolus par aucun continent isolé. Le panafricanisme ouvre ainsi sur une exigence plus vaste : le \cle{développement} comme exigence de \textbf{justice universelle}. Un monde où un quart de l'humanité demeure en sous-développement est un monde injuste et instable ; la solidarité entre les peuples n'est pas une charité, mais une condition de paix commune. L'unité africaine n'est donc pas un repli identitaire : elle est la contribution de l'Afrique à l'unité du genre humain.

\courssub{Méthode : rédiger la conclusion de la dissertation}

\begin{methode}[La conclusion en trois temps]
\begin{enumerate}
\item \textbf{Répondre à la problématique.} Reprendre la question posée dans l'introduction et y répondre clairement, en une ou deux phrases, en tranchant (ou en dépassant dialectiquement l'alternative). Exemple : « Le panafricanisme n'est donc ni un rêve irréalisable ni un programme déjà accompli, mais un idéal régulateur. »
\item \textbf{Faire le bilan de la réflexion.} Résumer en deux ou trois phrases le chemin parcouru : les deux thèses confrontées et la synthèse qui les dépasse. Ce bilan montre au correcteur que la réponse n'est pas gratuite, mais justifiée par toute la démarche.
\item \textbf{Ouvrir.} Élargir la réflexion vers une question voisine ou plus générale, sans la traiter : « Reste à savoir si cette logique d'unité ne vaut pas, demain, pour l'humanité tout entière. »
\end{enumerate}
\end{methode}

\begin{attention}[Pièges de la conclusion]
Une conclusion ne doit \textbf{jamais introduire un argument nouveau} : tout ce qu'elle affirme doit avoir été démontré dans le développement. Elle ne doit pas non plus être une simple répétition de l'introduction, ni se terminer par une formule vague (« chacun a son opinion », « seul l'avenir nous le dira »). Sa fonction est de \emph{justifier} la réponse et de \emph{clôturer} la démarche : c'est le moment où le correcteur vérifie la cohérence d'ensemble de la copie.
\end{attention}

\begin{exemplebox}[Conclusion type sur le panafricanisme]
« Le panafricanisme est-il une utopie ? Notre réflexion a montré que, pris comme projet d'unité politique immédiate, il se heurte à des divisions réelles et à des échecs répétés ; mais que, considéré comme condition du développement, il demeure une nécessité économique et diplomatique. Le panafricanisme n'est donc pas une utopie stérile mais un \textbf{idéal régulateur}, dont les avancées concrètes — Union africaine, intégrations régionales, ZLECAf — montrent la pertinence. Il appartient à la jeunesse africaine, par l'entreprise, le numérique et la culture, d'en poursuivre la réalisation patiente. Reste à savoir si cette exigence d'unité ne dessine pas, à l'échelle du monde, la forme future de la justice entre les peuples. »
\end{exemplebox}

\begin{exoresolu}[Sujet de dissertation : « Le panafricanisme est-il une utopie ? »]
Énoncé : rédiger la problématique, le plan détaillé et la conclusion d'une dissertation sur ce sujet.
\tcblower
\textbf{Problématique.} Le panafricanisme se présente comme la philosophie du développement de l'Afrique par l'unité ; mais soixante ans d'échecs et de divisions invitent à le juger irréaliste. Faut-il alors y voir un rêve sans prise sur le réel, ou l'horizon nécessaire de tout développement africain ?

\textbf{Plan dialectique.}
\begin{itemize}
\item \emph{Thèse} : le panafricanisme est une utopie — divisions linguistiques et politiques, échec du projet unioniste de 1963, dépendances persistantes, faiblesse du commerce intra-africain.
\item \emph{Antithèse} : le panafricanisme est une nécessité — seule l'unité confère un poids démographique (2,5 milliards d'habitants en 2050), économique (ZLECAf) et diplomatique (une voix à l'ONU).
\item \emph{Synthèse} : le panafricanisme est un idéal régulateur (référence possible : Kant, \emph{Idée d'une histoire universelle}) — un horizon qui guide des avancées concrètes et progressives, dont la jeunesse est aujourd'hui l'acteur principal (entrepreneuriat, numérique, culture).
\end{itemize}

\textbf{Conclusion.} Reprendre la conclusion type proposée ci-dessus : réponse nuancée à la problématique, bilan des trois moments, ouverture vers la solidarité mondiale comme exigence de justice universelle.
\end{exoresolu}

\begin{exercice}[À faire]
\begin{enumerate}
\item Définir en deux phrases chacun des termes suivants : utopie, idéal régulateur, développement.
\item Construire le tableau « utopie / nécessité / synthèse » pour le sujet : « L'unité africaine est-elle possible ? »
\item Rédiger une conclusion complète (8 à 12 lignes) pour le sujet : « Le développement de l'Afrique passe-t-il nécessairement par son unité ? »
\end{enumerate}
\end{exercice}

\begin{corrige}[Exercice 1]
\textbf{1.} L'\emph{utopie} est un projet idéal qui ne correspond à aucune réalité possible ; elle décrit un « nulle part ». L'\emph{idéal régulateur} est un but jamais pleinement atteint mais qui oriente l'action et produit des avancées réelles. Le \emph{développement} est le processus global (économique, social, politique, culturel) d'amélioration durable des conditions de vie d'une population.

\textbf{2.} \emph{Utopie} : frontières coloniales, conflits, égoïsmes des souverainetés. \emph{Nécessité} : marché unique, poids diplomatique, mutualisation des ressources. \emph{Synthèse} : l'unité est possible comme processus progressif (intégrations régionales, ZLECAf), non comme État unique immédiat.

\textbf{3.} Attendus : réponse claire à la problématique (oui, mais une unité progressive et non totalitaire), bilan des thèses confrontées, ouverture (par exemple vers la justice mondiale ou le rôle de la jeunesse). Aucun argument nouveau ne doit apparaître.
\end{corrige}

\newpage

\coursec{Activité d'intégration}

\coursec{Le Panafricanisme et le développement de l'Afrique}

\courssub{Objectifs de l'activité}

À l'issue de ce travail pratique guidé, tu seras capable de :
\begin{itemize}
\item mobiliser les définitions du \cle{développement} et du \cle{sous-développement} pour analyser une situation concrète ;
\item construire et interpréter une carte schématique et un diagramme en barres à partir de données chiffrées ;
\item appliquer les notions de la leçon (\cle{panafricanisme}, \cle{unité africaine}, \cle{intégration régionale}) à la situation du Cameroun ;
\item rédiger et justifier une conclusion argumentée (plaidoyer) comportant au moins deux arguments et un exemple camerounais précis.
\end{itemize}

\courssub{Situation}

\textbf{Contexte.} En mars 2025, le club de philosophie du Lycée technique de Douala participe à la « Semaine de l'Afrique » organisée par le MINRESI. Le thème retenu est : \emph{« L'Afrique unie en 2050 : utopie ou condition du développement ? »}. Ton groupe de quatre élèves doit présenter un panneau composé d'une carte, d'un diagramme et d'un plaidoyer écrit.

\textbf{Dossier documentaire remis au groupe.}

\emph{Document 1 — Quelques données sur les échanges commerciaux de l'Afrique (source : CNUCED, BAD, valeurs arrondies).}
\begin{center}
\begin{tabularx}{\linewidth}{|l|X|}\hline
\rowcolor{popblueL} \textbf{Indicateur} & \textbf{Valeur} \\ \hline
Part des échanges intra-africains dans le commerce total de l'Afrique & environ 15 \% \\ \hline
Part des échanges intra-européens dans le commerce de l'Union européenne & environ 60 \% \\ \hline
Part des échanges intra-asiatiques dans le commerce de l'Asie & environ 55 \% \\ \hline
Population de l'Afrique (marché potentiel de la ZLECAf) & environ 1,4 milliard d'habitants \\ \hline
Nombre de pays signataires de la ZLECAf (Zone de libre-échange continentale africaine) & 54 sur 55 \\ \hline
\end{tabularx}
\end{center}

\emph{Document 2 — Le Cameroun dans l'intégration africaine.} Le Cameroun est membre de la CEMAC (Communauté économique et monétaire de l'Afrique centrale, 6 États). Le corridor Douala-N'Djamena (environ \SI{2000}{km} de route et de voie ferrée) achemine l'essentiel des importations du Tchad via le port de Douala. Le Cameroun a signé la ZLECAf en 2018 et l'a ratifiée en 2020. Pourtant, les échanges entre pays de la CEMAC ne représentent qu'environ 3 \% de leur commerce total : le Cameroun exporte surtout pétrole brut, cacao et bois vers l'Europe et l'Asie, et importe des produits manufacturés.

\emph{Document 3 — Extrait du discours de Kwame Nkrumah (Accra, 1963) :} « L'indépendance politique n'a de sens que si elle s'accompagne de l'unité économique du continent. Aucun de nos pays, pris isolément, ne peut négocier à égalité avec les grandes puissances. »

\emph{Document 4 — Agenda 2063 de l'Union africaine :} « Nous voulons une Afrique intégrée, prospère et pacifique, dirigée par ses propres citoyens et représentant une force dynamique sur la scène internationale. »

\emph{Document 5 — Photographies :} le port en eau profonde de Kribi, le barrage de Nachtigal, une file de camions bloqués plusieurs jours à un poste frontière entre deux pays voisins.

\textbf{Tâche finale.} Rédiger un plaidoyer d'une demi-page répondant à la question : \emph{« L'unité africaine est-elle la condition du développement du Cameroun ? »}

\courssub{Consignes}

\begin{enumerate}
\item À partir des documents 1 et 2, dégage deux manifestations du \cle{sous-développement} de l'Afrique et du Cameroun, en citant les chiffres qui les illustrent.
\item Construis, sur la base du canevas fourni, une carte schématique de l'Afrique localisant : les principales ressources (pétrole du golfe de Guinée, cuivre de Zambie/RDC, cacao d'Afrique de l'Ouest, phosphate du Maroc), les blocs régionaux (CEMAC, CEDEAO, SADC, UMA) et le corridor Douala-N'Djamena. Ajoute titre, légende et flèche du nord.
\item Réalise un diagramme en barres comparant les parts d'échanges intra-régionaux de l'Afrique, de l'Europe et de l'Asie (document 1). Que constates-tu ? Quelle hypothèse formules-tu pour expliquer cet écart ?
\item En t'appuyant sur les documents 3 et 4, explique en quelques lignes en quoi le panafricanisme se présente comme une \emph{philosophie du développement} : quel problème diagnostique-t-il, et quelle solution propose-t-il ?
\item À l'aide du document 5 et de tes connaissances, identifie deux \cle{difficultés} concrètes qui freinent l'unité africaine (frontières, infrastructures, nationalismes, dépendance extérieure).
\item Rédige le plaidoyer d'une demi-page : annonce ta thèse (oui, non, ou réponse nuancée), justifie-la par au moins \textbf{deux arguments} et \textbf{un exemple camerounais concret} (CEMAC, ZLECAf, corridor Douala-N'Djamena), puis conclus.
\end{enumerate}

\courssub{Ressources à mobiliser}

\begin{aretenir}[Notions utiles de la leçon]
\begin{itemize}
\item \textbf{Développement} : processus global de transformation (économique, sociale, humaine) qui améliore durablement les conditions de vie d'un peuple ; il se mesure notamment par le PIB, l'IDH, l'industrialisation, l'intégration commerciale.
\item \textbf{Sous-développement} : état de dépendance et de blocage caractérisé par l'exportation de matières premières brutes, la faiblesse des échanges internes, l'insuffisance des infrastructures et la dépendance extérieure.
\item \textbf{Panafricanisme} : mouvement et philosophie visant l'unité politique et économique de l'Afrique (Nkrumah, Senghor, Nyerere), concrétisée par l'OUA (1963), puis l'Union africaine (2002), les communautés régionales (CEMAC, CEDEAO) et la ZLECAf (2019).
\item \textbf{Difficultés} : héritage des frontières coloniales, nationalismes, conflits, faiblesse des infrastructures de liaison, concurrence des puissances extérieures.
\end{itemize}
\end{aretenir}

\begin{methode}[Rédiger un plaidoyer argumenté]
\begin{enumerate}
\item Poser clairement la question et annoncer sa thèse dès la première phrase.
\item Développer l'argument 1 (théorique ou historique), puis l'argument 2 (économique ou chiffré).
\item Illustrer par un exemple camerounais précis et daté.
\item Envisager une objection (les difficultés de l'unité) et y répondre.
\item Conclure par une phrase ferme qui répond à la question posée.
\end{enumerate}
\end{methode}

\courssub{Démarche et corrigé commenté}

\begin{exoresolu}[TP guidé : « L'Afrique unie en 2050 »]
Énoncé : construire la carte et le diagramme, puis rédiger le plaidoyer répondant à la question : « L'unité africaine est-elle la condition du développement du Cameroun ? »

\tcblower

\textbf{Étape 1 — Manifestations du sous-développement (consigne 1).}
Première manifestation : la \emph{faiblesse des échanges intra-africains} (environ 15 \% du commerce total, contre 60 \% en Europe), signe que les économies africaines commercent davantage avec l'extérieur qu'entre elles. Deuxième manifestation : la \emph{spécialisation dans les matières premières brutes} (pétrole, cacao, bois exportés sans transformation) et l'insuffisance de l'intégration régionale : seulement 3 \% d'échanges au sein de la CEMAC. Le Cameroun, « Afrique en miniature », illustre ce paradoxe : pays riche en ressources, mais dépendant des marchés extérieurs.

\textbf{Étape 2 — La carte schématique (consigne 2).} Le canevas ci-dessous donne le modèle attendu : contour simplifié du continent, ressources localisées par symboles, blocs régionaux colorés, corridor matérialisé par un trait épais.

\begin{popfigure}
\begin{tikzpicture}[scale=0.85, every node/.style={font=\scriptsize}]
% Contour très simplifié de l'Afrique (amélioré)
\draw[thick, popdark, fill=popcream] 
    (5.5,7.0) -- (7.0,7.2) -- (8.2,6.5) -- (9.0,5.5) -- (9.5,4.5) -- (9.8,3.5) -- (9.5,2.5) -- (8.8,1.5) -- (7.8,0.5) -- (6.8,-0.2) -- (5.8,-0.5) -- (4.8,-0.3) -- (4.0,0.5) -- (3.2,1.2) -- (2.5,2.0) -- (2.0,2.8) -- (1.8,3.5) -- (1.8,4.5) -- (2.2,5.2) -- (3.0,6.0) -- (4.0,6.5) -- cycle;
% Flèche Nord
\draw[-{Stealth[length=3mm]}, popink, thick] (1.0,6.0) -- (1.0,7.0) node[above]{N};
% Blocs régionaux (zones approximatives avec coins arrondis)
\fill[popblue!35, rounded corners=4pt] (2.2,6.2) rectangle (4.8,7.0);
\node[align=center] at (3.5,6.6) {\textbf{UMA}\\(Union du Maghreb)};
\fill[popgreen!35, rounded corners=4pt] (1.8,3.5) rectangle (4.2,5.5);
\node[align=center] at (3.0,4.5) {\textbf{CEDEAO}\\(Afrique de l'Ouest)};
\fill[poporange!45, rounded corners=4pt] (4.0,2.8) rectangle (6.0,4.8);
\node[align=center] at (5.0,3.8) {\textbf{CEMAC}\\(Afrique centrale)};
\fill[poppurple!30, rounded corners=4pt] (4.5,0.0) rectangle (7.5,2.5);
\node[align=center] at (6.0,1.25) {\textbf{SADC}\\(Afrique australe)};
% Ressources (symboles + étiquettes)
\node[popink, scale=1.5] at (3.8,3.6) {$\blacktriangle$};
\node[popink, left] at (3.7,3.6) {Pétrole (Golfe de Guinée)};
\node[popgold!80!black, scale=1.3] at (2.2,4.2) {$\bullet$};
\node[popgold!80!black, below] at (2.2,4.1) {Cacao};
\node[popdark, scale=1.2] at (5.8,1.8) {$\blacksquare$};
\node[popdark, right] at (5.9,1.8) {Cuivre (RDC--Zambie)};
\node[popmuted, scale=1.3] at (4.5,6.6) {$\diamond$};
\node[popmuted, right] at (4.6,6.6) {Phosphate (Maroc)};
% Corridor Douala-N'Djamena
\draw[line width=2.5pt, poppink] (4.1,3.5) -- (4.5,4.0) -- (5.0,4.5);
\fill[poppink] (4.1,3.5) circle (2pt) node[below left, poppink]{Douala};
\fill[poppink] (5.0,4.5) circle (2pt) node[right, poppink]{N'Djamena};
% Échelle indicative
\draw[popdark, thick] (1.0,0.0) -- (2.0,0.0) node[midway, below]{~2000 km};
\end{tikzpicture}
\legende{Croquis schématique de l'Afrique : 1. Blocs régionaux (UMA, CEDEAO, CEMAC, SADC) ; 2. Ressources majeures ($\blacktriangle$ pétrole, $\bullet$ cacao, $\blacksquare$ cuivre, $\diamond$ phosphate) ; 3. Corridor Douala--N'Djamena (trait rose). Contour approximatif, sans précision cartographique.}
\end{popfigure}

\textbf{Étape 3 — Le diagramme en barres (consigne 3).}

\begin{popfigure}
\begin{tikzpicture}
\begin{axis}[
  ybar, bar width=1.0cm, width=0.9\linewidth, height=5.5cm,
  ymin=0, ymax=70,
  ylabel={Part des échanges intra-régionaux (en \%)},
  symbolic x coords={Afrique, Asie, Europe}, xtick=data,
  nodes near coords={\pgfmathprintnumber\pgfplotspointmeta\,\%}, nodes near coords align={vertical},
  ymajorgrids, grid style={popline!40},
  every axis/.append style={font=\scriptsize, tick label style={font=\scriptsize}},
  enlarge x limits=0.25]
\addplot[fill=popblue!70, draw=popdark] coordinates {(Afrique,15) (Asie,55) (Europe,60)};
\end{axis}
\end{tikzpicture}
\legende{Comparaison des parts d'échanges intra-régionaux (données arrondies, CNUCED/BAD).}
\end{popfigure}

\textbf{Constat} : l'Afrique échange environ quatre fois moins entre ses propres pays (15 \%) que l'Europe (60 \%) ou l'Asie (55 \%). \textbf{Hypothèse explicative} : les économies africaines, héritières d'une organisation coloniale tournée vers l'exportation, produisent souvent les mêmes matières premières et manquent d'infrastructures de liaison (routes, rails, énergie) et d'industries de transformation ; elles ne peuvent donc pas commercer entre elles de manière significative.

\textbf{Étape 4 — Le panafricanisme comme philosophie du développement (consigne 4).} Le diagnostic de Nkrumah (document 3) est que la \emph{balkanisation} du continent rend chaque État trop faible pour négocier et se développer seul : l'indépendance politique sans unité économique reste une indépendance de façade. La solution proposée est l'\emph{unité} : mutualiser les ressources, créer un grand marché intérieur, industrialiser ensemble. L'Agenda 2063 (document 4) prolonge ce projet : faire de l'Afrique unie un acteur mondial. Le panafricanisme est donc bien une philosophie du développement : il lie la liberté politique à la prospérité collective.

\textbf{Étape 5 — Les difficultés (consigne 5).} Première difficulté : les \emph{frontières héritées de la colonisation} et les nationalismes, qui ralentissent la libre circulation des personnes et des biens (camions bloqués plusieurs jours aux frontières, document 5). Deuxième difficulté : le \emph{déficit d'infrastructures de liaison} et la dépendance extérieure : les routes et chemins de fer relient historiquement les mines aux ports (vers l'extérieur) plus que les pays africains entre eux. On peut ajouter les conflits armés et la concurrence des grandes puissances pour l'accès aux ressources.

\textbf{Étape 6 — Le plaidoyer rédigé (consigne 6), modèle de copie.}

\emph{« L'unité africaine n'est pas un rêve de poète : elle est la condition du développement du Cameroun. D'abord, parce que seul, notre pays est trop étroit : avec environ 30 millions d'habitants en 2025, son marché intérieur ne permet pas de rentabiliser de grandes industries de transformation ; uni à l'Afrique de la ZLECAf, il accède à un marché de 1,4 milliard de consommateurs, ce qui rend possibles la transformation locale du cacao et la création d'emplois qualifiés. Ensuite, parce que l'histoire et les chiffres le prouvent : Nkrumah l'avait annoncé en 1963, aucun pays africain isolé ne négocie à égalité avec les grandes puissances ; les données le confirment, l'Afrique n'échange que 15 \% de son commerce entre ses pays, contre 60 \% pour l'Europe unie. Le corridor Douala--N'Djamena montre la voie concrète : quand le port de Douala dessert le Tchad, le Cameroun devient le poumon économique de l'Afrique centrale, et la CEMAC, malgré ses 3 \% d'échanges internes actuels, dessine ce que pourrait être une intégration réussie. Certes, les frontières, les nationalismes et le manque de routes freinent encore l'unité ; mais ces obstacles ne sont pas des fatalités : ils désignent précisément le travail à accomplir. Nous concluons donc : oui, l'unité africaine est la condition du développement du Cameroun — non une unité miraculeuse, mais une unité construite pas à pas, corridor après corridor, traité après traité. »}

\textbf{Commentaire.} Ce plaidoyer respecte la méthode : thèse annoncée dès la première phrase ; argument 1 économique (taille du marché, ZLECAf) ; argument 2 historique et chiffré (Nkrumah, 15 \% contre 60 \%) ; exemple camerounais concret et daté (corridor Douala--N'Djamena, CEMAC, population 2025) ; objection envisagée et réfutée (obstacles = travail à faire) ; conclusion ferme et nuancée.
\end{exoresolu}

\courssub{Bilan de l'activité}

Cette activité a permis de mettre en évidence trois acquis essentiels. D'abord, les notions de \cle{développement} et de \cle{sous-développement} ne sont pas des définitions abstraites : elles se lisent dans des données concrètes (15 \% d'échanges intra-africains, 3 \% au sein de la CEMAC) et dans des réalités quotidiennes (matières premières exportées brutes, camions bloqués aux frontières). Ensuite, le \cle{panafricanisme} apparaît comme une réponse rationnelle et argumentée à un diagnostic précis : la division du continent entretient sa dépendance ; l'unité est pensée comme levier de développement, de Nkrumah à l'Agenda 2063 et à la ZLECAf. Enfin, l'activité a entraîné deux savoir-faire exigibles au Probatoire et au Baccalauréat technique : \emph{appliquer les notions à une situation concrète} (carte, diagramme, exemples camerounais) et \emph{rédiger une conclusion justifiée} selon la logique de la dissertation : thèse, arguments, exemple, objection, conclusion nuancée. Tu es désormais prêt à traiter un sujet tel que : « L'unité africaine est-elle une utopie ? » ou « Le développement de l'Afrique passe-t-il nécessairement par son unité ? »

\newpage

\coursec{Méthodes et exercices résolus}

\coursec{Le Panafricanisme et le développement de l'Afrique}

\courssub{Fiche méthode 1 : Définir et distinguer développement et sous-développement}

\begin{methode}[Définir rigoureusement les concepts de développement et de sous-développement]
\begin{enumerate}
\item \textbf{Définir le développement} en croisant deux dimensions : la dimension \cle{quantitative} (croissance économique, augmentation du PIB, industrialisation) et la dimension \cle{qualitative} (amélioration des conditions de vie : santé, éducation, libertés, dignité). Retenir la formule de François Perroux : le développement est « la combinaison des changements mentaux et sociaux qui rendent un peuple apte à faire croître son produit réel global ».
\item \textbf{Définir le sous-développement} non comme un simple retard, mais comme un \cle{état structurel} : économie extravertie (tournée vers l'extérieur), dépendance, dualisme (secteur moderne côtoyant un secteur traditionnel stagnant).
\item \textbf{Mobiliser les indicateurs} : PIB par habitant, IDH (Indice de Développement Humain : santé, éducation, niveau de vie), taux d'alphabétisation, espérance de vie.
\item \textbf{Éviter la confusion} : croissance $\neq$ développement. Un pays peut croître (exportation du pétrole, du cacao) sans se développer si les richesses ne profitent pas à la population.
\item \textbf{Illustrer} chaque définition par un exemple camerounais ou africain précis.
\end{enumerate}
\end{methode}

\begin{exoresolu}[Définir et distinguer]
\textbf{Énoncé.} Un élève affirme : « Le Cameroun exporte du pétrole et du cacao, son économie croît, donc le Cameroun est développé. » En vous appuyant sur les définitions du développement et du sous-développement, montrez en quatre étapes que ce raisonnement est insuffisant.
\tcblower
\textbf{Solution détaillée.}

\textbf{Étape 1 : identifier la thèse de l'élève.} L'élève identifie le développement à la seule \cle{croissance économique} mesurée par les exportations. C'est une réduction du développement à sa dimension quantitative.

\textbf{Étape 2 : rappeler la définition complète du développement.} Le développement est un processus global qui combine la croissance économique \emph{et} les transformations sociales, culturelles et mentales améliorant les conditions de vie de \emph{toute} la population (Perroux, Amartya Sen : le développement comme « processus d'élargissement des libertés réelles »).

\textbf{Étape 3 : confronter la thèse aux faits.} L'exportation de matières premières (pétrole, cacao, café, bois) caractérise précisément une économie \cle{extravertie} et \cle{dépendante}, signe de sous-développement : les richesses sont vendues brutes, sans transformation locale, et les revenus profitent souvent à une minorité. Or au Cameroun, une partie importante de la population vit sous le seuil de pauvreté, avec un accès inégal à la santé et à l'éducation.

\textbf{Étape 4 : conclure.} Le raisonnement de l'élève est donc insuffisant : la croissance est une \emph{condition} du développement, non sa \emph{totalité}. Tant que la richesse produite ne se traduit pas par une élévation durable du niveau de vie de l'ensemble de la population (IDH élevé), on ne peut parler de développement.

\textbf{Phrase de conclusion.} Le développement ne se mesure pas seulement en tonnes de cacao exportées, mais en libertés réelles et en dignité offertes à chaque citoyen.
\end{exoresolu}

\courssub{Fiche méthode 2 : Analyser une situation concrète de sous-développement}

\begin{methode}[Appliquer les notions de l'unité à une situation concrète de la vie courante]
\begin{enumerate}
\item \textbf{Lire la situation} et repérer les faits observables (prix, infrastructures, modes de vie, échanges).
\item \textbf{Classer les faits} selon les \cle{manifestations} du sous-développement : économiques (pauvreté, chômage, économie informelle), sociales (malnutrition, faible scolarisation), politiques (mauvaise gouvernance), culturelles (dépendance aux modèles étrangers).
\item \textbf{Identifier les causes} en distinguant causes \cle{internes} (gouvernance, corruption, guerres, faible industrialisation) et causes \cle{externes} (héritage colonial, échanges inégaux, endettement, domination des multinationales).
\item \textbf{Relier causes et manifestations} par un raisonnement explicite : « parce que..., alors... ».
\item \textbf{Proposer une piste de solution} cohérente avec le panafricanisme (unité, transformation locale, solidarité africaine).
\end{enumerate}
\end{methode}

\begin{exoresolu}[Analyser une situation concrète]
\textbf{Énoncé.} Dans un village de la région du Centre, les planteurs vendent le kilogramme de cacao brut à bas prix à des acheteurs étrangers ; le chocolat fabriqué en Europe est ensuite revendu au Cameroun à un prix dix fois supérieur. Les jeunes du village, sans emploi, rêvent de partir à l'étranger. Analysez cette situation à l'aide des notions de sous-développement (manifestations et causes) et dégagez une piste de solution.
\tcblower
\textbf{Solution détaillée.}

\textbf{Étape 1 : repérage des faits.} Vente de matière première brute à bas prix ; importation du produit transformé à prix élevé ; chômage des jeunes ; tentation de l'émigration.

\textbf{Étape 2 : manifestations du sous-développement.}
\begin{itemize}
\item \emph{Économique} : économie de rente et d'exportation de produits bruts, sans industrie de transformation ; chômage des jeunes.
\item \emph{Sociale} : appauvrissement des producteurs malgré leur travail ; exode des forces vives.
\item \emph{Culturelle} : dépendance aux produits et modèles étrangers (le chocolat consommé vient d'Europe).
\end{itemize}

\textbf{Étape 3 : causes.}
\begin{itemize}
\item \emph{Causes externes} : héritage de l'économie coloniale (l'Afrique fournit les matières premières, l'Occident transforme) ; échanges inégaux (les cours des matières premières sont fixés par les marchés étrangers).
\item \emph{Causes internes} : absence d'unités de transformation locales (usines de chocolat), manque d'investissement et de formation.
\end{itemize}

\textbf{Étape 4 : relation cause $\rightarrow$ manifestation.} Parce que le cacao est vendu brut, la \cle{valeur ajoutée} (le profit de la transformation) échappe au pays ; parce que le pays n'industrialise pas, les jeunes ne trouvent pas d'emploi et émigrent, ce qui prive le village de ses forces vives et aggrave le sous-développement : c'est un \cle{cercle vicieux}.

\textbf{Étape 5 : piste de solution.} Dans l'esprit du panafricanisme : créer des coopératives de planteurs, transformer le cacao sur place (chocolat \emph{made in Cameroon}), et s'appuyer sur le marché commun africain (ZLECAf) pour vendre le produit fini aux pays voisins.

\textbf{Phrase de conclusion.} Cette situation montre que le sous-développement n'est pas une fatalité : il résulte de structures que l'unité d'action africaine et la transformation locale peuvent briser.
\end{exoresolu}

\courssub{Fiche méthode 3 : Construire une frise historique du panafricanisme}

\begin{methode}[Restituer l'historique du panafricanisme]
\begin{enumerate}
\item \textbf{Fixer l'origine} : le panafricanisme naît au tournant du \textsc{xx}\textsuperscript{e} siècle dans la \cle{diaspora} africaine (Amérique, Caraïbes) : W.E.B. Du Bois, Marcus Garvey, puis Aimé Césaire et Léopold Sédar Senghor avec la \cle{Négritude}.
\item \textbf{Repérer les congrès panafricains} : premier congrès de Londres (1900), congrès de Manchester (1945) avec Kwame Nkrumah et Jomo Kenyatta, qui exige l'indépendance.
\item \textbf{Marquer le passage à l'action politique} : indépendances (années 1960), création de l'\cle{OUA} (Organisation de l'Unité Africaine) à Addis-Abeba en \cle{1963}.
\item \textbf{Actualiser} : transformation de l'OUA en \cle{Union Africaine} (UA) en 2002 ; projets d'intégration économique (ZLECAf, accord signé en 2018 à Kigali, entré en vigueur en mai 2019, échanges effectifs depuis janvier 2021).
\item \textbf{Présenter la frise} sous forme chronologique claire, avec dates, lieux et acteurs.
\end{enumerate}
\end{methode}

\begin{exoresolu}[Construire une frise historique]
\textbf{Énoncé.} « Le panafricanisme est né en Afrique au moment des indépendances. » Corrigez cette affirmation en retraçant, en quatre temps datés et argumentés, les grandes étapes de l'histoire du panafricanisme.
\tcblower
\textbf{Solution détaillée.}

\textbf{Étape 1 : corriger l'erreur.} L'affirmation est doublement inexacte : le panafricanisme n'est \emph{pas né en Afrique} mais dans la diaspora, et \emph{avant} les indépendances, dont il a été l'un des moteurs.

\textbf{Étape 2 : la naissance dans la diaspora (fin \textsc{xix}\textsuperscript{e} -- début \textsc{xx}\textsuperscript{e} siècle).} Descendants d'esclaves africains en Amérique et aux Caraïbes, des intellectuels comme W.E.B. Du Bois et Marcus Garvey dénoncent la colonisation et le racisme, et appellent à l'unité de tous les Noirs du monde. Le premier congrès panafricain se tient à \cle{Londres en 1900}.

\textbf{Étape 3 : la radicalisation (1945).} Le congrès de \cle{Manchester (1945)}, avec Nkrumah et Kenyatta, passe de la revendication culturelle à l'exigence politique : l'\cle{indépendance immédiate} des colonies africaines. Parallèlement, la Négritude (Césaire, Senghor) affirme la valeur des cultures noires.

\textbf{Étape 4 : l'institutionnalisation (1963 -- 2002).} Après les indépendances, l'\cle{OUA} est créée à Addis-Abeba en \cle{1963} pour consolider l'unité du continent ; elle devient l'\cle{Union Africaine} en \cle{2002}, avec l'ambition d'une intégration économique et politique approfondie (ZLECAf depuis 2019/2021).

\textbf{Phrase de conclusion.} Né de la souffrance de la diaspora, le panafricanisme a précédé, préparé puis prolongé les indépendances : il est la philosophie de l'unité africaine, non son produit tardif.
\end{exoresolu}

\courssub{Fiche méthode 4 : Expliquer les objectifs du panafricanisme comme philosophie du développement}

\begin{methode}[Présenter le panafricanisme comme philosophie du développement]
\begin{enumerate}
\item \textbf{Définir le panafricanisme} : mouvement de pensée et d'action visant l'\cle{unité} des Africains du continent et de la diaspora pour leur libération et leur épanouissement.
\item \textbf{Dégager les objectifs politiques} : décolonisation, souveraineté, unité des États africains (Nkrumah : « l'Afrique doit s'unir ou périr »).
\item \textbf{Dégager les objectifs économiques} : transformation locale des matières premières, marché commun, monnaie commune, indépendance vis-à-vis des puissances étrangères.
\item \textbf{Dégager les objectifs culturels} : revalorisation des cultures africaines (Négritude), confiance en soi, refus de l'imitation servile.
\item \textbf{Montrer le lien avec le développement} : seule l'unité donne à l'Afrique la taille critique (marché, ressources, force de négociation) pour échapper au sous-développement.
\end{enumerate}
\end{methode}

\begin{exoresolu}[Expliquer les objectifs du panafricanisme]
\textbf{Énoncé.} « Le panafricanisme n'est pas seulement un rêve d'unité politique : c'est une philosophie du développement. » Justifiez cette affirmation en dégageant les objectifs politiques, économiques et culturels du panafricanisme et leur lien avec le développement de l'Afrique.
\tcblower
\textbf{Solution détaillée.}

\textbf{Étape 1 : définir.} Le panafricanisme est le mouvement qui promeut la solidarité et l'unité de tous les peuples africains et d'ascendance africaine, en vue de leur libération et de leur développement.

\textbf{Étape 2 : objectifs politiques.} Le panafricanisme vise d'abord la \cle{libération} (fin de la colonisation et de l'apartheid) puis l'\cle{unité politique} : pour Nkrumah, des États africains divisés restent faibles et manipulables ; unis, ils pèsent dans les relations internationales. La souveraineté politique est la condition de toute politique de développement autonome.

\textbf{Étape 3 : objectifs économiques.} Divisée en plus de cinquante petits marchés, l'Afrique ne peut industrialiser. Le panafricanisme propose l'\cle{intégration économique} : marché commun (ZLECAf), transformation locale des ressources, coopération entre États. L'unité crée la \cle{taille critique} nécessaire au développement : négocier ensemble les prix des matières premières, mutualiser les infrastructures (routes, énergie).

\textbf{Étape 4 : objectifs culturels.} Avec la Négritude (Senghor, Césaire), le panafricanisme restaure la \cle{dignité culturelle} africaine. Or le développement exige la confiance en soi : un peuple qui méprise ses propres langues et savoirs ne peut développer que des copies. Le développement endogène suppose l'ancrage culturel.

\textbf{Étape 5 : synthèse.} Politique (souveraineté), économique (intégration) et culturel (dignité) : les trois objectifs convergent vers un même but, faire de l'Afrique un acteur de son propre développement et non un objet des puissances étrangères.

\textbf{Phrase de conclusion.} Le panafricanisme est bien une philosophie du développement, car il affirme que l'unité est la condition première de la liberté, de la prospérité et de la dignité africaines.
\end{exoresolu}

\courssub{Fiche méthode 5 : Évaluer les difficultés et limites du panafricanisme}

\begin{methode}[Analyser les difficultés du panafricanisme]
\begin{enumerate}
\item \textbf{Classer les difficultés} en catégories : politiques (souveraineté jalouse des États, frontières héritées de la colonisation, conflits), économiques (économies concurrentes exportant les mêmes produits bruts, faibles échanges intra-africains, dépendance financière), culturelles et linguistiques (diversité des langues, héritages anglophone/francophone/lusophone), externes (ingérences étrangères, néocolonialisme).
\item \textbf{Illustrer chaque difficulté} par un fait précis (ex. : les échanges entre pays africains restent faibles comparés aux échanges avec l'Europe ou la Chine).
\item \textbf{Nuancer} : une difficulté n'est pas une impossibilité ; montrer les progrès (UA, ZLECAf, libre circulation dans certaines régions grâce au protocole de 2018).
\item \textbf{Éviter deux excès} : l'afro-pessimisme (le panafricanisme est mort) et l'angélisme (l'unité est déjà réalisée).
\end{enumerate}
\end{methode}

\begin{exoresolu}[Évaluer les difficultés du panafricanisme]
\textbf{Énoncé.} « Plus de soixante ans après la création de l'OUA, l'unité africaine reste un chantier inachevé. » Identifiez et expliquez quatre difficultés majeures qui freinent la réalisation du panafricanisme, puis montrez pourquoi ces difficultés ne condamnent pas le projet.
\tcblower
\textbf{Solution détaillée.}

\textbf{Étape 1 : difficulté politique.} Les États africains, jaloux de leur souveraineté récemment conquise, refusent de déléguer leur pouvoir à des institutions supranationales fortes ; les frontières héritées de la colonisation (conférence de Berlin, 1884-1885) divisent des peuples frères et nourrissent des conflits.

\textbf{Étape 2 : difficulté économique.} Les économies africaines sont \cle{concurrentes} plutôt que complémentaires : elles exportent toutes des matières premières brutes vers l'extérieur et échangent peu entre elles. La dépendance financière (aide, dette, institutions internationales) limite l'autonomie des décisions.

\textbf{Étape 3 : difficulté culturelle et linguistique.} Le partage linguistique hérité de la colonisation (pays anglophones, francophones, lusophones) et la diversité ethnique compliquent la construction d'une conscience commune ; certains régimes instrumentalisent même l'ethnicité.

\textbf{Étape 4 : difficulté externe.} Le \cle{néocolonialisme} : les anciennes puissances coloniales et les nouvelles puissances (multinationales, Chine) entretiennent la division africaine, plus profitable à leurs intérêts (accords bilatéraux, soutien à des régimes dociles).

\textbf{Étape 5 : nuance.} Ces difficultés n'invalident pas le projet : l'Union Africaine existe et s'affirme (missions de paix), la ZLECAf (2019/2021) crée un marché commun de plus d'un milliard de consommateurs, et la jeunesse africaine, connectée, fait revivre l'idéal panafricain. L'histoire montre que les grandes unifications sont toujours lentes (l'Europe elle-même a mis des décennies).

\textbf{Phrase de conclusion.} Le panafricanisme est un chantier inachevé, non un projet abandonné : ses difficultés mesurent l'ampleur de la tâche, non son impossibilité.
\end{exoresolu}

\courssub{Fiche méthode 6 : Rédiger une conclusion de dissertation justifiée}

\begin{methode}[Rédiger et justifier une démarche, une réponse ou une conclusion]
\begin{enumerate}
\item \textbf{Répondre directement à la question posée} en une phrase claire (thèse assumée), sans « peut-être » évasif.
\item \textbf{Résumer la démarche} : rappeler en deux ou trois phrases les étapes du raisonnement (thèse examinée, antithèse, dépassement).
\item \textbf{Justifier la position} par l'argument le plus fort du développement, reformulé brièvement.
\item \textbf{Nuancer si nécessaire} : préciser les conditions ou limites de la réponse (ex. : « à condition que... »).
\item \textbf{Ouvrir} éventuellement sur une question nouvelle liée au sujet, sans introduire de notion hors programme.
\end{enumerate}
\end{methode}

\begin{exoresolu}[Rédiger une conclusion justifiée]
\textbf{Énoncé.} Sujet de dissertation : « Le panafricanisme est-il une utopie ? » Rédigez une conclusion complète (8 à 12 lignes) qui réponde à la question, résume la démarche et justifie votre position.
\tcblower
\textbf{Solution détaillée (modèle de conclusion).}

Au terme de cette réflexion, nous pouvons affirmer que le panafricanisme n'est pas une utopie, mais un \cle{projet} dont la réalisation, difficile, est déjà partiellement engagée. Notre démarche a d'abord montré que le panafricanisme, né dans la diaspora avec Du Bois et la Négritude, répond à un besoin réel : dépasser la division coloniale qui maintient l'Afrique dans le sous-développement. Nous avons ensuite reconnu les obstacles considérables qui le freinent : souverainetés jalouses, économies concurrentes, ingérences étrangères — obstacles qui donnent au projet des allures de rêve inaccessible. Mais nous avons établi, enfin, qu'une difficulté n'est pas une impossibilité : la création de l'OUA puis de l'Union Africaine, l'entrée en vigueur de la ZLECAf et l'éveil de la jeunesse prouvent que l'unité africaine progresse, lentement mais sûrement. L'utopie est ce qui ne peut advenir nulle part ; or l'unité africaine advient déjà, par étapes. À condition que les États acceptent de partager une part de leur souveraineté et que les peuples s'approprient l'idéal de leurs pères fondateurs, le panafricanisme demeure la voie la plus sûre vers le développement du continent. Reste ouverte une question que l'avenir tranchera : l'Afrique saura-t-elle transformer son unité rêvée en puissance réelle ?

\textbf{Commentaire de la copie.} Cette conclusion respecte les cinq étapes : réponse directe (« n'est pas une utopie, mais un projet »), résumé de la démarche en trois temps, justification par l'argument décisif (distinction difficulté/impossibilité), nuance conditionnelle (« à condition que... »), ouverture pertinente.
\end{exoresolu}

\begin{attention}[Les 5 erreurs qui coûtent des points au Bac]
\begin{enumerate}
\item \textbf{Confondre croissance et développement.} Affirmer qu'un pays « est développé parce que son PIB augmente » est une erreur de définition : exigez toujours la dimension humaine (IDH, santé, éducation, libertés).
\item \textbf{Faire naître le panafricanisme en Afrique en 1960.} Erreur historique classique : le panafricanisme naît dans la \cle{diaspora} (Londres 1900, Manchester 1945) ; l'OUA (1963) n'en est que l'institutionnalisation.
\item \textbf{Citer des dates et des noms faux ou approximatifs.} Écrire « l'OUA créée en 1945 » ou « Nkrumah, président du Nigeria » discrédite toute la copie. Retenir : Londres 1900, Manchester 1945, OUA 1963 (Addis-Abeba), UA 2002 ; Nkrumah $=$ Ghana.
\item \textbf{Présenter le sous-développement comme une fatalité naturelle.} Oublier les causes historiques et structurelles (colonisation, échanges inégaux, néocolonialisme) revient à essentialiser la pauvreté africaine : c'est une faute d'analyse.
\item \textbf{Conclure sans répondre à la question.} Une conclusion qui se contente de répéter le plan (« nous avons vu d'abord..., ensuite... ») sans trancher (« le panafricanisme est / n'est pas... ») ni justifier ce choix perd l'essentiel des points de synthèse.
\end{enumerate}
\end{attention}

\newpage

\coursec{Exercices}

\courssub{Je vérifie mes connaissances}

\begin{exercice}[QCM : notions fondamentales]
Pour chaque question, une seule réponse est exacte. Entoure la lettre correspondante.
\begin{enumerate}
\item Le développement se définit d'abord comme :
\begin{enumerate}
\item[a)] la seule croissance du PIB d'un pays ;
\item[b)] un processus global de transformation économique, sociale, culturelle et politique améliorant les conditions de vie des populations ;
\item[c)] l'imitation du mode de vie occidental ;
\item[d)] l'augmentation des exportations de matières premières.
\end{enumerate}
\item Le panafricanisme est né historiquement :
\begin{enumerate}
\item[a)] en Afrique du Sud après 1994 ;
\item[b)] dans la diaspora africaine (Amériques et Caraïbes) à la fin du XIX\textsuperscript{e} siècle ;
\item[c)] à la création de la ZLECAf en 2018 ;
\item[d)] lors de la Conférence de Berlin en 1885.
\end{enumerate}
\item L'OUA (Organisation de l'unité africaine) a été créée en :
\begin{enumerate}
\item[a)] 1945 à Manchester ; \quad b) 1963 à Addis-Abeba ; \quad c) 2002 à Durban ; \quad d) 1957 à Accra.
\end{enumerate}
\item Le sous-développement de l'Afrique s'explique NOTAMMENT par :
\begin{enumerate}
\item[a)] l'absence totale de ressources naturelles ;
\item[b)] la paresse des populations ;
\item[c)] un ensemble de causes internes (mauvaise gouvernance) et externes (héritage colonial, échanges inégaux) ;
\item[d)] le climat uniquement.
\end{enumerate}
\end{enumerate}
\end{exercice}

\begin{exercice}[Vrai ou faux justifié]
Indique si chaque affirmation est vraie ou fausse, puis justifie ta réponse en deux ou trois lignes.
\begin{enumerate}
\item « Le PIB par habitant suffit à mesurer le développement d'un pays. »
\item « Le panafricanisme vise uniquement l'unité politique des États africains. »
\item « Les échanges commerciaux entre pays africains sont aujourd'hui plus importants que les échanges de l'Afrique avec le reste du monde. »
\item « La ZLECAf (Zone de libre-échange continentale africaine) s'inscrit dans la logique panafricaniste. »
\end{enumerate}
\end{exercice}

\begin{exercice}[Définitions]
Définis en quelques lignes chacune des notions suivantes, en donnant pour chacune un exemple concret tiré du contexte camerounais ou africain :
\begin{enumerate}
\item le développement ;
\item le sous-développement ;
\item le panafricanisme ;
\item le développement endogène.
\end{enumerate}
\end{exercice}

\begin{exercice}[Repères chronologiques]
Replace sur une frise chronologique les événements suivants et indique pour chacun son importance pour le panafricanisme :
\begin{enumerate}
\item le 5\textsuperscript{e} Congrès panafricain de Manchester (1945) ;
\item l'indépendance du Ghana de Kwame Nkrumah (1957) ;
\item la création de l'OUA (1963) ;
\item la transformation de l'OUA en Union africaine (2002) ;
\item l'entrée en vigueur de la ZLECAf (2021).
\end{enumerate}
\end{exercice}

\courssub{Je m'entraîne}

\begin{exercice}[Indicateurs du développement]
Le tableau ci-dessous présente quelques indicateurs (valeurs approchées, sources : PNUD et Banque mondiale, données récentes).

\begin{center}
\begin{tabularx}{\linewidth}{|l|X|X|X|}
\hline
\rowcolor{popblueL} \textbf{Pays} & \textbf{PIB/habitant (USD)} & \textbf{IDH} & \textbf{Espérance de vie (ans)} \\
\hline
Cameroun & \num{1600} & 0,58 & 61 \\
\hline
Afrique du Sud & \num{6800} & 0,71 & 65 \\
\hline
France & \num{44000} & 0,90 & 83 \\
\hline
\end{tabularx}
\end{center}

\begin{enumerate}
\item À partir du tableau, dégage deux manifestations du sous-développement du Cameroun.
\item Montre, à l'aide d'un exemple, qu'un PIB par habitant élevé ne suffit pas à garantir le bien-être de toute la population.
\item Propose deux autres indicateurs (sociaux ou culturels) qui permettraient de compléter ce tableau.
\end{enumerate}
\end{exercice}

\begin{exercice}[Causes du sous-développement]
Voici une liste de facteurs : (a) la faible transformation locale des matières premières ; (b) la corruption et la mauvaise gouvernance ; (c) les frontières héritées de la colonisation ; (d) les termes inégaux de l'échange international ; (e) la faiblesse de l'industrialisation ; (f) les conflits armés.
\begin{enumerate}
\item Classe ces facteurs en deux catégories : causes internes et causes externes du sous-développement africain.
\item Choisis deux facteurs et explique précisément comment chacun freine le développement du Cameroun.
\item Pour chacun de ces deux facteurs, propose une solution concrète inspirée du panafricanisme.
\end{enumerate}
\end{exercice}

\begin{exercice}[Objectifs du panafricanisme]
À partir du cours, rédige un texte argumenté de 15 à 20 lignes répondant à la question : « En quoi le panafricanisme peut-il être considéré comme une philosophie du développement de l'Afrique ? » Tu mobiliseras au moins deux objectifs du panafricanisme (unité politique, intégration économique, dignité culturelle, solidarité entre Africains et diaspora) et un auteur cité au programme (Kwame Nkrumah, Léopold Sédar Senghor, Cheikh Anta Diop ou Marcus Garvey).
\end{exercice}

\begin{exercice}[Situation concrète : la transformation du cacao]
Le Cameroun exporte l'essentiel de son cacao sous forme de fèves brutes, puis importe du chocolat transformé à prix élevé.
\begin{enumerate}
\item Explique en quoi cette situation illustre le paradoxe d'une Afrique « riche mais sous-développée ».
\item Identifie la cause du sous-développement que cette situation révèle.
\item Propose deux mesures concrètes, à l'échelle nationale puis à l'échelle continentale (panafricaniste), pour transformer cette situation.
\end{enumerate}
\end{exercice}

\courssub{Je me prépare au Bac}

\begin{exercice}[Dissertation : le panafricanisme est-il une utopie ?]
\textbf{Sujet :} « Le panafricanisme est-il une utopie ? »

Travail demandé :
\begin{enumerate}
\item Rédige une introduction complète (accroche, présentation du sujet, définition des termes « panafricanisme » et « utopie », problématique, annonce du plan).
\item Rédige le développement en deux parties :
\begin{itemize}
\item \textbf{Thèse :} le panafricanisme apparaît comme une utopie (difficultés : divisions des États, frontières coloniales, faiblesse des échanges intra-africains, dépendance extérieure, échecs partiels de l'OUA) ;
\item \textbf{Antithèse :} le panafricanisme est une voie réaliste et nécessaire (réussites de l'intégration, Union africaine, ZLECAf, exemples d'intégrations régionales, force du marché continental).
\end{itemize}
\item Rédige une conclusion dans laquelle tu prends clairement position (synthèse : le panafricanisme comme projet en construction, condition du développement endogène de l'Afrique).
\end{enumerate}
\emph{Barème indicatif : introduction 4 pts ; développement 12 pts ; conclusion 4 pts.}
\end{exercice}

\begin{exercice}[Commentaire de texte : Kwame Nkrumah]
\textbf{Texte :} « L'Afrique doit s'unir. Nous avons toutes les ressources nécessaires : les minerais, les terres fertiles, les hommes. Divisés, nous restons faibles et exploités ; unis, nous pouvons parler d'égal à égal avec les autres puissances du monde. L'indépendance politique n'est qu'un premier pas : sans l'unité économique du continent, elle demeure une indépendance en trompe-l'œil. Aucun de nos pays, pris isolément, n'est assez fort pour négocier seul son avenir. » (inspiré de Kwame Nkrumah, \emph{L'Afrique doit s'unir}, 1963)

Travail demandé :
\begin{enumerate}
\item Dégage la thèse de l'auteur et formule-la en une phrase.
\item Releve et explique trois arguments qui soutiennent cette thèse.
\item Montre en quoi ce texte fait du panafricanisme une philosophie du développement et non seulement un projet politique.
\item Discussion : cet appel à l'unité, lancé en 1963, est-il encore d'actualité à l'heure de la ZLECAf ? Justifie ta réponse par deux arguments personnels.
\end{enumerate}
\emph{Barème indicatif : thèse 3 pts ; arguments 6 pts ; analyse 5 pts ; discussion 6 pts.}
\end{exercice}

\begin{exercice}[Analyse de documents et étude de cas : la ZLECAf]
\textbf{Document 1.} Quelques données sur le commerce africain (valeurs approchées, sources : CNUCED, Banque africaine de développement) :

\begin{center}
\begin{tabularx}{\linewidth}{|X|X|}
\hline
\rowcolor{popblueL} \textbf{Indicateur} & \textbf{Valeur} \\
\hline
Part des échanges intra-africains dans le commerce total de l'Afrique & environ 15 \% \\
\hline
Part des échanges intra-européens dans le commerce de l'Europe & environ 60 \% \\
\hline
Part des matières premières dans les exportations africaines & environ 80 \% \\
\hline
Population du marché de la ZLECAf & environ 1,3 milliard d'habitants \\
\hline
Nombre de pays signataires de la ZLECAf & 54 \\
\hline
\end{tabularx}
\end{center}

\textbf{Document 2.} Entrée en vigueur en 2021, la Zone de libre-échange continentale africaine (ZLECAf) supprime progressivement les droits de douane entre pays africains afin de créer un marché unique de biens et de services.

Travail demandé :
\begin{enumerate}
\item À partir du document 1, dégage le paradoxe de l'Afrique : un continent riche en ressources mais sous-développé et peu intégré à son propre marché.
\item Explique comment la ZLECAf concrétise deux objectifs historiques du panafricanisme.
\item Étude de cas : « La ZLECAf peut-elle transformer l'économie camerounaise ? » Argumente en mobilisant les notions d'\textbf{unité} et de \textbf{développement endogène}, en envisageant deux opportunités (marché élargi, industrialisation locale) et deux difficultés (concurrence, infrastructures, gouvernance).
\item Conclus en indiquant deux conditions nécessaires pour que la ZLECAf serve réellement le développement du Cameroun.
\end{enumerate}
\emph{Barème indicatif : analyse du document 4 pts ; objectifs 4 pts ; argumentation 8 pts ; conclusion 4 pts.}
\end{exercice}

\newpage

\coursec{Corrigés des exercices}

\coursec{Le Panafricanisme et le développement de l'Afrique}

\courssub{Exercices corrigés et méthodes}

\begin{corrige}[Exercice 1 — QCM : notions fondamentales]
\begin{enumerate}
\item \textbf{Réponse b).} Le développement est un processus \emph{global} et \emph{multidimensionnel} : il dépasse la seule croissance économique (réponse a insuffisante) pour inclure les transformations sociales (santé, éducation), culturelles et politiques (démocratie, bonne gouvernance). Il ne se réduit ni à l'imitation de l'Occident (c), ni à l'exportation de matières premières (d), qui est au contraire un trait du sous-développement.
\item \textbf{Réponse b).} Le panafricanisme naît dans la \cle{diaspora} africaine (Amériques, Caraïbes) à la fin du XIX\textsuperscript{e} siècle, avec des figures comme Marcus Garvey et W. E. B. Du Bois, avant de revenir sur le continent après 1945.
\item \textbf{Réponse b).} L'OUA est créée le 25 mai 1963 à Addis-Abeba (Éthiopie). Manchester (1945) est le 5\textsuperscript{e} Congrès panafricain ; Durban (2002) marque la naissance de l'Union africaine ; Accra (1957) est l'indépendance du Ghana.
\item \textbf{Réponse c).} Le sous-développement a des causes \emph{internes} (mauvaise gouvernance, corruption, conflits) et \emph{externes} (héritage colonial, échanges inégaux, faible transformation des matières premières). L'Afrique regorge de ressources (a faux) ; les réponses b et d relèvent de préjugés déterministes réfutés par l'analyse historique et économique.
\end{enumerate}
\textbf{Points de vigilance :} au Bac, un QCM peut être assorti d'une demande de justification ; il faut toujours être capable d'expliquer pourquoi les autres réponses sont fausses.
\end{corrige}

\begin{corrige}[Exercice 2 — Vrai ou faux justifié]
\begin{enumerate}
\item \textbf{Faux.} Le PIB par habitant est une moyenne qui ignore les inégalités de répartition et les dimensions sociales (santé, éducation). C'est pourquoi le PNUD a créé l'\cle{IDH}, qui combine revenu, espérance de vie et scolarisation. Un pays peut avoir un PIB élevé et une majorité de population pauvre.
\item \textbf{Faux.} Le panafricanisme vise certes l'unité politique, mais aussi l'\cle{intégration économique} (marché commun, ZLECAf), la \cle{dignité culturelle} (revalorisation des cultures africaines, Cheikh Anta Diop) et la solidarité avec la diaspora. C'est un projet global de développement.
\item \textbf{Faux.} Les échanges intra-africains ne représentent qu'environ 15 \% du commerce total de l'Afrique, contre environ 60 \% pour l'Europe. L'Afrique commerciale reste tournée vers l'extérieur : c'est précisément le problème que la ZLECAf veut corriger.
\item \textbf{Vrai.} En créant un marché continental unique de 1,3 milliard de consommateurs, la ZLECAf réalise l'objectif panafricaniste d'intégration économique défendu depuis Nkrumah : unir les économies africaines pour développer le continent de l'intérieur.
\end{enumerate}
\textbf{Points de vigilance :} une justification de deux ou trois lignes doit contenir un argument précis (donnée chiffrée, notion du cours, exemple), pas une simple reformulation de l'affirmation.
\end{corrige}

\begin{corrige}[Exercice 3 — Définitions]
\begin{enumerate}
\item \textbf{Le développement} : processus global de transformation économique, sociale, culturelle et politique d'une société, qui améliore durablement les conditions de vie de l'ensemble de la population. \emph{Exemple :} la construction d'écoles, d'hôpitaux et de routes au Cameroun, accompagnée d'une industrialisation croissante.
\item \textbf{Le sous-développement} : état d'une société caractérisé par la pauvreté, la dépendance extérieure, la faible industrialisation et l'insuffisance des services sociaux, malgré l'abondance possible des ressources. \emph{Exemple :} le Cameroun exporte son pétrole et son cacao bruts mais importe des produits finis coûteux.
\item \textbf{Le panafricanisme} : mouvement de pensée et d'action né dans la diaspora, visant l'unité politique, économique et culturelle de l'Afrique et de ses diasporas, afin de rendre au continent sa dignité et de fonder son développement. \emph{Exemple :} la création de l'OUA en 1963, puis de l'Union africaine en 2002.
\item \textbf{Le développement endogène} : développement fondé sur les ressources propres (humaines, naturelles, culturelles) d'un pays ou d'un continent, et non sur les modèles et l'aide extérieurs. \emph{Exemple :} transformer localement le cacao camerounais en chocolat au lieu d'exporter les fèves brutes.
\end{enumerate}
\textbf{Points de vigilance :} une bonne définition au Bac comporte le genre proche (processus, état, mouvement) et les attributs essentiels ; l'exemple doit être précis et correctement relié à la définition.
\end{corrige}

\begin{corrige}[Exercice 4 — Repères chronologiques]
Frise chronologique attendue :

\begin{popfigure}
\begin{tikzpicture}[scale=1, every node/.style={font=\small}]
\draw[-{Stealth}, thick, popdark] (0,0) -- (14.5,0) node[right]{temps};
\foreach \x/\d/\e in {1.2/1945/{Manchester}, 4.2/1957/{Ghana}, 7.2/1963/{OUA}, 10.4/2002/{UA}, 13.4/2021/{ZLECAf}}{
  \draw[thick, popblue] (\x,-0.15) -- (\x,0.15);
  \node[below, popdark] at (\x,-0.15) {\d};
  \node[above, popblue] at (\x,0.15) {\e};
}
\end{tikzpicture}
\legende{Frise chronologique des grandes étapes du panafricanisme (1945--2021).}
\end{popfigure}

Importance de chaque événement :
\begin{enumerate}
\item \textbf{Manchester (1945)} : le 5\textsuperscript{e} Congrès panafricain marque le passage du panafricanisme de la diaspora au continent ; Nkrumah et Kenyatta y réclament l'indépendance immédiate des colonies.
\item \textbf{Ghana (1957)} : première indépendance de l'Afrique noire colonisée ; Nkrumah en fait la base du combat pour l'unité africaine (« l'indépendance du Ghana est incomplète sans la libération de toute l'Afrique »).
\item \textbf{OUA (1963)} : première organisation continentale ; elle consacre l'unité politique mais aussi le respect des frontières héritées de la colonisation, ce qui limitera son ambition.
\item \textbf{Union africaine (2002)} : relance du projet panafricaniste avec des objectifs d'intégration économique et de gouvernance plus ambitieux (Parlement panafricain, Conseil de paix et de sécurité).
\item \textbf{ZLECAf (2021)} : concrétisation économique majeure du panafricanisme, créant un marché unique de 1,3 milliard de consommateurs.
\end{enumerate}
\textbf{Points de vigilance :} ne pas confondre Manchester (congrès de la diaspora et des nationalistes) avec Addis-Abeba (organisation des États indépendants).
\end{corrige}

\begin{corrige}[Exercice 5 — Indicateurs du développement]
\begin{enumerate}
\item \textbf{Deux manifestations du sous-développement du Cameroun :}
\begin{itemize}
\item un PIB par habitant de \num{1600} USD, soit environ 27 fois moins que celui de la France (\num{44000} USD) : $\dfrac{\num{44000}}{\num{1600}} = 27,5$ ;
\item un IDH de 0,58 et une espérance de vie de 61 ans, contre 0,90 et 83 ans en France : faiblesse des conditions sanitaires et éducatives.
\end{itemize}
\item \textbf{Limite du PIB par habitant :} l'Afrique du Sud affiche un PIB/habitant de \num{6800} USD, plus de 4 fois celui du Cameroun, pourtant son espérance de vie (65 ans) reste très inférieure à celle de la France (83 ans), car la richesse y est très inégalement répartie. Le PIB moyen masque les inégalités : une minorité riche peut coexister avec une majorité pauvre.
\item \textbf{Deux indicateurs complémentaires :} le taux de scolarisation (ou le taux d'alphabétisation des adultes) et le taux d'accès à l'eau potable (ou à l'électricité, ou la mortalité infantile).
\end{enumerate}
\textbf{Points de vigilance :} dans le calcul du rapport, bien vérifier l'ordre de grandeur : $\num{44000}/\num{1600} = 27,5$ et non 2,75. Toujours citer les chiffres du tableau pour étayer la réponse.
\end{corrige}

\begin{corrige}[Exercice 6 — Causes du sous-développement]
\begin{enumerate}
\item \textbf{Classement :}
\begin{itemize}
\item \emph{Causes internes :} (b) la corruption et la mauvaise gouvernance ; (f) les conflits armés ; (e) la faiblesse de l'industrialisation (insuffisance des politiques nationales d'industrie).
\item \emph{Causes externes :} (a) la faible transformation locale des matières premières (imposée par la division internationale du travail héritée de la colonie) ; (c) les frontières héritées de la colonisation ; (d) les termes inégaux de l'échange international.
\end{itemize}
\emph{Remarque :} (a) et (e) sont liées ; on accepte (a) en externe (structure héritée du système colonial d'échange) et (e) en interne (choix de politique économique), à condition de justifier.
\item \textbf{Explication de deux facteurs :}
\begin{itemize}
\item \emph{Mauvaise gouvernance (b) :} la corruption détourne les ressources publiques destinées aux écoles, hôpitaux et routes ; elle décourage les investisseurs et freine la croissance camerounaise.
\item \emph{Échanges inégaux (d) :} le Cameroun vend ses matières premières brutes (cacao, pétrole, bois) à bas prix et achète les produits manufacturés à prix élevé ; la valeur ajoutée s'accumule à l'extérieur, appauvrissant le pays.
\end{itemize}
\item \textbf{Solutions panafricanistes :}
\begin{itemize}
\item Contre (b) : coopération continentale de lutte contre la corruption (mécanismes de l'Union africaine, évaluation entre pairs), transparence dans la gestion des ressources.
\item Contre (d) : intégration économique par la ZLECAf pour créer des chaînes de valeur africaines : transformer le cacao sur le continent et vendre le chocolat sur le marché commun de 1,3 milliard de consommateurs.
\end{itemize}
\end{enumerate}
\textbf{Points de vigilance :} ne pas opposer mécaniquement interne/externe : les deux types de causes s'entremêlent (la colonisation a installé des structures que les élites postcoloniales ont parfois maintenues). Le correcteur attend cette nuance.
\end{corrige}

\begin{corrige}[Exercice 7 — Objectifs du panafricanisme]
\textbf{Texte argumenté de référence (15 à 20 lignes) :}

Le panafricanisme peut être considéré comme une philosophie du développement de l'Afrique parce qu'il ne se contente pas de rêver d'une unité politique : il fait de cette unité la \emph{condition} du développement du continent. D'abord, l'unité politique est pour Kwame Nkrumah la condition de la souveraineté réelle : dans \emph{L'Afrique doit s'unir} (1963), il affirme qu'aucun État africain isolé n'est assez fort pour négocier seul son avenir ; divisés, les pays restent faibles et exploités, unis ils peuvent parler d'égal à égal avec les puissances du monde. L'indépendance politique sans unité économique n'est qu'une « indépendance en trompe-l'œil ». Ensuite, l'intégration économique est au cœur du projet : en unissant leurs marchés, les pays africains peuvent industrialiser le continent, transformer leurs matières premières et échapper aux échanges inégaux qui les appauvrissent. La ZLECAf, entrée en vigueur en 2021, concrétise cet objectif en créant un marché unique de 1,3 milliard de consommateurs. Enfin, le panafricanisme est aussi une philosophie de la dignité culturelle : avec Cheikh Anta Diop, il réhabilite les cultures et les langues africaines, condition d'un développement endogène, pensé par et pour les Africains plutôt qu'imité de l'extérieur. Ainsi, le panafricanisme articule unité politique, intégration économique et fierté culturelle : il propose une voie proprement africaine vers le développement, fondée sur les forces propres du continent et la solidarité entre Africains et diaspora.

\textbf{Points de vigilance :} le texte doit mobiliser au moins deux objectifs (ici trois) et un auteur du programme (ici Nkrumah et Diop) ; chaque objectif doit être relié explicitement à la notion de développement.
\end{corrige}

\begin{corrige}[Exercice 8 — Situation concrète : la transformation du cacao]
\begin{enumerate}
\item \textbf{Le paradoxe « riche mais sous-développé » :} le Cameroun possède une richesse agricole précieuse (le cacao, dont il est l'un des premiers producteurs mondiaux), mais il en exporte la matière brute à bas prix et réimporte le produit fini (le chocolat) à prix élevé. La plus grande partie de la valeur ajoutée (transformation, emballage, commercialisation) est captée à l'étranger : le pays est riche en ressources mais pauvre en revenus.
\item \textbf{Cause révélée :} la \cle{faible transformation locale des matières premières}, héritée de la division coloniale du travail (la colonie fournit les matières brutes, la métropole transforme), entretenue par les \cle{termes inégaux de l'échange} international.
\item \textbf{Mesures :}
\begin{itemize}
\item \emph{À l'échelle nationale :} construire des unités locales de transformation (usines de chocolat, de beurre de cacao), soutenues par la formation de techniciens et des incitations fiscales aux investisseurs dans l'agro-industrie.
\item \emph{À l'échelle continentale (panafricaniste) :} utiliser la ZLECAf pour vendre le chocolat africain sans droits de douane sur le marché continental de 1,3 milliard de consommateurs, et coordonner entre pays producteurs (Cameroun, Côte d'Ivoire, Ghana) une politique commune de prix et de transformation, à l'image d'un « cartel du cacao ».
\end{itemize}
\end{enumerate}
\textbf{Points de vigilance :} la solution continentale doit être explicitement reliée au panafricanisme (unité économique, marché commun, solidarité entre producteurs) et non rester une simple mesure commerciale.
\end{corrige}

\begin{corrige}[Exercice 9 — Dissertation : le panafricanisme est-il une utopie ?]
\textbf{1. Introduction (4 pts) :}

\emph{Accroche :} En 2021, l'entrée en vigueur de la ZLECAf a relancé un vieux rêve africain : celui d'un continent uni. \emph{Présentation :} Né dans la diaspora à la fin du XIX\textsuperscript{e} siècle, le panafricanisme est le mouvement de pensée et d'action qui vise l'unité politique, économique et culturelle de l'Afrique et de ses diasporas. \emph{Définition des termes :} le \cle{panafricanisme} désigne ce projet d'unité continentale au service du développement ; l'\cle{utopie}, au sens courant, désigne un projet séduisant mais irréalisable, un rêve sans prise sur le réel. \emph{Problématique :} le panafricanisme est-il une utopie ? Autrement dit : l'unité africaine relève-t-elle du rêve impossible, ou constitue-t-elle au contraire la voie réaliste du développement du continent ? \emph{Annonce du plan :} nous verrons d'abord que les échecs et les obstacles font paraître le panafricanisme utopique ; nous montrerons ensuite qu'il est une voie réaliste et nécessaire ; nous conclurons qu'il est un projet en construction, condition du développement endogène de l'Afrique.

\textbf{2. Développement (12 pts) :}

\emph{Thèse — le panafricanisme comme utopie (6 pts) :}
\begin{itemize}
\item \textbf{Argument 1 :} les États africains restent divisés par les frontières héritées de la colonisation (Conférence de Berlin, 1885), que l'OUA elle-même a consacrées en 1963 : l'unité politique rêvée par Nkrumah (« États-Unis d'Afrique ») n'a jamais été réalisée.
\item \textbf{Argument 2 :} la faiblesse des échanges intra-africains (environ 15 \% du commerce total, contre 60 \% en Europe) montre que le continent reste économiquement tourné vers l'extérieur et dépendant.
\item \textbf{Argument 3 :} les échecs partiels de l'OUA (incapacité à prévenir les conflits, non-ingérence, rivalités entre chefs d'État) illustrent l'écart entre le rêve et la réalité.
\end{itemize}

\emph{Antithèse — une voie réaliste et nécessaire (6 pts) :}
\begin{itemize}
\item \textbf{Argument 1 :} des réussites concrètes existent : la transformation de l'OUA en Union africaine (2002) a doté le continent d'institutions plus ambitieuses (Conseil de paix et de sécurité, Parlement panafricain).
\item \textbf{Argument 2 :} la ZLECAf (2021), signée par 54 pays, crée un marché unique de 1,3 milliard de consommateurs : c'est la concrétisation économique du projet de Nkrumah.
\item \textbf{Argument 3 :} l'unité est une \emph{nécessité} : aucun pays africain isolé ne peut négocier seul face aux grandes puissances ; les intégrations régionales (CEMAC, CEDEAO) prouvent que la coopération fonctionne.
\end{itemize}

\textbf{3. Conclusion (4 pts) :}

Le panafricanisme n'est donc ni un rêve accompli ni une pure illusion : c'est un \emph{projet en construction}. Ses difficultés sont réelles, mais ses avancées (Union africaine, ZLECAf) montrent qu'il progresse. Surtout, il demeure la condition d'un développement endogène de l'Afrique : seule l'unité permettra au continent de transformer ses richesses et de parler d'égal à égal avec le monde. L'utopie d'hier peut devenir la réalité de demain, à condition que les États et les jeunes générations s'en emparent.

\textbf{Points de vigilance :} l'introduction doit définir les DEUX termes du sujet ; la problématique doit reprendre la question ; la conclusion doit trancher clairement. Éviter le hors-sujet : il ne s'agit pas de décrire le sous-développement, mais de juger le caractère utopique ou réaliste du panafricanisme.
\end{corrige}

\begin{corrige}[Exercice 10 — Commentaire de texte : Kwame Nkrumah]
\textbf{1. Thèse de l'auteur (3 pts) :} L'unité politique et surtout économique de l'Afrique est la condition nécessaire de sa liberté réelle et de son développement : sans elle, les indépendances demeurent fictives et le continent reste exploité.

\textbf{2. Trois arguments (6 pts) :}
\begin{itemize}
\item \emph{Argument de la faiblesse de la division :} « Divisés, nous restons faibles et exploités » — la fragmentation héritée de la colonisation rend chaque État vulnérable face aux puissances extérieures.
\item \emph{Argument des ressources :} « Nous avons toutes les ressources nécessaires : les minerais, les terres fertiles, les hommes » — l'Afrique possède en elle-même les moyens de son développement ; l'unité permettrait de les mobiliser collectivement.
\item \emph{Argument de l'indépendance trompe-l'œil :} « sans l'unité économique du continent, elle demeure une indépendance en trompe-l'œil » — la souveraineté politique sans puissance économique est une souveraineté formelle, vidée de son contenu.
\end{itemize}

\textbf{3. Une philosophie du développement (5 pts) :} Nkrumah ne réclame pas seulement une union diplomatique ou symbolique : il subordonne l'unité politique à l'\emph{unité économique}, seule capable de transformer les ressources du continent en développement réel. Le panafricanisme devient ainsi une doctrine du développement endogène : l'Afrique doit compter sur ses propres forces (ressources, hommes) et sur son marché uni plutôt que sur l'extérieur. L'enjeu n'est pas seulement de gouverner, mais de « négocier son avenir », c'est-à-dire de maîtriser son destin économique.

\textbf{4. Discussion (6 pts) :} L'appel de 1963 reste d'actualité.
\begin{itemize}
\item \emph{Argument 1 :} la ZLECAf (2021) montre que le diagnostic de Nkrumah est enfin pris au sérieux : 54 pays créent un marché unique de 1,3 milliard de consommateurs, concrétisant l'unité économique qu'il appelait de ses vœux. Les échanges intra-africains (environ 15 \%) restent cependant faibles : le combat n'est pas terminé.
\item \emph{Argument 2 :} aucun pays africain, même le Cameroun, ne pèse seul dans les négociations commerciales mondiales ; seule une voix continentale unie peut défendre les intérêts africains (prix des matières premières, transformation locale). L'actualité confirme donc la pertinence du texte.
\end{itemize}

\textbf{Points de vigilance :} distinguer clairement thèse (ce que l'auteur soutient) et arguments (ce qui la justifie) ; dans la discussion, produire des arguments \emph{personnels} appuyés sur des faits actuels, et non répéter le texte.
\end{corrige}

\begin{corrige}[Exercice 11 — Analyse de documents et étude de cas : la ZLECAf]
\textbf{1. Le paradoxe africain (4 pts) :} le document 1 montre que l'Afrique exporte environ 80 \% de matières premières : elle est riche de ressources mais n'en tire pas la valeur ajoutée de la transformation, d'où son sous-développement. De plus, ses échanges internes ne représentent que 15 \% de son commerce total, contre 60 \% pour l'Europe : le continent est mal intégré à son propre marché, chaque pays commerçant davantage avec l'extérieur qu'avec ses voisins. Paradoxe : un continent riche, mais pauvre et divisé.

\textbf{2. Deux objectifs panafricanistes concrétisés par la ZLECAf (4 pts) :}
\begin{itemize}
\item \emph{L'intégration économique :} en supprimant progressivement les droits de douane entre 54 pays, la ZLECAf crée le marché commun rêvé par Nkrumah dès 1963.
\item \emph{L'unité du continent et le développement endogène :} en unissant 1,3 milliard de consommateurs, elle permet à l'Afrique de compter sur ses propres forces (marché intérieur, transformation locale) plutôt que sur l'extérieur.
\end{itemize}

\textbf{3. Étude de cas : la ZLECAf peut-elle transformer l'économie camerounaise ? (8 pts) :}

\emph{Opportunités :}
\begin{itemize}
\item \textbf{Marché élargi :} les entreprises camerounaises accèdent à 1,3 milliard de consommateurs sans droits de douane, au lieu du seul marché national (environ 27 millions d'habitants) : débouchés nouveaux pour le café, le cacao transformé, le bois ouvré.
\item \textbf{Industrialisation locale :} le marché continental rend rentables les usines de transformation (chocolat, meubles, produits agroalimentaires) : le Cameroun peut passer de l'exportation de matières brutes à celle de produits finis, créant emplois et valeur ajoutée — c'est la voie du \cle{développement endogène}.
\end{itemize}

\emph{Difficultés :}
\begin{itemize}
\item \textbf{Concurrence :} les produits industriels de pays plus avancés (Afrique du Sud, Égypte, Nigeria) peuvent inonder le marché camerounais et étouffer les jeunes industries locales.
\item \textbf{Infrastructures et gouvernance :} routes, ports et voies ferrées insuffisants (le corridor Douala--N'Djamena reste coûteux), lourdeurs administratives et corruption renchérissent les échanges et peuvent annuler le bénéfice de la suppression des droits de douane.
\end{itemize}

\textbf{4. Conclusion : deux conditions nécessaires (4 pts) :}
\begin{enumerate}
\item investir dans les infrastructures (routes, ports, énergie) et former une main-d'œuvre qualifiée pour rendre les entreprises camerounaises compétitives ;
\item améliorer la gouvernance (lutte contre la corruption, simplification des procédures douanières) et soutenir l'industrialisation locale, afin que l'\cle{unité} continentale serve réellement le \cle{développement endogène} du pays.
\end{enumerate}

\textbf{Points de vigilance :} mobiliser explicitement les deux notions exigées (\emph{unité} et \emph{développement endogène}) ; équilibrer les deux opportunités et les deux difficultés ; appuyer l'analyse sur les chiffres du document 1 (15 \%, 80 \%, 1,3 milliard).
\end{corrige}

\newpage

\coursec{Fiche bilan}

\begin{popfigure}
\begin{tikzpicture}[mindmap, grow cyclic, every node/.style={concept, font=\small}, concept color=popblue, text=white, level 1 concept/.append style={level distance=3.6cm, sibling angle=60, font=\scriptsize}, level 2 concept/.append style={level distance=2.2cm, font=\tiny}]
\node[concept, concept color=popdark, font=\small\bfseries] {Panafricanisme et développement de l'Afrique}
  child[concept color=popblue] { node {Développement}
    child { node {Croissance + structures + progrès humain} }
    child { node {Indicateurs : PIB, IDH, alphabétisation} }
  }
  child[concept color=poporange] { node {Sous-développement}
    child { node {Pauvreté, dépendance, inégalités} }
    child { node {Causes : colonialisme, matières premières, gouvernance} }
  }
  child[concept color=popgreen] { node {Historique}
    child { node {1900 : Congrès de Londres (Sylvester-Williams)} }
    child { node {Du Bois, Garvey, Nkrumah, Senghor} }
    child { node {1963 : création de l'OUA ; 2002 : UA} }
  }
  child[concept color=poppurple] { node {Objectifs}
    child { node {Unité politique et économique} }
    child { node {Indépendance, dignité, développement endogène} }
  }
  child[concept color=poppink] { node {Difficultés}
    child { node {Nationalismes, frontières héritées} }
    child { node {Dépendance économique, ingérences} }
  }
  child[concept color=popgold] { node {Débat}
    child { node {Utopie ou voie d'avenir ?} }
  };
\end{tikzpicture}
\legende{Carte mentale de la leçon : le panafricanisme comme philosophie du développement de l'Afrique (définitions, historique, objectifs, difficultés, débat).}
\end{popfigure}

\begin{bilan}
\textbf{Définitions clés à connaître par c\oe ur}
\begin{itemize}
\item \cle{Développement} : processus global de transformation (économique, sociale, culturelle, politique) qui combine la croissance quantitative et le changement qualitatif des structures, en vue de l'amélioration durable des conditions de vie et de l'épanouissement de l'homme (F. Perroux : « combinaison des changements mentaux et sociaux »).
\item \cle{Sous-développement} : état de retard économique et social caractérisé par la pauvreté, la dépendance extérieure, l'insuffisance industrielle et la faiblesse des indicateurs humains ; il n'est pas une fatalité mais le produit de causes historiques et structurelles.
\item \cle{Panafricanisme} : mouvement de pensée et d'action politique visant l'unité, la solidarité et la libération des peuples africains et de la diaspora, conçu comme philosophie du développement de l'Afrique.
\item \cle{Développement endogène} : développement pensé à partir des ressources, des valeurs et des besoins propres de l'Afrique, et non par simple imitation des modèles extérieurs.
\end{itemize}

\textbf{Repères historiques à retenir}
\begin{center}
\begin{tabularx}{\linewidth}{|l|X|}
\hline
\rowcolor{popblueL} \textbf{Date} & \textbf{Événement} \\
\hline
1900 & Premier congrès panafricain de Londres (H. Sylvester-Williams). \\
\hline
1919--1945 & Congrès panafricains animés par W. E. B. Du Bois ; mouvement Garvey ; négritude (Senghor, Césaire). \\
\hline
1945 & Congrès de Manchester : radicalisation, exigence d'indépendances. \\
\hline
1957 & Indépendance du Ghana ; K. Nkrumah, figure de l'unité africaine. \\
\hline
1963 & Création de l'OUA à Addis-Abeba. \\
\hline
2002 & Transformation de l'OUA en Union africaine (UA). \\
\hline
\end{tabularx}
\end{center}

\textbf{Objectifs et difficultés du panafricanisme}
\begin{itemize}
\item \textbf{Objectifs} : unité politique et économique du continent ; indépendance réelle ; restauration de la dignité africaine ; développement solidaire et endogène.
\item \textbf{Difficultés} : nationalismes étroits et frontières héritées de la colonisation ; dépendance économique et financière ; ingérences extérieures ; conflits internes ; faiblesse des institutions panafricaines.
\end{itemize}

\textbf{Méthode express (dissertation)}
\begin{enumerate}
\item Définir les deux concepts du sujet (développement, panafricanisme) dans l'introduction.
\item Problématiser : le panafricanisme est-il une utopie ou une voie réaliste de développement ?
\item Plan dialectique : thèse (nécessité de l'unité), antithèse (obstacles et échecs), synthèse (unité rénovée, intégration régionale, développement endogène).
\item Illustrer avec des exemples camerounais et africains (CEMAC, UA, ZLECAf).
\end{enumerate}
\end{bilan}

\begin{aretenir}[Checklist avant le Bac]
\begin{itemize}
\item[$\square$] Je sais définir le développement en distinguant croissance économique et progrès humain.
\item[$\square$] Je connais au moins deux indicateurs du développement (PIB, IDH).
\item[$\square$] Je sais définir le sous-développement et citer ses manifestations (pauvreté, dépendance, faible industrialisation).
\item[$\square$] Je distingue les causes externes (colonialisme, échange inégal) et internes (gouvernance) du sous-développement.
\item[$\square$] Je connais les grandes étapes du panafricanisme (Londres 1900, Manchester 1945, OUA 1963, UA 2002).
\item[$\square$] Je sais citer au moins trois figures du panafricanisme (Du Bois, Nkrumah, Senghor).
\item[$\square$] Je sais énoncer les objectifs du panafricanisme comme philosophie du développement.
\item[$\square$] Je connais les principales difficultés et limites du panafricanisme.
\item[$\square$] Je sais construire un plan dialectique sur le thème « panafricanisme : utopie ou voie d'avenir ? ».
\item[$\square$] Je sais illustrer mon devoir avec des exemples africains et camerounais actuels (UA, CEMAC, ZLECAf).
\end{itemize}
\end{aretenir}

\findecours

\end{document}
